\documentclass[12pt]{book}
% Idioma y codificación
\usepackage[spanish, es-tabla]{babel}       %es-tabla para que se titule "Tabla"
\usepackage[utf8]{inputenc}

% Márgenes
\usepackage[a4paper,top=3cm,bottom=2.5cm,left=3cm,right=3cm]{geometry}

% Comentarios de bloque
\usepackage{verbatim}

% Paquetes de links
\usepackage[hidelinks]{hyperref}    % Permite enlaces
\usepackage{url}                    % redirecciona a la web

% Más opciones para enumeraciones
\usepackage{enumitem}

% Personalizar la portada
\usepackage{titling}


% Paquetes de tablas
\usepackage{multirow}


%------------------------------------------------------------------------

%Paquetes de figuras
\usepackage{caption}
\usepackage{subcaption} % Figuras al lado de otras
\usepackage{float}      % Poner figuras en el sitio indicado H.


% Paquetes de imágenes
\usepackage{graphicx}       % Paquete para añadir imágenes
\usepackage{transparent}    % Para manejar la opacidad de las figuras

% Paquete para usar colores
\usepackage[dvipsnames]{xcolor}
\usepackage{pagecolor}      % Para cambiar el color de la página

% Habilita tamaños de fuente mayores
\usepackage{fix-cm}


%------------------------------------------------------------------------

% Paquetes de matemáticas
\usepackage{mathtools, amsfonts, amssymb, mathrsfs}
\usepackage[makeroom]{cancel}     % Simplificar tachando
\usepackage{polynom}    % Divisiones y Ruffini
\usepackage{units} % Para poner fracciones diagonales con \nicefrac

\usepackage{pgfplots}   %Representar funciones
\pgfplotsset{compat=1.18}  % Versión 1.18

\usepackage{tikz-cd}    % Para usar diagramas de composiciones
\usetikzlibrary{calc}   % Para usar cálculo de coordenadas en tikz

%Definición de teoremas, etc.
\usepackage{amsthm}
%\swapnumbers   % Intercambia la posición del texto y de la numeración

\theoremstyle{plain}

\makeatletter
\@ifclassloaded{article}{
  \newtheorem{teo}{Teorema}[section]
}{
  \newtheorem{teo}{Teorema}[chapter]  % Se resetea en cada chapter
}
\makeatother

\newtheorem{coro}{Corolario}[teo]           % Se resetea en cada teorema
\newtheorem{prop}[teo]{Proposición}         % Usa el mismo contador que teorema
\newtheorem{lema}[teo]{Lema}                % Usa el mismo contador que teorema

\theoremstyle{remark}
\newtheorem*{observacion}{Observación}

\theoremstyle{definition}

\makeatletter
\@ifclassloaded{article}{
  \newtheorem{definicion}{Definición} [section]     % Se resetea en cada chapter
}{
  \newtheorem{definicion}{Definición} [chapter]     % Se resetea en cada chapter
}
\makeatother

\newtheorem*{notacion}{Notación}
\newtheorem*{ejemplo}{Ejemplo}
\newtheorem*{ejercicio*}{Ejercicio}             % No numerado
\newtheorem{ejercicio}{Ejercicio} [section]     % Se resetea en cada section


% Modificar el formato de la numeración del teorema "ejercicio"
\renewcommand{\theejercicio}{%
  \ifnum\value{section}=0 % Si no se ha iniciado ninguna sección
    \arabic{ejercicio}% Solo mostrar el número de ejercicio
  \else
    \thesection.\arabic{ejercicio}% Mostrar número de sección y número de ejercicio
  \fi
}


% \renewcommand\qedsymbol{$\blacksquare$}         % Cambiar símbolo QED
%------------------------------------------------------------------------

% Paquetes para encabezados
\usepackage{fancyhdr}
\pagestyle{fancy}
\fancyhf{}

\newcommand{\helv}{ % Modificación tamaño de letra
\fontfamily{}\fontsize{12}{12}\selectfont}
\setlength{\headheight}{15pt} % Amplía el tamaño del índice


%\usepackage{lastpage}   % Referenciar última pag   \pageref{LastPage}
\fancyfoot[C]{\thepage}

%------------------------------------------------------------------------

% Conseguir que no ponga "Capítulo 1". Sino solo "1."
\makeatletter
\@ifclassloaded{book}{
  \renewcommand{\chaptermark}[1]{\markboth{\thechapter.\ #1}{}} % En el encabezado
    
  \renewcommand{\@makechapterhead}[1]{%
  \vspace*{50\p@}%
  {\parindent \z@ \raggedright \normalfont
    \ifnum \c@secnumdepth >\m@ne
      \huge\bfseries \thechapter.\hspace{1em}\ignorespaces
    \fi
    \interlinepenalty\@M
    \Huge \bfseries #1\par\nobreak
    \vskip 40\p@
  }}
}
\makeatother

%------------------------------------------------------------------------
% Paquetes de cógido
\usepackage{minted}
\renewcommand\listingscaption{Código fuente}

\usepackage{fancyvrb}
% Personaliza el tamaño de los números de línea
\renewcommand{\theFancyVerbLine}{\small\arabic{FancyVerbLine}}

% Estilo para C++
\newminted{cpp}{
    frame=lines,
    framesep=2mm,
    baselinestretch=1.2,
    linenos,
    escapeinside=||
}



\usepackage{listings} % Para incluir código desde un archivo

\renewcommand\lstlistingname{Código Fuente}
\renewcommand\lstlistlistingname{Índice de Códigos Fuente}

% Definir colores
\definecolor{vscodepurple}{rgb}{0.5,0,0.5}
\definecolor{vscodeblue}{rgb}{0,0,0.8}
\definecolor{vscodegreen}{rgb}{0,0.5,0}
\definecolor{vscodegray}{rgb}{0.5,0.5,0.5}
\definecolor{vscodebackground}{rgb}{0.97,0.97,0.97}
\definecolor{vscodelightgray}{rgb}{0.9,0.9,0.9}

% Configuración para el estilo de C similar a VSCode
\lstdefinestyle{vscode_C}{
  backgroundcolor=\color{vscodebackground},
  commentstyle=\color{vscodegreen},
  keywordstyle=\color{vscodeblue},
  numberstyle=\tiny\color{vscodegray},
  stringstyle=\color{vscodepurple},
  basicstyle=\scriptsize\ttfamily,
  breakatwhitespace=false,
  breaklines=true,
  captionpos=b,
  keepspaces=true,
  numbers=left,
  numbersep=5pt,
  showspaces=false,
  showstringspaces=false,
  showtabs=false,
  tabsize=2,
  frame=tb,
  framerule=0pt,
  aboveskip=10pt,
  belowskip=10pt,
  xleftmargin=10pt,
  xrightmargin=10pt,
  framexleftmargin=10pt,
  framexrightmargin=10pt,
  framesep=0pt,
  rulecolor=\color{vscodelightgray},
  backgroundcolor=\color{vscodebackground},
}

%------------------------------------------------------------------------

% Comandos definidos
\newcommand{\bb}[1]{\mathbb{#1}}
\newcommand{\cc}[1]{\mathcal{#1}}

% I prefer the slanted \leq
\let\oldleq\leq % save them in case they're every wanted
\let\oldgeq\geq
\renewcommand{\leq}{\leqslant}
\renewcommand{\geq}{\geqslant}

% Si y solo si
\newcommand{\sii}{\iff}

% Letras griegas
\newcommand{\eps}{\epsilon}
\newcommand{\veps}{\varepsilon}
\newcommand{\lm}{\lambda}

\newcommand{\ol}{\overline}
\newcommand{\ul}{\underline}
\newcommand{\wt}{\widetilde}
\newcommand{\wh}{\widehat}

\let\oldvec\vec
\renewcommand{\vec}{\overrightarrow}

% Derivadas parciales
\newcommand{\del}[2]{\frac{\partial #1}{\partial #2}}
\newcommand{\Del}[3]{\frac{\partial^{#1} #2}{\partial^{#1} #3}}
\newcommand{\deld}[2]{\dfrac{\partial #1}{\partial #2}}
\newcommand{\Deld}[3]{\dfrac{\partial^{#1} #2}{\partial^{#1} #3}}


\newcommand{\AstIg}{\stackrel{(\ast)}{=}}
\newcommand{\Hop}{\stackrel{L'H\hat{o}pital}{=}}

\newcommand{\red}[1]{{\color{red}#1}} % Para integrales, destacar los cambios.

% Método de integración
\newcommand{\MetInt}[2]{
    \left[\begin{array}{c}
        #1 \\ #2
    \end{array}\right]
}

% Declarar aplicaciones
% 1. Nombre aplicación
% 2. Dominio
% 3. Codominio
% 4. Variable
% 5. Imagen de la variable
\newcommand{\Func}[5]{
    \begin{equation*}
        \begin{array}{rrll}
            #1:& #2 & \longrightarrow & #3\\
               & #4 & \longmapsto & #5
        \end{array}
    \end{equation*}
}

%------------------------------------------------------------------------


% Para poder incluir árboles
\usepackage{forest}
\usepackage{booktabs}

% Para poder añadir autómatas
% https://www3.nd.edu/~kogge/courses/cse30151-fa17/Public/other/tikz_tutorial.pdf
\usetikzlibrary{automata} %, positioning, arrows}
\tikzset{
    -Stealth,
    node distance=3cm, % specifies the minimum distance between two nodes. Change if necessary.
    every state/.style={thick, fill=gray!10}, % sets the properties for each ’state’ node
    initial text=$ $, % sets the text that appears on the start arrow
    % Un tipo de nodo, que es error, que lo pone rojo
    error/.style={thick, fill=red!20},
}

\begin{document}

    % 1. Foto de fondo
    % 2. Título
    % 3. Encabezado Izquierdo
    % 4. Color de fondo
    % 5. Coord x del titulo
    % 6. Coord y del titulo
    % 7. Fecha
    % 8. Autor

    
    % 1. Foto de fondo
% 2. Título
% 3. Encabezado Izquierdo
% 4. Color de fondo
% 5. Coord x del titulo
% 6. Coord y del titulo
% 7. Fecha

\newcommand{\portada}[7]{

    \portadaBase{#1}{#2}{#3}{#4}{#5}{#6}{#7}
    \portadaBook{#1}{#2}{#3}{#4}{#5}{#6}{#7}
}

\newcommand{\portadaExamen}[7]{

    \portadaBase{#1}{#2}{#3}{#4}{#5}{#6}{#7}
    \portadaArticle{#1}{#2}{#3}{#4}{#5}{#6}{#7}
}




\newcommand{\portadaBase}[7]{

    % Tiene la portada principal y la licencia Creative Commons
    
    % 1. Foto de fondo
    % 2. Título
    % 3. Encabezado Izquierdo
    % 4. Color de fondo
    % 5. Coord x del titulo
    % 6. Coord y del titulo
    % 7. Fecha
    
    
    \thispagestyle{empty}               % Sin encabezado ni pie de página
    \newgeometry{margin=0cm}        % Márgenes nulos para la primera página
    
    
    % Encabezado
    \fancyhead[L]{\helv #3}
    \fancyhead[R]{\helv \nouppercase{\leftmark}}
    
    
    \pagecolor{#4}        % Color de fondo para la portada
    
    \begin{figure}[p]
        \centering
        \transparent{0.3}           % Opacidad del 30% para la imagen
        
        \includegraphics[width=\paperwidth, keepaspectratio]{assets/#1}
    
        \begin{tikzpicture}[remember picture, overlay]
            \node[anchor=north west, text=white, opacity=1, font=\fontsize{60}{90}\selectfont\bfseries\sffamily, align=left] at (#5, #6) {#2};
            
            \node[anchor=south east, text=white, opacity=1, font=\fontsize{12}{18}\selectfont\sffamily, align=right] at (9.7, 3) {\textbf{\href{https://losdeldgiim.github.io/}{Los Del DGIIM}}};
            
            \node[anchor=south east, text=white, opacity=1, font=\fontsize{12}{15}\selectfont\sffamily, align=right] at (9.7, 1.8) {Doble Grado en Ingeniería Informática y Matemáticas\\Universidad de Granada};
        \end{tikzpicture}
    \end{figure}
    
    
    \restoregeometry        % Restaurar márgenes normales para las páginas subsiguientes
    \pagecolor{white}       % Restaurar el color de página
    
    
    \newpage
    \thispagestyle{empty}               % Sin encabezado ni pie de página
    \begin{tikzpicture}[remember picture, overlay]
        \node[anchor=south west, inner sep=3cm] at (current page.south west) {
            \begin{minipage}{0.5\paperwidth}
                \href{https://creativecommons.org/licenses/by-nc-nd/4.0/}{
                    \includegraphics[height=2cm]{assets/Licencia.png}
                }\vspace{1cm}\\
                Esta obra está bajo una
                \href{https://creativecommons.org/licenses/by-nc-nd/4.0/}{
                    Licencia Creative Commons Atribución-NoComercial-SinDerivadas 4.0 Internacional (CC BY-NC-ND 4.0).
                }\\
    
                Eres libre de compartir y redistribuir el contenido de esta obra en cualquier medio o formato, siempre y cuando des el crédito adecuado a los autores originales y no persigas fines comerciales. 
            \end{minipage}
        };
    \end{tikzpicture}
    
    
    
    % 1. Foto de fondo
    % 2. Título
    % 3. Encabezado Izquierdo
    % 4. Color de fondo
    % 5. Coord x del titulo
    % 6. Coord y del titulo
    % 7. Fecha


}


\newcommand{\portadaBook}[7]{

    % 1. Foto de fondo
    % 2. Título
    % 3. Encabezado Izquierdo
    % 4. Color de fondo
    % 5. Coord x del titulo
    % 6. Coord y del titulo
    % 7. Fecha

    % Personaliza el formato del título
    \pretitle{\begin{center}\bfseries\fontsize{42}{56}\selectfont}
    \posttitle{\par\end{center}\vspace{2em}}
    
    % Personaliza el formato del autor
    \preauthor{\begin{center}\Large}
    \postauthor{\par\end{center}\vfill}
    
    % Personaliza el formato de la fecha
    \predate{\begin{center}\huge}
    \postdate{\par\end{center}\vspace{2em}}
    
    \title{#2}
    \author{\href{https://losdeldgiim.github.io/}{Los Del DGIIM}}
    \date{Granada, #7}
    \maketitle
    
    \tableofcontents
}




\newcommand{\portadaArticle}[7]{

    % 1. Foto de fondo
    % 2. Título
    % 3. Encabezado Izquierdo
    % 4. Color de fondo
    % 5. Coord x del titulo
    % 6. Coord y del titulo
    % 7. Fecha

    % Personaliza el formato del título
    \pretitle{\begin{center}\bfseries\fontsize{42}{56}\selectfont}
    \posttitle{\par\end{center}\vspace{2em}}
    
    % Personaliza el formato del autor
    \preauthor{\begin{center}\Large}
    \postauthor{\par\end{center}\vspace{3em}}
    
    % Personaliza el formato de la fecha
    \predate{\begin{center}\huge}
    \postdate{\par\end{center}\vspace{5em}}
    
    \title{#2}
    \author{\href{https://losdeldgiim.github.io/}{Los Del DGIIM}}
    \date{Granada, #7}
    \thispagestyle{empty}               % Sin encabezado ni pie de página
    \maketitle
    \vfill
}
    \portada{etsiitA4.jpg}{Modelos de\\Computación}{Modelos de Computación}{MidnightBlue}{-8}{27}{2024-2025}{Arturo Olivares Martos}

    %\chapter{Ensamblador}
\section{Historia}
La linea de procesadores Intel, comúnmente conocida como \textit{x86}
ha seguido un largo proceso de evolución, manteniendo compatibilidad
hasta 8086, introducido en 1978. La familia de procesadores x86 ha seguido una evolución hasta
el día de hoy:
\begin{itemize}
	\item 8086 (1978) $\rightarrow$ 29K transistores, 16 bits, 5 MHz. \quad [Intel 16-bit, IBM PC]
	\item 386 (1985) $\rightarrow$ 275K transistores, 32 bits, 16 MHz. \quad [Intel 32-bit]
	\item Pentium 4E (2004) $\rightarrow$ 125M transistores, 64 bits, 3.4 GHz. \quad [Intel 64-bit, Hyperthreading, Implementación de x86-64 de AMD]
	\item Core i7 Haswell (2013) $\rightarrow$ 1.4G transistores, 64 bits, 3.9 Ghz \quad[Multicore, AVX2]
	\item Core i9 Raptor Lake (2022) $ \rightarrow$ +30G transistores, 64 bits, 5.2 GHz \quad[DDR5, PCIe 5.0, AVX-512]
\end{itemize}
A lo largo de los años, muchas compañias han creado procesadores
que se puedan comparar con los de Intel. Entre ellas, AMD ha sido
la más destacada, cayendo ligeramente detrás de Intel en cuanto a rendimiento
y ofreciendo un precio más competitivo. Sin embargo, en los últimos años, AMD ha
sido capaz de ofrecer procesadores con un rendimiento similar o superior a los de Intel,
concretamente, en 2024, Intel ha sufrido una gran decadencia por problemas de estabilidad
en los ultimos lanzamientos.
Nosotros nos centraremos en la arquitectura x86-64, que es la que se utiliza en el libro de
clase, sin embargo, haremos pequeñas menciones a IA32, cuando sea conveniente.
\section{Lenguaje C, ensamblador y máquina}
Las computadoras modernas utilizan distintas capas de abstracción, ocultando la complejidad
de la implementación. Dos de estas capas son realmente importantes para nosotros,
el formato y comportamiento de las instrucciones de máquina, que viene determinado
por la ISA (Instruction Set Architecture), esta define el formato de las instrucciones,
y el resultado que cada una de estas tendrá sobre el estado de la máquina.
La mayoria de ISAs, incluyendo x86-64, describen el comportamiento de un programa de
forma secuencial, como si se ejecutase una instrucción tras otra, sin embargo, en realidad,
el hardware de la CPU es mucho mas complejo, ejecutando instrucciones de forma paralelizada
o en otro orden, sin embargo, esto se verá con mas detenimiento en la asignatura \textit{Arquitectura de Computadores}.
La otra capa de abstracción ocurre sobre las direcciones de memoria, se nos presentan
como \textit{virtual addresses}, proporcionando un modelo de memoria que aparenta ser
un gran array de bytes, sin embargo, en realidad, la memoria se maneja de otra forma,
como se ve en la asignatura \textit{Sistemas Operativos}.
Recalcamos aunque no vayan a ser estudiadas en profundidad, que existen otras
ISA, como ARM (utilizada en dispositivos móviles), RISC-V (Nueva ISA open-source),
o la misma IA32. Aunque no se mencione más, se recomienda al lector la profundización
sobre ARM, que actualmente está ampliando su uso en el mercado de portatiles. (Intel Ultra, M1, etc.)
Recordamos que podemos clasificar en general las distintas ISAs en dos tipos:
\begin{itemize}
	\item \textbf{CISC} (Complex Instruction Set Computer): Conjunto de instrucciones complejas.
	\item \textbf{RISC} (Reduced Instruction Set Computer): Conjunto de instrucciones reducidas.
\end{itemize}
Realmente la mayoría de ISAs actuales son híbridas, convirtiendo
las instrucciones de CISC en instrucciones de RISC, para mejorar el rendimiento.
\newline
\newline
El codigo x86-64 difiere en gran medida del código C, destacamos las siguientes
partes del estado del procesador que normalmente no se ven en C:
\begin{itemize}
	\item \textbf{PC (Program Counter)}: Registra la dirección de la siguiente instrucción a ejecutar.
	      Se denomina \textit{\%rip} en x86-64.
	\item \textbf{Archivo de registros}: Contiene 16 registros etiquetados con valores de 64 bits.
	      Pueden ser utilizados para almacenar valores temporales, punteros, etc.
	\item \textbf{Codigos de condición}: Contienen información sobre la instrucción mas reciente.
	      Se utilizan para realizar saltos condicionales.
	\item \textbf{Pila}: Se utiliza para almacenar enteros o flotantes. Es un array direccionable por bytes.
\end{itemize}
Ahora, vamos a ver cómo se relaciona el código C con el código ensamblador. Para ello,
habrá que tener en cuenta que el código generado por el compilador puede variar
según el compilador, la versión, las optimizaciones, etc.
\subsection{Proceso de compilación}
Partimos de un número arbitrario de ficheros \textit{.c}, que contienen código C.
Para compilarlos, se pueden utilizar distintos compiladores, nosotros usaremos
\textit{gcc}:
% codigo en bash:
\begin{minted}{bash}
    gcc -Og programa.c -o programa
\end{minted}
Cuando lo ejecutemos:
\begin{center}
	\begin{tikzpicture}[node distance=1.6cm]

		% Styles within tikzpicture
		\tikzstyle{process} = [rectangle, rounded corners, minimum width=3cm, minimum height=1cm, text centered, draw=black, fill=green!30]
		\tikzstyle{arrow} = [thick,->,>=stealth]

		% Nodes
		\node (source) [process, fill=green!60] {Source File};
		\node (preprocessor) [process, below of=source] {Preprocessor};
		\node (compiler) [process, below of=preprocessor] {Compiler};
		\node (assembler) [process, below of=compiler] {Assembler};
		\node (linker) [process, below of=assembler] {Linker};
		\node (executable) [process, fill=green!60, below of=linker] {Executable};

		% Arrows
        \draw [-Stealth] (source) -- (preprocessor);
        \draw [-Stealth](preprocessor) -- node[right]{Preprocessed File (.i file)} (compiler);
        \draw [-Stealth](compiler) -- node[right]{Assembly code (.s file)} (assembler);
        \draw [-Stealth](assembler) -- node[right]{Object code (.o file)} (linker);
        \draw [-Stealth](linker) -- (executable);
	\end{tikzpicture}
\end{center}
En este proceso, gcc primero preprocesa el código, generando un archivo \textit{.i}, que
posteriormente compila, generando un archivo \textit{.s}, que se ensambla para generar
\textit{.o}, finalmente, se enlazan los distintos archivos \textit{.o} para generar el ejecutable.
Vamos a explicar un poco más detenidamente cada uno de los pasos:
\begin{itemize}
	\item \textbf{Preprocesador}: C proporciona ciertas herramientas del lenguaje mediante un preprocesador,
	      que es el primer paso en la compilación. Los más utilizados son \textit{\#include}, para incluir el contenido de un archivo
	      durante la compilación, y \textit{\#define}, para reemplazar un token por una secuencia arbitraria de texto
	      o valores. Además, elimina los comentarios de un archivo fuente. En gcc, podemos generar el código preprocesado con el flag \textit{-E}.
	\item \textbf{Compilación}: El compilador traduce el codigo preprocesado y genera código ensamblador, con gcc, podemos
	      generar el código ensamblador con el flag \textit{-S}.
	\item \textbf{Ensamblado}: El ensamblador traduce el código ensamblador a código de máquina, generando un archivo \textit{.o}.
	      En gcc, podemos generar el código de máquina con el flag \textit{-c}.
	\item \textbf{Enlazado}: El enlazador combina los distintos archivos \textit{.o} generados en un solo ejecutable. También se encarga de combinar
	      las librerias que se hayan utilizado en el código. En gcc, podemos generar el ejecutable con el flag \textit{-o}.
\end{itemize}


\begin{ejemplo}
	Veamos un ejemplo de cómo se relaciona el código C con el código ensamblador (solo se muestra la parte interesante):
	\begin{center} % Centramos el contenido
		\begin{minipage}{0.45\textwidth}
			\begin{minted}[fontsize=\Large]{c}
long plus(long x, long y);

void sumStore(long x, long y,
        long *dest) {
    long t = plus(x, y);
    *dest = t;
}
                \end{minted}
		\end{minipage}
		\begin{minipage}{0.45\textwidth}
			\begin{minted}[fontsize=\Large]{asm}
sumstore:
    pushq %rbx
    movq  %rdx, %rbx
    call  plus
    movq  %raw, (%rbx)
    popq  %rbx
    ret
                \end{minted}
		\end{minipage}
	\end{center}
\end{ejemplo}
Si lo compilamos en nuestra propia máquina, encontraremos muchas directivas que comienzan por . (ignorar)
\subsection{Representación Datos C, IA32, x86-64}
A continuación, mostramos los distintos tamaños en Bytes para los tipos de dato mas comúnes en C. Así como
su representación en Assembly. Debido a sus origenes como una arquitectura de 16 bits que expandió a 32 bits,
Intel usa el termino \textit{word} para referirse a un dato de 16 bits, y \textit{doubleword} para referirse a
una palabra de 32 bits. En x86-64, se ha añadido el termino \textit{quadword} para referirse a una palabra de 64 bits.
\begin{center}
	\begin{tabular}{|c|c|c|c|c|}
		\hline
		\textbf{Tipo de dato en C} & \textbf{Assembly sufijo} & \textbf{Normal 32 bits} & \textbf{IA32} & \textbf{x86-64} \\
		\hline
		\text{unsigned}            &                          & 4                       & 4             & 4               \\
		\text{int}                 & l                        & 4                       & 4             & 4               \\
		\text{long}                & q                        & 4                       & 4             & 8               \\
		\text{char}                & b                        & 1                       & 1             & 1               \\
		\text{short}               & w                        & 2                       & 2             & 2               \\
		\text{float}               & s                        & 4                       & 4             & 4               \\
		\text{double}              & l                        & 8                       & 8             & 8               \\
		\text{long double}         &                          & 8                       & 10/12         & 16              \\
		\text{char *}              &                          & 4                       & 4             & 8               \\
		\hline
	\end{tabular}
\end{center}
Los números de punto flotante se representan en el formato IEEE 754, ya visto en TOC. Se incluye
un diagrama de la representación de un número de punto flotante en simple y doble precision:
\begin{center}
	\includesvg[width=0.8\textwidth]{Imagenes/IEEE}
\end{center}
También cabe mencionar que existe un tipo de dato cuyo tamaño es 80 bits (long double) sin embargo su uso está
completamente desaconsejado debido a su portabilidad entre sistemas.
\section{Registros, operandos y operaciones}
\subsection{Registros}
Un procesador x86-64 tiene 16 registros de propósito general, que almacenan valores de 64 bits.
Estos registros se usan para almacenar valores temporales, punteros, etc. Sus nombres comienzan
con $\%r$, sin embargo, siguen convenciones de nombres debido a su legado. El primer 8086 tenía
8 registros de 16 bits, cada uno tenía un uso específico, de ahí su nombre, con la extensión a
IA32, los registros se extendieron a 32 bits, y con x86-64, a 64 bits, incluyendo además 8
nuevos registros, cuyos nombres van desde $\%r8$ hasta $\%r15$.

\begin{center}
	\includesvg[width=0.8\textwidth]{Imagenes/registers}
\end{center}
\subsection{Operandos}
La mayoría de instrucciones tienen uno o más operandos que especifican el \textit{source} y el destino de la operación.
En x86-64, los operandos pueden ser de tres tipos, (1) Inmediatos, (2) Registros, (3) Memoria. Veamos la sintaxis
para cada uno de ellos junto con la operación que se realiza:
\begin{center}
	\begin{tabular}{|l|l|l|l|}
		\hline
		\textbf{Type} & \textbf{Form}               & \textbf{Operand value}                      & \textbf{Name}       \\
		\hline
		Immediate     & \$\textit{Imm}              & \textit{Imm}                                & Immediate           \\
		Register      & $r_a$                       & $R[r_a]$                                    & Register            \\
		Memory        & \textit{Imm}                & $M[\textit{Imm}]$                           & Absolute            \\
		Memory        & $(r_a)$                     & $M[R[r_a]]$                                 & Indirect            \\
		Memory        & $\textit{Imm}(r_b)$         & $M[\textit{Imm} + R[r_b]]$                  & Base + displacement \\
		Memory        & $(r_b, r_i)$                & $M[R[r_b] + R[r_i]]$                        & Indexed             \\
		Memory        & $\textit{Imm}(r_b, r_i)$    & $M[\textit{Imm} + R[r_b] + R[r_i]]$         & Indexed             \\
		Memory        & $(r_i, s)$                  & $M[R[r_i] \cdot s]$                         & Scaled indexed      \\
		Memory        & $\textit{Imm}(, r_i, s)$    & $M[\textit{Imm} + R[r_i] \cdot s]$          & Scaled indexed      \\
		Memory        & $(r_b, r_i, s)$             & $M[R[r_b] + R[r_i] \cdot s]$                & Scaled indexed      \\
		Memory        & $\textit{Imm}(r_b, r_i, s)$ & $M[\textit{Imm} + R[r_b] + R[r_i] \cdot s]$ & Scaled indexed      \\
		\hline
	\end{tabular}
\end{center}
Los tipos de datos inmediatos se usan para valores constantes. En AT\&T se representan con un signo de dolar, siguiendo la notación
usual de C, por ejemplo \textit{\$0xFF} o \textit{\$323}. Los registros se representan con un signo de porcentaje, por ejemplo, $\%rax$.
Como ya hemos visto, los registros son de 64 bits, sin embargo, se pueden realizar operaciones de 32 bits, 16 bits o 8 bits, más adelante veremos cómo.
Los operandos de memoria se representan con paréntesis, por ejemplo, \textit{(\%rax)}, \textit{0x8(\%rax)}, etc. Podemos interpretarlo como una
desreferenciación de un puntero, es decir, acceder al valor almacenado en la dirección de memoria que se indica.
\begin{ejemplo}
	Veamos un ejemplo de cómo se realiza un swap en x86-64:
	\begin{center} % Centramos el contenido
		\begin{minipage}{0.45\textwidth}
			\begin{minted}[fontsize=\Large]{c}
void swap(long *xp, long *yp) 
{
    long x = *xp;
    *xp = *yp;
    *yp = x;
}
                \end{minted}
		\end{minipage}
		\begin{minipage}{0.45\textwidth}
			\begin{minted}[fontsize=\Large]{asm}
swap:
    movq (%rdi), %rax
    movq (%rsi), %rdx
    movq %rax, (%rsi)
    movq %rdx, (%rdi)
    ret
                \end{minted}
		\end{minipage}
	\end{center}
	Podemos ver que inicialmente, la disposición era $\% rdi \rightarrow xp$, $\% rsi \rightarrow yp$.
	Posteriormente, se almacena el \textbf{contenido} de $\%rdi$ en el registro $\%rax$, y el contenido de $\%rsi$ en $\%rdx$.
	Finalmente, se almacena el contenido de $\%rax$ en la dirección de memoria que apunta $\%rsi$, y el contenido de $\%rdx$ en la dirección de memoria que apunta $\%rdi$.
	Veamos un ejemplo para entender mejor el funcionamiento de los registros, supongamos que tenemos los siguientes valores:
	\begin{itemize}
		\item \textbf{Registro }$\%$rdi: 0x120
		\item \textbf{Registro }$\%$rsi: 0x130
		\item \textbf{Memoria en 0x120: } 125
		\item \textbf{Memoria en 0x130: } 450
	\end{itemize}
	Veamos ahora paso por paso el contenido de cada registro( $\rightarrow$ representa el contenido de la dirección de memoria):
	\[
		\begin{array}{|c|c|c|c|c|c|}
			\hline
			\textbf{Registro} & \textbf{Valor} & \textbf{movq}           & \textbf{movq}           & \textbf{movq}                                 & \textbf{movq}                         \\
			                  &                & \textbf{(\%rdi), \%rax} & \textbf{(\%rsi), \%rdx} & \textbf{\%rax, (\%rsi)}                       & \textbf{\%rdx, (\%rdi)}               \\
			\hline
			\%rdi             & 0x120          & 0x120                   & 0x120                   & 0x120                         \rightarrow 125 & 0x120 \textcolor{red}{\rightarrow450} \\
			\%rsi             & 0x130          & 0x130                   & 0x130                   & 0x130\textcolor{red}{\rightarrow125}          & 0x130 \rightarrow 125                 \\
			\%rax             & -              & \textcolor{red}{125}    & 125                     & 125                                           & 125                                   \\
			\%rdx             & -              & -                       & \textcolor{red}{450}    & 450                                           & 450                                   \\
			\hline
		\end{array}
	\]
	Es posible que el lector no entienda el por qué de las cuatro instrucciones en vez de tres como en C, esto se debe a que, como hemos visto antes, no se pueden realizar
	operaciones de memoria a memoria, por lo que se necesita un registro para almacenar el valor de una de las direcciones de memoria.
\end{ejemplo}
\subsection{Operaciones}
Veamos algunas de las operaciones más comunes en x86-64, notemos que algunas de ellas además pueden tener
variantes indicando el tamaño de los operandos. Antes de ver las operaciones de forma general, recalcamos
la instrucción \textbf{\textit{leaq}}:
\subsubsection{leaq}
La instrucción (\textit{load effective address}) se utiliza como variante de la instrucción \textit{movq},
la instrucción lee de la memoria la dirección efectiva de un operando y la almacena en un registro, sin acceder a la memoria en sí.
Por ejemplo \textit{leaq 0x8(\%rax), \%rdx} almacena en $\%rdx$ la dirección de memoria $\%rax + 0x8$. Esta instrucción es a menudo utilizada
por los compiladores para optimizaciones.
\begin{ejemplo}
	Veamos una forma para multiplicar un numero por 12 usando leaq
	\begin{center} % Centramos el contenido
		\begin{minipage}{0.3\textwidth}
			\begin{minted}[fontsize=\large]{c}
void m12(long x) 
{
    return x*12;
}
                \end{minted}
		\end{minipage}
		\begin{minipage}{0.69\textwidth}
			\begin{minted}[fontsize=\large]{asm}
swap:
    leaq (%rdi, %rdi, 2), %rax # t = x + 2x
    salq 2, %rax             # return t<<2
                \end{minted}
		\end{minipage}
	\end{center}
\end{ejemplo}

Veamos algunas de las operaciones del tipo aritmetico-lógicas
\begin{center}
    \begin{tabular}{|c|c|c|}
    \hline
    \textbf{Instrucción} & \textbf{Efecto}                       & \textbf{Descripción}     \\ \hline
    \texttt{lea S, D}         & $D \leftarrow \&S$                   & Load effective address   \\ \hline
    \texttt{inc D}            & $D \leftarrow D + 1$                 & Increment                \\ \hline
    \texttt{dec D}            & $D \leftarrow D - 1$                 & Decrement                \\ \hline
    \texttt{neg D}            & $D \leftarrow -D$                    & Negate                   \\ \hline
    \texttt{not D}            & $D \leftarrow \sim D$                & Complement               \\ \hline
    \texttt{add S, D}         & $D \leftarrow D + S$                 & Add                      \\ \hline
    \texttt{sub S, D}         & $D \leftarrow D - S$                 & Subtract                 \\ \hline
    \texttt{imul S, D}        & $D \leftarrow D * S$                 & Multiply                 \\ \hline
    \texttt{xor S, D}         & $D \leftarrow D \oplus S$            & Exclusive-or             \\ \hline
    \texttt{or S, D}          & $D \leftarrow D \vee S$              & Or                       \\ \hline
    \texttt{and S, D}         & $D \leftarrow D \wedge S$            & And                      \\ \hline
    \texttt{sal k, D}         & $D \leftarrow D \ll k$               & Left shift               \\ \hline
    \texttt{shl k, D}         & $D \leftarrow D \ll k$               & Left shift (same as sal) \\ \hline
    \texttt{sar k, D}         & $D \leftarrow D \gg_A k$             & Arithmetic right shift   \\ \hline
    \texttt{shr k, D}         & $D \leftarrow D \gg_L k$             & Logical right shift      \\ \hline
    \end{tabular}
\end{center}

Destacamos una característica peculiar de las operaciones de shift, en el caso de que hagamos un shift a la derecha sobre un
registro de 32 bits \texttt{shr \$16 $\%$edx}, $\%$rdx se rellenará con 0 en los primeros 32 bits también, de esta forma si teniamos $\$-1$
por ejemplo y realizamos \texttt{shr \$16 $\%$edx}, nos quedaremos con 0x0000 0000 0000 FFFF en vez de 0xFFFF FFFF 0000 FFFF.
\subsubsection{Operaciones aritméticas especiales}
Multiplicar dos enteros con/sin signo de 64 bits puede tener un resultado que necesite 128 bits para representarlo
es para ello que se incluyen las siguientes instrucciones especiales:
\begin{center}
    \begin{tabular}{|c|c|c|}
    \hline
        \textbf{Instruction} & \textbf{Effect} & \textbf{Description} \\ \hline
    \texttt{imulq S} & \texttt{R[\%rdx]:R[\%rax] ← S × R[\%rax]} & Signed full multiply \\ \hline
    \texttt{mulq S} & \texttt{R[\%rdx]:R[\%rax] ← S × R[\%rax]} & Unsigned full multiply \\ \hline
    \texttt{cqto} & \texttt{R[\%rdx]:R[\%rax] ← SignExtend(R[\%rax])} & Convert to oct word \\ \hline
    \texttt{idivq S} & 
    \begin{tabular}[t]{@{}l@{}}
        \texttt{R[\%rdx] ← R[\%rdx]:R[\%rax] mod S;} \\ 
        \texttt{R[\%rax] ← R[\%rdx]:R[\%rax] ÷ S}
    \end{tabular} & Signed divide \\ \hline
    \texttt{divq S} & 
    \begin{tabular}[t]{@{}l@{}}
        \texttt{R[\%rdx] ← R[\%rdx]:R[\%rax] mod S;} \\ 
        \texttt{R[\%rax] ← R[\%rdx]:R[\%rax] ÷ S}
    \end{tabular} & Unsigned divide \\ \hline
    \end{tabular}
\end{center}
Además de esto, añadimos ciertas peculiaridades que nos pueden ser útiles en la práctica:

\begin{itemize}
    \item Las instrucciones \texttt{movzbl} y \texttt{movslq} (la \texttt{z} indica \textit{zero extension} y la \texttt{s} \textit{sign extension}) se utilizan para extender cantidades a registros. Lo remarcable es que, mientras que \texttt{movslq} mueve 32 bits a 64 bits con extensión de signo, \texttt{movzlq} no existe ya que al mover 32 bits a 64 bits (\texttt{movl}) ya se realiza una extensión de ceros.
    \item Cuando se tiene el código desensamblado de un programa, podemos ver la dirección en hexadecimal, junto con varios pares de números en hexadecimal y a la derecha las instrucciones. Los pares de números en hexadecimal representan los bytes de la instrucción. Por ejemplo, si tenemos la línea siguiente:
    \begin{center}
        400544: \quad \texttt{e8 08 00 00 00} \quad \texttt{callq 400551 <f>}
    \end{center}
    Podremos ver que \texttt{e8} corresponde con la instrucción \texttt{callq} y los siguientes 4 bytes son la dirección a la que se llama. En este caso, \texttt{400544} + 8 = \texttt{400551}, como se puede ver en la instrucción a la derecha.
\end{itemize}


\section{Control}
Hasta ahora hemos considerado como se comporta el código cuando se ejecuta de forma secuencial, sin embargo,
algunas estructuras de $C$ como los bucles, las estructuras de control, etc. requieren de una ejecución condicional, esto es,
la secuencia de instrucciones a ejecutar depende de una condición.
En x86-64, veremos dos formas de implementar estas estructuras de control más adelante, primero, veamos
los códigos de condición que se utilizan en x86-64:
\subsection{Codigos de condición}
Además de los registros de propósito general, la CPU mantiene un conjunto de bits de estado que se utilizan para
describir atributos de la última operación aritmético-lógica realizada.
Veamos los códigos de condición que vamos a usar:
\begin{itemize}
	\item \textbf{CF (Carry Flag)}: El flag de acarreo, se activa si hay un acarreo en la última operación. Se utiliza para detectar overflow en operaciones sin signo.
	\item \textbf{ZF (Zero Flag)}: El flag de cero, se activa si el resultado de la última operación es cero.
	\item \textbf{SF (Sign Flag)}: El flag de signo, se activa si el resultado de la última operación es negativo.
	\item \textbf{OF (Overflow Flag)}: El flag de desbordamiento, se activa si hay un desbordamiento en la última operación en complemento a 2.
\end{itemize}
Estos códigos de condición se ajustan implicitamente por las operaciones aritméticas, de esta forma, \textit{leaq} no afecta a los códigos de condición.
Además de ajustarse por las operaciones que ya conocemos, existen instrucciones que ajustan explicitamente los códigos de condición sin alterar los registros de propósito general.
\begin{itemize}
	\item El conjunto de instrucciones \textbf{\textit{Compare}} (cmp) ajustan los códigos de condición de acuerdo a las diferencias entre sus dos operandos. Debemos tener en cuenta que en $AT\&T$ la sintaxis
	      es al revés: \textit{cmpq S1, S2} compara $S1$ con $S2$ mediante $S2- S1$.
	\item El conjunto de instrucciones \textbf{\textit{Test}} (test) ajustan los códigos de condición de acuerdo a las operaciones lógicas entre sus dos operandos.
	      Se realiza la operación $S1 \& S2$ y se ajustan los códigos de condición. Es común que se utilicen ambos operandos iguales, por ejemplo, para
	      comprobar si un número es 0.
\end{itemize}
\subsubsection{Consultar códigos de condición}
En vez de consultarlos directamente, se utilizan mediante tres formas:
\begin{enumerate}
	\item Se establece un único byte a 0 o 1 en función de una combinación de los códigos de condición.
	\item Se utilizan para realizar saltos condicionales.
	\item Se utilizan para transferir datos condicionalmente.
\end{enumerate}
Para la primera forma, se utilizan las instrucciones \textit{setcc} siguientes:
\begin{center}
	\begin{tabular}{|c|c|}
		\hline
		\textbf{Instruction} & \textbf{Effect}                                                          \\ \hline
		\texttt{sete D}      & D $\leftarrow$ ZF (Equal / zero)                                         \\ \hline
		\texttt{setne D}     & D $\leftarrow$ $\sim$ ZF (Not equal / not zero)                          \\ \hline
		\texttt{sets D}      & D $\leftarrow$ SF (Negative)                                             \\ \hline
		\texttt{setns D}     & D $\leftarrow$ $\sim$ SF (Nonnegative)                                   \\ \hline
		\texttt{setg D}      & D $\leftarrow$ $\sim$ (SF $\oplus$ OF) \& $\sim$ ZF (Greater, signed >)  \\ \hline
		\texttt{setge D}     & D $\leftarrow$ $\sim$ (SF $\oplus$ OF) (Greater or equal, signed $\geq$) \\ \hline
		\texttt{setl D}      & D $\leftarrow$ SF $\oplus$ OF (Less, signed <)                           \\ \hline
		\texttt{setle D}     & D $\leftarrow$ (SF $\oplus$ OF) $\lor$ ZF (Less or equal, signed $\leq$) \\ \hline
		\texttt{seta D}      & D $\leftarrow$ $\sim$ CF \& $\sim$ ZF (Above, unsigned >)                \\ \hline
		\texttt{setae D}     & D $\leftarrow$ $\sim$ CF (Above or equal, unsigned $\geq$)               \\ \hline
		\texttt{setb D}      & D $\leftarrow$ CF (Below, unsigned <)                                    \\ \hline
		\texttt{setbe D}     & D $\leftarrow$ CF $\lor$ ZF (Below or equal, unsigned $\leq$)            \\ \hline
	\end{tabular}
\end{center}
\begin{ejemplo}
	Veamos un ejemplo de su uso:
	\begin{center}
		\begin{minipage}{0.39\textwidth}
			\begin{minted}[fontsize=\Large]{c}
    long gt(long a, long b)
    {
        return a > b;
    }
            \end{minted}
		\end{minipage}
		\begin{minipage}{0.6\textwidth}
			\begin{minted}[fontsize=\Large]{asm}
    gt:
        cmpq %rsi, %rdi
        setg %al
        movzbq %al, %rax
        ret

            \end{minted}
		\end{minipage}
	\end{center}
	Donde movzbq es una instrucción que significa \textit{ Move with Zero-extend Byte to Long}.
\end{ejemplo}

Veamos ahora cómo se realizan los saltos condicionales.
Estos provocan que la ejecución del programa pueda saltar a otra parte del código en función de los códigos de condición, es muy posible que
en algún momento nos encontremos una instrucción en C de tipo goto, su uso está desaconsejado, sin embargo, por su similitud con el ensamblador,
puede resultar útil en ciertas situaciones. Veamos una tabla con las instrucciones de salto condicional y después un ejemplo de uso.

\begin{center}
	\begin{tabular}{|c|c|c|}
		\hline
		\textbf{Instruction} & \textbf{Jump Condition}                 & \textbf{Description}         \\ \hline
		jmp                  & Label                                   & Direct jump                  \\ \hline
		jmp                  & *Operand                                & Indirect jump                \\ \hline
		je                   & ZF                                      & Equal / zero                 \\ \hline
		jne                  & $\sim$ZF                                & Not equal / not zero         \\ \hline
		js                   & SF                                      & Negative                     \\ \hline
		jns                  & $\sim$SF                                & Nonnegative                  \\ \hline
        jg                   & $\sim$(SF \^{} OF) $\&$ $\sim$ZF & Greater (signed >)           \\ \hline
        jge                  & $\sim$(SF \^{} OF)                  & Greater or equal (signed >=) \\ \hline
        jl                   & SF \^{} OF                          & Less (signed <)              \\ \hline
        jle                  & (SF \^{}  OF) $|$ ZF              & Less or equal (signed <=)    \\ \hline
        ja                   & $\sim$CF $\&$ $\sim$ZF               & Above (unsigned >)           \\ \hline
		jae                  & $\sim$CF                                & Above or equal (unsigned >=) \\ \hline
		jb                   & CF                                      & Below (unsigned <)           \\ \hline
		jbe                  & CF $|$ ZF                            & Below or equal (unsigned <=) \\ \hline
	\end{tabular}
\end{center}
\begin{ejemplo}
	Veamos ahora un ejemplo de cómo hacer por ejemplo un valor absoluto de una diferencia:
	\begin{center}
		\begin{minipage}{0.39\textwidth}
			\begin{minted}[fontsize=\Large]{c}
    long abs(long x, long y)
    {
        if (x > y)
            return x - y;
        else
            return y - x;
    }
            \end{minted}
		\end{minipage}
		\begin{minipage}{0.6\textwidth}
			\begin{minted}[fontsize=\Large]{asm}
        abs:
            cmpq %rsi, %rdi
            jle .L1
            movq    %rdi, %rax
            subq    %rsi, %rax
            ret
        .L1:
            movq %rsi, %rax
            subq %rdi, %rax
            ret
			\end{minted}
		\end{minipage}
	\end{center}
	Si tratamos de \href{https://godbolt.org/z/aj7f7os8d}{replicar} el ejemplo anterior, veremos que en general, el compilador no produce este código,
	esto se debe a las optimizaciones que el compilador aplica, evitando en la medida de lo posible los saltos condicionales, ya que muchas veces,
	es mejor calcular ambas ramas, a tener que predecir y calcular cual de las ramas hay que seguir.
	Para obtener este mismo código, deberemos especificar la opción de compilación \textit{-fno-if-conversion}.
\end{ejemplo}
Nos detenemos brevemente en las optimizaciones mencionadas, lo que hace el compilador en vez de saltar ramas es asignar valores a las variables
en función de un test: \texttt{let val = test ? x : y}. \textit{gcc} intenta utilizar este tipo de optimizaciones siempre que sea seguro.
De esta forma, el código anterior se traduciría a:
\begin{center}
	\begin{minted}[fontsize=\Large]{asm}
        abs:
            movq %rdi, %rax # x
            subq %rsi, %rax # x - y
            movq %rsi, %rdx 
            subq %rdi, %rdx # y - x
            cmpq %rsi, %rdi # x : y
            cmovle %rdx, %rax # si x <= y, rax = rdx
    \end{minted}
\end{center}
Esto se verá como una optimización en la asignatura de \textit{Arquitectura de Computadores}, sin embargo, adelantamos
que por ejemplo \texttt{ let val = test ? p : *q} no se puede optimizar de la misma forma, ya que pudiera tener efectos no deseables.
Notemos además que este tipo de optimizaciones solo tienen sentido cuando los cálculos son sencillos, en cuyo caso es preferible calcular las
dos ramas a predecir cual de ellas seguir. Es por este tipo de opciones que \textit{gcc} es posible que genere un código mas lento cuando se activan
las optimizaciones \textit{-O3} ya que si la rama era muy predecible, la predicción de saltos es mas rápida que calcular las dos ramas.
Tampoco se debe utilizar en casos donde una rama u otra tiene efectos colaterales, ya que eso pudiera cambiar el comportamiento del programa.
\subsection{Bucles}
La mayoría de lenguajes ofrecen varias formas de implementar bucles: \textit{ do-while, while, for, loop, etc.} En x86-64, no existen instrucciones
similares, en su lugar, se combinan diferentes tests y saltos condicionales para implementar los bucles.
\subsubsection{Bucles Do-While}
En C, un bucle \textit{do-while} se implementa de la siguiente forma:
\begin{center}
	\begin{minipage}{0.45\textwidth}
		\begin{minted}[fontsize=\Large]{c}
long pcount(long x)
{
    long result = 0;
    do
    {
        result += x & 0x1;
        x >>= 1;
    } while (x);
    return result;
}
        \end{minted}
		\descriptionlabel{Versión C}
	\end{minipage}
	\begin{minipage}{0.45\textwidth}
		\begin{minted}[fontsize=\Large]{c}
long pcount(long x)
{
    long result = 0;
    loop:
        result += x & 0x1;
        x >>= 1;
        if(x) goto loop;
    return result;
}
        \end{minted}
		\descriptionlabel{Versión goto}

	\end{minipage}
\end{center}
En este ejemplo se ha hecho un clásico programa popcount que cuenta el número de bits a 1 en un número.
El nombre \textit{popcount} viene de \textit{population count}, que es el nombre que se le da a esta operación.
Como ya sabemos, el efecto del loop es ejecutar el cuerpo del bucle, evaluar la condición y si el resultado de la condición es distinto de 0,
seguir ejecutando el bucle. Notemos que se ejecuta al menos una vez.
\subsubsection{Bucles While}
\begin{ejemplo}
	Veamos ahora cómo se implementa un bucle \textit{while}:
	\begin{center}
		\begin{minipage}{0.45\textwidth}
			\begin{minted}[fontsize=\Large]{c}
long pcount(long x)
{
    long result = 0;
    while(x) {
        result += x & 0x1;
        x >>= 1;
    }

    return result;
}
\end{minted}
			\descriptionlabel{Versión C}
		\end{minipage}
		\begin{minipage}{0.45\textwidth}
			\begin{minted}[fontsize=\Large]{c}

long pcount(long x)
{
        long result = 0;
        goto test;
    loop:
        result += x & 0x1;
        x >>= 1;
    test:
        if (x) goto loop;
        return result;
}
        \end{minted}
			\descriptionlabel{Versión goto}

		\end{minipage}
	\end{center}
\end{ejemplo}
En general, while evalúa la condición antes de ejecutar el bucle, por lo que si la condición no se cumple, el bucle no se ejecuta.
\subsubsection{Bucles For}
Sabemos ya que en general un bucle \textit{for} se compone de tres partes, la inicialización, la condición y la actualización.
No adjuntamos el código en C ya que es trivial sabiendo los anteriores, notemos que un bucle \textit{for} es equivalente a un bucle \textit{while} de la siguiente forma:
\textit{for (init; test; update) body} es equivalente a \textit{init; while(test) {body; update;}}
Por último, veamos las secuencias switch.
\subsection{Secuencias Switch}
Una secuencia switch proporciona una ramificación condicional múltiple basada en el valor de una expresión entera.
Son muy útiles cuando trabajamos con tests que pueden tener una gran cantidad de resultados, haciendo el código más legible
y eficiente. Este último punto puede sorprender, por ello, vamos a desarrollarlo.
\\ \\
Una \textit{tabla de salto} es una estructura de datos en forma de array donde $i$ es la dirección de un segmento de codigo que ejecuta la acción
que el programa debería seguir cuando el valor de la expresión es $i$. La diferencia entre multiples \textit{if-else} y \textit{switch} radica en
que el tiempo que se tarda en ejecutar un \textit{switch} es constante, mientras que en un \textit{if-else} es lineal en función del número de casos.

\begin{center}
	\begin{minipage}{0.45\textwidth}
		\centering
		\textbf{Código en C}

		\begin{minted}[fontsize=\large]{c}

void switch_eg(long x) {
    switch (x) {
        case 1:
            printf("Code 1\n");
            break;
        case 2:
            printf("Code 2\n");
            break;
        case 3:
            printf("Code 3\n");
            break;
        default:
            printf("Code 4\n");
    }
}

    \end{minted}

	\end{minipage}
	\hfill
	\begin{minipage}{0.45\textwidth}
		\centering
		\textbf{Assembly Code}

		\begin{minted}[fontsize=\large]{asm}
cmpq $3, %rdi
ja .L4 # Si x > 3, Default
jmp *.L1(,%rdi,8) # Dirección L1 + 8x 
.L1:
.L2:
.L3:
.L4 # Default
        \end{minted}

		\vspace{1em}
		\textbf{Branch Table (Jump Table)}

		\begin{tabular}{|c|c|}
			\hline
			\textbf{Índice} & \textbf{Dirección de Salto} \\
			\hline
			1               & Code Block 1                \\
			\hline
			2               & Code Block 2                \\
			\hline
			3               & Code Block 3                \\
			\hline
			Default         & Code Block 4                \\
			\hline
		\end{tabular}
	\end{minipage}
\end{center}
Como podemos ver, la tabla de saltos comienza en este caso en .L1, para acceder
a el resto de bloques de código, se suma 8 veces el valor de x, ya que cada dirección
ocupa 8 bytes. Si x es mayor que 3, se salta a la dirección .L4, que es el bloque por defecto.
Si por ejemplo en vez de hacer break, se quisiera hacer un fall through, obtendríamos algo
del siguiente estilo:
\begin{center}
	\begin{minipage}{0.45\textwidth}
		\begin{minted}[fontsize=\large]{c}
long w = 1;
switch(x) {
    ...
  case 2:
    w = y/z;
    /* Fall through */ 
case 3:
    w += z;
    break;
    ...
}
\end{minted}
	\end{minipage}
	\begin{minipage}{0.45\textwidth}
		\begin{equation*}
			\%rdi \rightarrow x \quad \%rsi \rightarrow y \quad \%rdx \rightarrow z
		\end{equation*}
		\begin{minted}[fontsize=\large]{asm}
.L5: # Case 2
    movq %rsi, %rax
    cqto # Convert Quad To Oct
    idivq %rdx 
    jmp .L6
.L9: # Case 3
    movl $1, %eax
.L6: # Merge
    addq %rcx, %rax # w += z
    ret 

\end{minted}
	\end{minipage}
\end{center}

\section{Procedimientos}
Los procedimientos son una abstracción clave en el desarrollo de software, permiten dividir un programa en pequeñas partes
donde cada una de ellas proporciona una funcionalidad concreta al proporcionarle una serie de argumentos.
Pueden ser llamados desde cualquier parte del programa, lo que permite reutilizar el código y hacerlo más modular. Los procedimientos vienen
de distintas formas en cada lenguaje: funciones, métodos, subrutinas, lambdas, etc, pero todos ellos comparten ciertos elementos comunes.
\\ \\
Existen muchos atributos que hay que manejar cuando se proporciona soporte en lenguaje ensamblador para procedimientos, para verlos, supongamos un proceso
$P$ que llama a un procesor $Q$, que se ejecuta y tras esto vuelve a $P$, estas acciones incluyen:
\begin{itemize}
	\item \textbf{Pasar Control}: El PC tiene que fijarse en la dirección de comienzo de $Q$, y tras esto, cuando $Q$ termina, el PC tiene que volver a la dirección de la instrucción que sigue a la llamada.
	\item \textbf{Parámetros}: $P$ tiene que pasar los parámetros a $Q$, y $Q$ tiene que devolver un valor a $P$.
	\item \textbf{Reserva y Limpieza de Memoria}: $Q$ puede necesitar reservar memoria para sus propias variables, y tras la ejecución, limpiarla.
\end{itemize}
La implementación de estos procedimientos en x86-64 involucra una combinación de instrucciones especiales y convenciones de cómo usar los recursos de la máquina.
Detrás de esto hay un gran esfuerzo por minimizar el overhead que involucra una llamada, esto provoca que al final, solo se ejecuta lo mínimo indispensable mencionado
antes. Veamos ahora diferentes mecanismos paso por paso, comenzando por la pila:
\subsection{Stack (Pila)}
Es una región de memoría gestionada con aspectos de LIFO, es decir, el último elemento en entrar es el primero en salir, crece hacia direcciones inferiores, en x86-64
el registro $\%rsp$ contiene la dirección mas baja de la pila. Los datos, pueden ser accedidos y extraidos de la pila usando las instrucciones \textit{pushq} y \textit{popq}.
Por ejemplo para reservar espacio para datos sin valor inicial, se puede decrementar la pila la cantidad necesaria, análogamente, se puede liberar el espacio incrementando el stack pointer.
En x86-64, cuando un procedimiento requiere más espacio de lo que puede guardar en registros, reserva espacio en el stack. A esta región se la llama \textit{procedure space} para las variables locales.
En interés del espacio y la eficiencia, los procedimientos de x86-64 reservan solo el espacio necesario.

\subsection{Pasando el Control}
Pasar el control de una función $P$ a otra $Q$ requiere únicamente de poner el PC en la dirección de inicio de $Q$,
sin embargo, cuando $Q$ termina, el PC tiene que tener alguna manera de saber dónde debe reanudar la ejecución de $P$, esta información se guarda en las máquinas
x86-64 usando la instrucción call $Q$. Esta instrucción añade una dirección $A$ al stack y pone el PC en el comienzo de $Q$. La dirección $A$ se denomina \textit{return address} y se calcula como
la instrucción que sigue inmediatamente a call. La instrucción contraria ret, elimina una dirección $A$ de la pila y pone el PC en $A$.
(Notemos que estas instrucciones nos las encontraremos como callq y retq, donde el sufijo $q$ simplemente enfatiza que son versiones de x86-64 no de IA32).
Veamos un esquema con la estructura general del stack:


\begin{center}
	\includesvg[width=0.6\textwidth]{Imagenes/stack}
	\label {fig:stack}
\end{center}
\subsection{Pasando los datos}
Además de pasar el control a un procedimiento cuando es invocado, y de vuelta cuando termina, las llamadas a procedimientos
pueden involucrar pasar datos como argumentos, además un procedimiento puede devolver también un valor. En x86-64, la mayoria de
este flujo de datos es llevada a cabo mediante registros, ya hemos visto numerosos ejemplos en los que $\%rdi$ y $\%rsi$ se pasan como
parámetros y valores que se devuelven mediante $\%rax$. En x86-64, se pueden pasar hasta 6 argumentos mediante registros, que se usan en un orden
específico con el nombre del registro en función del tamaño del dato que contiene.
Cuando una función tiene más de 6 argumentos, los argumentos restantes se pasan en la pila tal que el argumento 7 se coloca abajo, el 8 encima, y así sucesivamente.
\subsection{Gestionando datos locales}
Hasta ahora, la mayoría de ejemplos que hemos visto no han requerido de variables locales, ya sea por optimización del compilador o por la sencillez del código,
sin embargo, sabemos que muchas veces es necesario almacenar variables locales.
Normalmente, en lenguajes basados en pila, o que soportan la recursión (e.g C, C++, Java, Rust, etc), las variables locales se almacenan en la pila.
Esta pila, se reserva en marcos (allocated in frames), que son bloques de memoria que contienen las variables locales de un procedimiento, estos marcos de memoria
se organizan tal que contienen información de retorno, almacenamiento local si es necesario y espacio temporal si es necesario, además hay un puntero de marco opcional:
$\%rbp$. Veamos un ejemplo de cadena de llamadas a procedimientos:
% \begin{center}
% 	\begin{tikzpicture}[>=Stealth, scale=1.5]
%
% 		\matrix (m1) [matrix of nodes, nodes={draw, minimum size=1cm}, column sep=0cm] {
% 			\node[fill=green!20] {\phantom{a}}; \\
% 			\node[fill=green!20] {\phantom{a}}; \\
% 			\node[fill=green!20] {\phantom{a}}; \\
% 			\node[fill=green!20] {\phantom{a}}; \\
% 		};
%
% 		\matrix (m2) [matrix of nodes, right=of m1, nodes={draw, minimum size=1cm}, column sep=0cm] {
% 			\node[fill=green!20] {\phantom{a}}; \\
% 			\node[fill=green!40] {fn A};                    \\
% 			\node[fill=green!20] {\phantom{a}}; \\
% 			\node[fill=green!20] {\phantom{a}}; \\
% 		};
%
% 		\draw[->] (m1) -- (m2);
%
% 		\matrix (m3) [matrix of nodes, right=of m2, nodes={draw, minimum size=1cm}, column sep=0cm] {
% 			\node[fill=green!20] {\phantom{a}}; \\
% 			\node[fill=green!40] {fn A};                    \\
% 			\node[fill=green!70] {fn B};                    \\
% 			\node[fill=green!0] {\phantom{a}};  \\
% 		};
%
% 		\draw[->] (m2) -- (m3);
% 		\matrix (m4) [matrix of nodes, right=of m3, nodes={draw, minimum size=1cm}, column sep=0cm] {
% 			\node[fill=green!20] {\phantom{a}}; \\
% 			\node[fill=green!40] {fn A};                    \\
% 			\node[fill=green!70] {fn B};                    \\
% 			\node[fill=green!70] {fn B};                    \\
% 		};
%
% 		\draw[->] (m3) -- (m4);
%
% 		\matrix (m5) [matrix of nodes, right=of m4, nodes={draw, minimum size=1cm}, column sep=0cm] {
% 			\node[fill=green!20] {\phantom{a}}; \\
% 			\node[fill=green!40] {fn A};                    \\
% 			\node[fill=green!70] {fn B};                    \\
% 			\node[fill=green!20] {\phantom{a}}; \\
% 		};
%
% 		\draw[->] (m4) -- (m5);
% 		\matrix (m6) [matrix of nodes, right=of m5, nodes={draw, minimum size=1cm}, column sep=0cm] {
% 			\node[fill=green!20] {\phantom{a}}; \\
% 			\node[fill=green!40] {fn A};                    \\
% 			\node[fill=green!20] {\phantom{a}}; \\
% 			\node[fill=green!20] {\phantom{a}}; \\
% 		};
%
% 		% Flecha entre la cuarta y quinta tabla
% 		\draw[->] (m5) -- (m6);
% 	\end{tikzpicture}
% \end{center}
Donde en cada uno de los pasos, $\%rsp$ apunta en la función más abajo existente y $\%rbp$ apunta al marco de la función inmediata anterior.
Los registros de programa actuan como un único recurso compartido por todos los procedimientos, aunque un solo procedimiento puede estar activo en un momento dado,
tenemos que asegurarnos de que cuando un procedimiento (caller) llama otro procedimiento (callee), el callee no sobreescriba los registros del caller. Veamos un ejemplo de esto:
\begin{center}
	\begin{minipage}{0.45\textwidth}
		\begin{minted}[fontsize=\large]{asm}
funcA:
    movq $15213, %rdx
    call funcB
    addq %rdx, %rax 
    ret
\end{minted}
	\end{minipage}
	\begin{minipage}{0.45\textwidth}
		\begin{minted}[fontsize=\large]{asm}
funcB: 
    subq $18213, %rdx
    ret 
    \end{minted}

	\end{minipage}
\end{center}
Como podemos ver, pudiera pasar que $\%rdx$ fuese sobreescrito en la función $funcB$, para evitar esto,
existen convenciones para coordinar el uso de registros.
Los registros $\%rbx, \%rbp, \%r12-\%r15$ se denominan \textit{callee-saved} ya que el callee debe guardar su valor si los modifica.
Para este fin, el callee puede o bien no modificar los registros, o guardando su valor en la pila y restaurándolo antes de salir, a este espacio se le llama
\textit{Saved Registers}, podemos verlo en el diagrama \ref{fig:stack}.
El resto de registros, excepto $\%rsp$, se clasifican como \textit{Caller-saved}, lo que
significa que pueden ser modificados por cualquier función, el nombre viene de un procedimiento $P$ que llama a un $Q$, como
$Q$ puede modificar libremente estos registros, $P$ debe guardar su valor si los necesita.
\subsection{Recursividad}
Las convenciones que hemos visto hasta ahora son suficientes para permitir a los procedimientos de x86-64 llamarse a sí mismos.
Cada llamada a procedimiento tiene su propio espacio en el stack, luego las variables locales de las múltiples llamadas no interfieren entre ellas.
Veamos un ejemplo del código C y ensamblador generado para una función recursiva:
\begin{center}
	\begin{minipage}{0.45\textwidth}
		\begin{minted}[fontsize=\large]{c}
long pcount_r(unsigned long x) {
    if (x == 0) {
        return 0;
    } else {
        return (x & 1) 
            + pcount_r(x >> 1);
    }
}
\end{minted}
	\end{minipage}
	\begin{minipage}{0.45\textwidth}
		\begin{minted}[fontsize=\large]{asm}
pcount_r:
    movl $0, %eax
    testq %rdi, %rdi
    je .L2
    pushq %rbx 
    movq %rdi, %rbx
    andl $1, %ebx
    shrq %rdi
    call pcount_r
    addq %rbx, %rax
    popq %rbx
  .L2
    rep; ret # Para branch prediction 
  \end{minted}
	\end{minipage}
\end{center}
Veamos más detenidamente lo que está pasando:
%Error 3
% \begin{itemize}
% 	\item Las primeras tres instrucciones sirven para saber cuando acaba la recursión, si x es 0, se devuelve 0.
% 	\item El pushq está ya que se va a modificar $\%rbx$ y se necesita restaurar su valor al salir de la función.
% 	\item Desde movq hasta shrq se calcula x $\&$ 1 y x >> 1.
% 	\item x >> 1 se pasa como argumento a la función $pcount_r$ (llamada recursiva)
% 	\item Se suma el resultado de la llamada recursiva con el valor de x \& 1.
% \end{itemize}
\section{Datos}
% Marker [DE AQUI HACIA ABAJO TODO CORRECTO]
\subsection{Arrays}
Los arrays en C se usan para almacenar colecciones de datos escalares en variables más grandes, la implementación particular
de C de los arrays es muy directa, una característica importante es que en C podemos manipular punteros a elementos de un array
y realizar operaciones con ellos. Los compiladores son extremadamente buenos simplificando el cálculo de direcciones 
usado cuando accedemos a un array, es por ello que la correspondencia entre C y ensamblador puede llegar a ser dificil de descifrar.
\subsubsection{Principios Básicos}
Sea $T$ un tipo de dato cualquiera, y $N$ una constante, consideramos la declaración $T A[N]$, entonces $A$ es un array de $N$ elementos.
La declaración reserva una región contigua de memoria de $T \cdot N$ bytes e introduce un identificador $A$ que se usa como puntero
al comienzo de esta región de memoria. Los elementos de este array se acceden mediante el uso de un índice entre 0 y N-1.
\begin{ejemplo}
    Consideramos \texttt{int arr[5]}, veamos los siguientes ejemplos:
    \begin{center}
        \includesvg[width=0.5\textwidth]{Imagenes/array}
        \end{center}
    \begin{itemize}
        \item \texttt{arr[4]} (int) se refiere al quinto elemento del array (valor 3).
        \item \texttt{arr} (int *) comienzo del array (posición arr)
        \item \texttt{arr+1} (int *) comienzo del array + 1 (posición arr + 4) 
        \item \texttt{\&arr[2]} (int *) dirección de val[2] (posición arr + 8)
        \item \texttt{arr[5]} (int) fuera de los límites del array, \href{https://en.cppreference.com/w/c/language/behavior}{undefined behavior.}
        \item \texttt{*(arr+1)} (int) valor de la dirección val + 1 (valor 5)
    \end{itemize}
\end{ejemplo}
Destacamos que declaramos arrays inmediatamente después de otro, no está garantizado que estos se almacenen en memoria de forma contigua. Así como hemos podido ver en el ejemplo anterior,
que C permite operaciones con los punteros, así, si $p$ es un puntero a un array de tipo $T$ y la dirección de memoria de $p$ es $x_p$, entonces la expresión $p + i$ tiene el valor \texttt{$x_p$ +  $T \cdot i$ = $i \cdot$ sizeof(T)}.
Recordamos que los operadores $\& $ y $*$ permiten la generación y dereferenciación de punteros respectivamente.
\begin{ejemplo}
    Ejemplo de acceso a elementos de array:
    \begin{center}
        \begin{minipage}{0.45\textwidth}
            \begin{minted}[fontsize=\large]{c}
int get_num(int *arr, int i) 
{
    return arr[i];
}
        \end{minted}
        \end{minipage}
        \begin{minipage}{0.45\textwidth}
        \begin{minted}[fontsize=\large]{asm}
    get_num:
        movslq  %esi, %rsi
        movl (%rdi, %rsi, 4), %eax
        ret
            \end{minted}
        \end{minipage}
    \end{center}
    Recordamos que \texttt{movslq} es una instrucción que significa \textit{Move with Sign-extend Long to Quad}, se utiliza ya que para hacer 
    operaciones con la dirección de memoria, es necesario que los registros sean de 64 bits.
\end{ejemplo}

\subsubsection{Arrays Anidados}
Los principios de los arrays multidimensionales son los mismos que los de los arrays unidimensionales, de esta manera, la declaración
\texttt{T A[R][C]} declara un array de $R$ elementos, cada uno de los cuales es un array de $C$ elementos. En general, un array bidimensional de elementos de tipo T. De otra forma,
podemos decir que tiene $R$ filas y $C$ columnas. Notemos que si el tamaño de $T$ es $K$ bytes, entonces el tamaño de $A$ es $ C \cdot R \cdot K$ bytes.
La disposición de los elementos en memoria es por filas, \textit{(row-major-order)}:
\begin{center}
\includesvg[width=0.8\textwidth]{Imagenes/array2d.svg}
\end{center}
De esta forma, si tenemos un array \texttt{T A[R][C]} el elemento 
\texttt{T[i][j]} está en la posición de memoria \texttt{$\&$A[i][j] = $x_A$ + L(C $\cdot i + j$)}.
\begin{ejemplo}
    Veamos un ejemplo en el que el flag --no-pie debe ser considerado
    \begin{center}
        \begin{minipage}{0.45\textwidth}
            \begin{minted}[fontsize=\large]{c}
    int *get_row(int index) 
    {
        return arr[index];
    }
            \end{minted}
        \end{minipage}
        \begin{minipage}{0.45\textwidth}
            \begin{minted}[fontsize=\large]{c}
#define ROWS 4
// Zip_dig -> Array de 5 elementos
zip_dig arr[ROWS] = {
    {1, 2, 7, 4, 5},
    {6, 7, 2, 9, 1},
    {1, 9, 3, 1, 5},
    {1, 7, 8, 9, 2}};
            \end{minted}
        \end{minipage}
        \vspace{0.5cm}
        \vspace{0.5cm}
        \begin{minipage}{0.45\textwidth}
            \begin{minted}[fontsize=\large]{asm}
    get_num:
        leaq (%rdi, %rdi, 4), %rdx
        leaq arr(,%rdx, 4), %rax
        ret
            \end{minted}
        \end{minipage}
        \begin{minipage}{0.45\textwidth}
            \begin{minted}[fontsize=\large]{asm}
    get_num:
        movslq %edi, %rdi
        leaq arr(%rdi ,%rdi, 4), %rdx
        leaq arr(%rip), %rax
        leaq (%rax, %rdx, 4), %rax
        ret
        \end{minted}
        \end{minipage}
    \end{center}
    A la izquierda, como se usa --no-pie que desactiva el \textit{Position Independent Code}, la dirección de memoria es conocida en tiempo de compilación, por lo que
    se puede acceder directamente con $ 5 \cdot \text{index}$ por ser 5 elementos por fila y $4 \cdot \text{index}$ por ser un tipo de dato entero de 4 bytes.
    En el segundo caso, el direccionamiento es relativo a contador de programa $\%rip$. Notemos que --no-pie resulta más eficiente sin embargo es más vulnerable a ataques de buffer overflow.
\end{ejemplo}
\subsubsection{Array Multinivel}
De manera similar a los arrays bidimensionales, tenemos los arrays multinivel, que tienen la peculiaridad de que 
se tiene un array de punteros a arrays, de esta forma, para acceder a un dato son necesarios dos accesos a memoria, uno para acceder al puntero y otro para acceder al dato.
\begin{ejemplo}
    Un ejemplo de array multinivel en C: 
    \begin{center}
        \begin{minted}{c}
zip_dig cmu = {0, 5, 2, 7, 3};
zip_dig mit = {1, 2, 3, 4, 5};
zip_dig ucb = {9, 3, 7, 5, 1};

# define UCOUNT 3
int *univ[UCOUNT] = {cmu, mit, ucb};
        \end{minted}
    \end{center}
\end{ejemplo}
\subsection{Estructuras}
C permite dos formas de definir tipos de datos compuestos: \textit{struct} que agrupan múltiples datos en un 
solo objeto y \textit{union} que permiten que un objeto sea referenciado usando diferentes tipos de datos. Veamos en detalle las estructuras.
\\ \\ 
El struct en C crea un tipo de dato que agrupa múltiples variables que pueden ser de diferentes tipos en un solo objeto,
los distintos elementos de la estructura se referencian mediante nombres, con la sintaxis \texttt{struct.field, struct->field}.
El struct se parece a los arrays en el sentido de que se garantiza que los datos de un mismo struct están contiguos en memoria.
Es el compilador el que mantiene información sobre cada estructura para indicar el offset de cada campo. El tamaño de los 
structs es el tamaño de todos sus campos sumados, teniendo en cuenta que es posible que el compilador reserve espacio adicional tras un 
tipo de dato para alinear la memoria.
\begin{ejemplo}
    Veamos un ejemplo de struct en C y cómo se distribuye su memoria:
    \begin{center}
        \begin{minipage}{0.45\textwidth}
            \begin{minted}[fontsize=\large]{c}
        struct rec {
            int a[4];
            size_t i;
            struct rec *next;
        }
            \end{minted}
        \end{minipage}
        \begin{minipage}{0.48\textwidth}
            \includesvg[width=0.8\textwidth]{Imagenes/struct}
        \end{minipage}
    \end{center}
        Notemos que si en el struct el array de enteros fuese de tamaño 3, el tamaño del struct se mantendría en 32
        B por el alineamiento de la memoria. El puntero $r$ indica el primer byte de la estructura, se accede
        mediante el offset, que como hemos dicho lo determina el compilador en tiempo de compilación.
        \begin{center}
        \begin{minipage}{0.48\textwidth}
            \begin{minted}[fontsize=\large]{c}
    void set_i(struct rec *r, size_t val)
    {
        r->i = val;
    }
            \end{minted}
        \end{minipage}
        \begin{minipage}{0.45\textwidth}
            \begin{minted}[fontsize=\large]{asm}
        set_i:
            movq %rsi, (%rdi, 16)
            ret
            \end{minted} 
        \end{minipage}
    \end{center}
\end{ejemplo}
Hemos introducido ya el alineamiento de los datos, detallamos los principios que rigen su funcionamiento:
\begin{itemize}
    \item Si el tipo de dato primitivo requiere $K$ bytes, la dirección debe ser múltiplo de $K$.
    \item El alineamiento es requisito en algunas maquinas, en x86-64 es recomendado por motivos de eficiencia (Es ineficiente cargar o almacenar datos que cruzan fronteras de quad word y
        la memoria virtual es delicada cuando un dato se extiende a 2 páginas).
    \item En la mayoría de lenguajes compilados (C, C++, Rust, etc) es el compilador el que inserta huecos en la estructura para alinear los campos.
    \item El alineamiento se hace:
        \begin{itemize}
            \item Para el struct, se toma el máximo valor miembro para alinear, esto se llama alineamiento externo.
            \item Para los datos primitivos, se toma el tamaño del dato, es decir se añade padding antes del dato miembro si fuese necesario para que se encuentre en una posición múltiplo de su tamaño.
        \end{itemize}
\end{itemize}
El compilador no reordena los datos con el fin de ahorrar espacio, sin embargo
el programador debería considerar siempre cómo ordenar los datos para ocupar el menor espacio posible, esto es 
especialmente importante ya que a diferencia de otras optimizaciones que el compilador hace por nosotros, esta no se puede hacer ya que rompería la coherencia y predictibilidad del código.
\subsection{Uniones}
Las uniones proporcionan una manera de circunvalar el sistema de tipos de C, permitiendo a un solo objeto ser referenciado por múltiples tipos. La sintaxis de la unión es idéntica a la de un struct, pero 
sus semánticas son diferentes. En una unión, todos los campos comparten la misma dirección de memoria, lo que significa que si se modifica un campo, se modifica el resto, para hacer esto, el compilador reserva
tamaño de acuerdo al elemento más grande.
\begin{ejemplo}
    Veamos un ejemplo de una unión en comparación con un struct:
    \begin{center}
        \begin{minipage}{0.45\textwidth}
            \begin{minted}[fontsize=\large]{c}
    union U1 {
        char c;
        int i[2];
        double v;
    } *up;
    \end{minted}
    \end{minipage}
    \begin{minipage}{0.45\textwidth}
        \begin{minted}[fontsize=\large]{c}
    struct S1 {
        char c;
        int i[2];
        double v;
    } *sp;
        \end{minted}
    \end{minipage}
\end{center} 
Si compilamos ambos códigos en una máquina Linux en x86-64, el offset y tamaño total de las estructuras sería el siguiente:
\begin{center}
    \begin{tabular}{ccccc}
        Tipo & c & i & v & Total Size \\ 
        \hline 
        U1 & 0 & 0 & 0 & 8 \\ 
        S1 & 0 & 4 & 16 & 24 \\ 
    \end{tabular}
\end{center}
Este gran poder conlleva una gran responsabilidad, ya que el programador debe ser consciente de que 
si usa datos como \texttt{float} e \texttt{int} en la misma unión, 
e intenta interpretar los bits del entero como si fuese un float,
el resultado no será el esperado, ya que no tendremos el mismo efecto que el cast (float) sobre el entero.
\end{ejemplo}
Recordamos que las palabras se almacenan en memoria como un número N de bytes, en esta organización, el byte mas 
significativo puede ser el que está en la dirección más baja (Big Endian) o en la dirección más alta (Little Endian).
Esta diferencia puede causar problemas al intercambiar datos binarios entre máquinas, también mencionamos Middle endian, que son  
aquellas arquitecturas que se pueden configurar de ambas maneras, aparece en ARM, sin embargo, la mayoría de las máquinas modernas son Little Endian.
\\ \\ 
Para finalizar este capítulo, presentamos un resumen de los tipos compuestos en C:
\begin{itemize}
    \item \textbf{Arrays}
        \begin{itemize}
            \item Reservan memoria contigua para almacenar elementos
            \item Se usa aritmética de indexación para localizar elementos individuales
            \item Puntero al primer elemento
            \item No comprueban límites
        \end{itemize}
\item \textbf{Structs}
    \begin{itemize}
        \item Reserva de una sola región de memoria, los campos en el mismo orden declarado 
        \item Se accede usando desplazamientos determinados por el compilador
        \item Alineamiento interno y externo mediante relleno
    \end{itemize}
\item \textbf{Unions}
    \begin{itemize}
        \item Declaraciones superpuestas
        \item Proveen de una manera para saltarse el sistema de tipos de C
    \end{itemize}
\end{itemize}


    \fancyhead[R]{\helv \nouppercase{\rightmark}}
    \chapter{Relaciones de Problemas}
    \section{Introducción a la Computación}


\begin{ejercicio}
    Sea la gramática $G=\left(V,T,P,S\right)$ dada por:
    \begin{align*}
        V &= \{S, X, Y\} \\
        T &= \{a,b\} \\
        S &= S
    \end{align*}
    \begin{enumerate}
        \item Describe el lenguaje generado por la gramática teniendo en cuenta que $P$ viene descrito por:
        \begin{align*}
            S &\rightarrow XYX \\
            X &\rightarrow aX \mid bX \mid \veps \\
            Y &\rightarrow bbb
        \end{align*}

        Sea $L=\{ubbbv\mid u,v\in\{a,b\}^\ast\}$. Demostraremos mediante doble inclusión que $L=\cc{L}(G)$.
        \begin{description}
            \item[$\subset)$] Sea $w\in L$. Entonces, $w=ubbbv$ con $u,v\in\{a,b\}^\ast$. Veamos que
            $S \stackrel{\ast}{\Longrightarrow} w$:
            \begin{equation*}
                S \Longrightarrow XYX \Longrightarrow XbbbX
            \end{equation*}

            Además, es fácil ver que la regla de producción $X\rightarrow aX \mid bX \mid \veps$ nos permite generar cualquier palabra $u\in\{a,b\}^\ast$. Por tanto, tenemos que $X \stackrel{\ast}{\Longrightarrow} u$ y $X \stackrel{\ast}{\Longrightarrow} v$; teniendo así que $S \stackrel{\ast}{\Longrightarrow} ubbbv$.

            \item[$\supset)$] Sea $w\in\cc{L}(G)$. Veamos la forma de $w$:
            \begin{equation*}
                S \Longrightarrow XYX \Longrightarrow XbbbX \Longrightarrow ubbbv \mid u,v\in\{a,b\}^\ast
            \end{equation*}
            donde en el último paso hemos empleado lo visto en el apartado anterior de la regla de producción $X\rightarrow aX \mid bX \mid \veps$. Por tanto, $w\in L$.
        \end{description}


        \item Describe el lenguaje generado por la gramática teniendo en cuenta que $P$ viene descrito por:
        \begin{align*}
            S &\rightarrow aX \\
            X &\rightarrow aX \mid bX \mid \veps
        \end{align*}

        Sea $L=\{au \mid u\in\{a,b\}^\ast\}$. Demostraremos mediante doble inclusión que $L=\cc{L}(G)$.
        \begin{description}
            \item[$\subset)$] Sea $w\in L$. Entonces, $w=au$ con $u\in\{a,b\}^\ast$. Veamos que
            $S \stackrel{\ast}{\Longrightarrow} w$:
            \begin{equation*}
                S \Longrightarrow aX \Longrightarrow au
            \end{equation*}
            donde en el último paso hemos empleado lo visto respecto a la regla de producción $X\rightarrow aX \mid bX \mid \veps$. Por tanto, $w\in \cc{L}(G)$.

            \item[$\supset)$] Sea $w\in\cc{L}(G)$. Veamos la forma de $w$:
            \begin{equation*}
                S \Longrightarrow aX \Longrightarrow au \mid u\in\{a,b\}^\ast
            \end{equation*}
            donde en el último paso hemos empleado lo visto respecto a la regla de producción $X\rightarrow aX \mid bX \mid \veps$. Por tanto, $w\in L$.
        \end{description}

        \item Describe el lenguaje generado por la gramática teniendo en cuenta que $P$ viene descrito por:
        \begin{align*}
            S &\rightarrow XaXaX \\
            X &\rightarrow aX \mid bX \mid \veps
        \end{align*}

        Sea $L=\{uavawa \mid u,v,w\in\{a,b\}^\ast\}$. Demostraremos mediante doble inclusión que $L=\cc{L}(G)$.
        \begin{description}
            \item[$\subset)$] Sea $z\in L$. Entonces, $z=uavawa$ con $u,v,w\in\{a,b\}^\ast$. Veamos que
            $S \stackrel{\ast}{\Longrightarrow} z$:
            \begin{equation*}
                S \Longrightarrow XaXaX \Longrightarrow uavawa
            \end{equation*}
            donde en el último paso hemos empleado lo visto respecto a la regla de producción $X\rightarrow aX \mid bX \mid \veps$. Por tanto, $z\in \cc{L}(G)$.

            \item[$\supset)$] Sea $z\in\cc{L}(G)$. Veamos la forma de $z$:
            \begin{equation*}
                S \Longrightarrow XaXaX \Longrightarrow uavawa \mid u,v,w\in\{a,b\}^\ast
            \end{equation*}
            donde en el último paso hemos empleado lo visto respecto a la regla de producción $X\rightarrow aX \mid bX \mid \veps$. Por tanto, $z\in L$.
        \end{description}

        \item Describe el lenguaje generado por la gramática teniendo en cuenta que $P$ viene descrito por:
        \begin{align*}
            S &\rightarrow SS \mid XaXaX \mid \veps \\
            X &\rightarrow bX \mid \veps
        \end{align*}

        Sea el lenguaje $L=\{b^i a b^j a b^k \mid i,j,k\in \bb{N}\cup \{0\}\}$. Demostraremos mediante doble inclusión que $L^\ast=\cc{L}(G)$.
        \begin{description}
            \item[$\subset)$] Sea $z\in L^\ast=\bigcup\limits_{i\in \bb{N}} L^i$.
            Sea $n$ el menor número natural tal que $z\in L^n$.
            Notando por $n_a(z)$ al número de $a$'s en $z$, tenemos que $n_a(z)=2n$.
            Entonces, $z\in L\cdot \ldots \cdot L$ ($n$ veces), por lo que existen
            $i_1,j_1,k_1,\ldots,i_n,j_n,k_n\in \bb{N}\cup \{0\}$ tales que $z=b^{i_1} a b^{j_1} a b^{k_1} \cdot \ldots \cdot b^{i_n} a b^{j_n} a b^{k_n}$. Veamos que
            $S \stackrel{\ast}{\Longrightarrow} z$:
            \begin{itemize}
                \item Para conseguir el número de $a$'s deseado, empleamos la regla de producción $S \rightarrow SS$ y reemplazamos una de las $S$ por $XaXaX$. Esto lo hacemos $n$ veces.
                \item Posteriormente, cada $X$ la sustituiremos tantas veces como sea necesario por $bX$ para conseguir el número de $b$'s deseado en cada posición, y finalizaremos con $X\rightarrow \veps$.
            \end{itemize}

            \item[$\supset)$] Sea $z\in\cc{L}(G)$, y sea $n_a(z)$ el número de $a$'s en $z$. Entonces, como el número de $a$ siempre aumenta de dos en dos, tenemos que $n_a(z)=2n$ para algún $n\in \bb{N}\cup \{0\}$.
            Veamos la forma de $z$:
            \begin{itemize}
                \item Para llegar a $z$, hemos tenido que emplear la regla de producción $S \rightarrow SS\rightarrow SXaXaX$ $n$ veces. Una vez llegados aquí, para eliminar la $S$ (ya que habremos llegado a $n_a(z)$ $a$'s), empleamos la regla de producción $S\rightarrow \veps$.
                \item Posteriormente, para cada $X$, tan solo podemos emplear la regla de producción $X\rightarrow bX \mid \veps$ para conseguir el número de $b$'s deseado en cada posición.
            \end{itemize}
            Por tanto, es directo ver que $z\in L^n\subseteq L^\ast$.
        \end{description}
    \end{enumerate}
\end{ejercicio}



\begin{ejercicio} \label{ej:1.2}
    Sea la gramática $G=\left(V,T,P,S\right)$. Determinar en cada caso el lenguaje generado por la gramática.
    \begin{enumerate}
        \item Tenga en cuenta que:
        \begin{align*}
            V &= \{S,A\}\\
            T &= \{a,b\}\\
            S &= S \\
            P &= \left\{
                \begin{array}{rcl}
                    S &\rightarrow & abAS \mid a \\
                    abA &\rightarrow & baab \\
                    A &\rightarrow & b
                \end{array}
            \right\}
        \end{align*}

        Sea $L=\{ua \mid u\in \{abb, baab\}^\ast\}$. Demostraremos mediante doble inclusión que $L=\cc{L}(G)$.
        \begin{description}
            \item[$\subset)$] Sea $w\in L$. Entonces, $w=ua$ con $u\in \{abb, baab\}^\ast$. Veamos que
            $S \stackrel{\ast}{\Longrightarrow} w$. Para ello, sabemos que $u\in \{abb, baab\}^\ast=\bigcup\limits_{i\in \bb{N}} \{abb, baab\}^i$.
            Sea $n$ el menor número natural tal que $u\in \{abb, baab\}^n$, es decir, es una concatenación de $n$ subcadenas, cada una de las cuales es o bien $abb$ o bien $baab$. Veamos que $S$ produce ambas subcadenas:
            \begin{itemize}
                \item Para producir $abb$, tenemos que $S\rightarrow abAS \rightarrow abbS$.
                \item Para producir $baab$, tenemos que $S\rightarrow abAS \rightarrow baabS$.
            \end{itemize}
            Como vemos, en cada caso podemos concatenar la subcadena necesaria, pero siempre nos quedará una $S$ al final. Usamos la regla de producción $S\rightarrow a$ para eliminarla, llegando así a $w$, por lo que $S \stackrel{\ast}{\Longrightarrow} w$ y $w\in \cc{L}(G)$.

            \item[$\supset)$] Sea $w\in\cc{L}(G)$. Veamos la forma de $w$, para lo cual hay dos opciones:
            \begin{itemize}
                \item $S\rightarrow a$: En este caso, habremos finalizado la palabra con $a$, por lo que habremos añadido la subcadena $a$ a la palabra al final.
                \item $S \rightarrow abAS$: En este caso, también hay dos opciones:
                \begin{itemize}
                    \item $S \rightarrow abAS \rightarrow baabS$: En este caso, habremos concatenado $baab$ con $S$, por lo que habremos añadido la subcadena $baab$ a la palabra.
                    \item $S \rightarrow abAS \rightarrow abbS$: En este caso, habremos concatenado $abb$ con $S$, por lo que habremos añadido la subcadena $abb$ a la palabra.
                \end{itemize}
            \end{itemize}
            Por tanto, $w$ es de la forma $ua$ con $u$ una concatenación de $abb$'s y $baab$'s, es decir, $u\in\{abb, baab\}^\ast$.
            Por tanto, $w\in L$.
        \end{description}

        \item \label{ej:1.2.b} Tenga en cuenta que:
        \begin{align*}
            V &= \{\langle \text{número} \rangle, \langle \text{dígito} \rangle\} \\
            T &= \{0,1,2,3,4,5,6,7,8,9\} \\
            S &= \langle \text{número} \rangle \\
            P &= \left\{
                \begin{array}{rcl}
                    \langle \text{número} \rangle &\rightarrow & \langle \text{número} \rangle \langle \text{dígito} \rangle \\
                    \langle \text{número} \rangle &\rightarrow & \langle \text{dígito} \rangle \\
                    \langle \text{dígito} \rangle &\rightarrow & 0 \mid 1 \mid 2 \mid 3 \mid 4 \mid 5 \mid 6 \mid 7 \mid 8 \mid 9
                \end{array}
            \right\}
        \end{align*}

        Tenemos que $\cc{L}(G)$ es el conjunto de los números naturales, permitiendo
        tantos ceros a la izquierda como se quiera. Es decir (usando la notación de potencia y concatenación vista para lenguajes):
        \begin{equation*}
            L = \{0^i n \mid i\in \bb{N}\cup\{0\},~n\in \bb{N}\cup \{0\}\}
        \end{equation*}
        Demostrémoslo mediante doble inclusión que $L=\cc{L}(G)$.
        \begin{description}
            \item[$\subset)$] Sea $w\in L$. Entonces, $w=0^i n$ con $i\in \bb{N}\cup\{0\}$ y $n\in \bb{N}\cup \{0\}$. Veamos que
            $\langle \text{número} \rangle \stackrel{\ast}{\Longrightarrow} w$:
            \begin{itemize}
                \item En primer lugar, aplicamos $|w|-1$ veces la regla de producción $\langle \text{número} \rangle \rightarrow \langle \text{número} \rangle \langle \text{dígito} \rangle$ y la regla
                que lleva de $\langle \text{dígito} \rangle$ a uno de los símbolos terminales, consiguiendo así en cada etapa reemplazar
                la última variable presente en la cadena por un dígito.
                \item Finalmente, aplicamos la regla de producción $\langle \text{número} \rangle \rightarrow \langle \text{dígito} \rangle$ para reemplazar la última variable por un dígito, que será el primero del número formado.
            \end{itemize}
            Por tanto, $\langle \text{número} \rangle \stackrel{\ast}{\Longrightarrow} w$, teniendo que $w\in \cc{L}(G)$.

            \item[$\supset)$] Sea $w\in\cc{L}(G)$. Como la única regla que
            aumenta la longitud es la regla de producción $\langle \text{número} \rangle \rightarrow \langle \text{número} \rangle \langle \text{dígito} \rangle$, tenemos que $w$ tiene la forma:
            \begin{align*}
                \langle \text{número} \rangle &\Longrightarrow \langle \text{número} \rangle \langle \text{dígito} \rangle \stackrel{|w|-1\text{\ veces}}{\Longrightarrow} \\
                &\Longrightarrow
                \langle \text{número} \rangle \langle \text{dígito} \rangle \langle \text{dígito} \rangle \stackrel{|w|-1\text{\ veces}}{\cdots} \langle \text{dígito} \rangle
                \Longrightarrow \\& \Longrightarrow
                \langle \text{dígito} \rangle \stackrel{|w|\text{\ veces}}{\cdots} \langle \text{dígito} \rangle
            \end{align*}
            Por tanto, tenemos que se trata una sucesión de $|w|$ dígitos, lo que nos lleva a que $w\in L$.
        \end{description}

        \item Tenga en cuenta que:
        \begin{align*}
            V &= \{A,S\} \\
            T &= \{a,b\} \\
            S &= S \\
            P &= \left\{
                \begin{array}{rcl}
                    S &\rightarrow & aS \mid aA \\
                    A &\rightarrow & bA \mid b
                \end{array}
            \right\}
        \end{align*}

        Sea $L=\{a^nb^m \in \{a,b\}^\ast \mid n, m \in \bb{N}\}$. Demostraremos mediante doble inclusión que $L=\cc{L}(G)$.
        \begin{description}
            \item[$\subset)$] Sea $w\in L$. Entonces, $w=a^nb^m$ con $n,m\in \bb{N}$. Veamos que
            $S \stackrel{\ast}{\Longrightarrow} w$:
            \begin{itemize}
                \item En primer lugar, aplicamos $n-1$ veces la regla de producción $S \rightarrow aS$ para obtener $a^{n-1}S$,
                \begin{equation*}
                    S \stackrel{\ast}{\Longrightarrow} a^{n-1}S
                \end{equation*}

                \item Para cambiar a la etapa de añadir $b$'s, aplicamos la regla de producción $S \rightarrow aA$, obteniendo así $a^{n}A$,
                \item Después, aplicamos $m-1$ veces la regla de producción $A \rightarrow bA$ para obtener $a^nb^{m-1}A$.
                \item Para finalizar, aplicamos la regla de producción $A \rightarrow b$ para obtener $a^nb^m$.
            \end{itemize}
            Por tanto, $S \stackrel{\ast}{\Longrightarrow} w$, teniendo que $w\in \cc{L}(G)$.

            \item[$\supset)$] Sea $w\in\cc{L}(G)$. Vemos que en la palabra siempre
            va a haber tan solo una variable (ya sea $S$ o $A$). Se empezará con la $S$, y en cierto momento se cambiará a la $A$,
            sin poder entonces volver a la $S$.
            \begin{itemize}
                \item Cuando se está en la etapa en la que hay $S$, tan solo se pueden añadir $a$'s,
                o bien cambiar a la $A$.
                \item Cuando se está en la etapa en la que hay $A$, tan solo se pueden añadir $b$'s.
            \end{itemize}
            Por tanto, tenemos que $w$ estará formada por una sucesión de
            $a$'s seguida de una sucesión de $b$'s, lo que nos lleva a que $w\in L$.
        \end{description}
    \end{enumerate}
\end{ejercicio}

\begin{ejercicio}
    Encontrar gramáticas de tipo 2 para los siguientes lenguajes sobre el alfabeto $\{a, b\}$. En cada caso determinar si los lenguajes generados son de tipo 3, estudiando si existe una gramática de tipo 3 que los genera.
    \begin{enumerate}
        \item Palabras en las que el número de $b$ no es tres.
        
        Tenemos varias opciones:
        \begin{itemize}
            \item Que no tenga $b$'s.
            \item Que tenga una $b$.
            \item Que tenga dos $b$'s.
            \item Que tenga $4$ o más $b$'s.
        \end{itemize}

        Sea la gramática $G=\left(V,T,P,S\right)$ dada por:
        \begin{align*}
            V &= \{S, A, X\} \\
            T &= \{a,b\} \\
            S &= S \\
            P &= \left\{
                \begin{array}{rcl}
                    S &\rightarrow & A \mid AbA \mid AbAbA \mid XbXbXbXbX \\
                    A &\rightarrow & aA \mid \veps \\
                    X &\rightarrow & aX \mid bX \mid \veps
                \end{array}
            \right\}
        \end{align*}

        Esta gramática no obstante es de tipo $2$. Busquemos otra que sea de tipo 3.
        Sea la gramática $G'=\left(V',T',P',S'\right)$ dada por:
        \begin{align*}
            V' &= \{S, X,Y,Z, W\} \\
            T' &= \{a,b\} \\
            S' &= S \\
            P' &= \left\{
                \begin{array}{rcl}
                    S &\rightarrow & \veps \mid aS \mid bX \\
                    X &\rightarrow & \veps \mid aX \mid bY \\
                    Y &\rightarrow & \veps \mid aY \mid bZ \\
                    Z &\rightarrow & aZ \mid bW \\
                    W &\rightarrow & \veps \mid aW \mid bW
                \end{array}
            \right\}
        \end{align*}

        Esta sí es de tipo $3$, y genera el lenguaje deseado.



        \item Palabras que tienen 2 ó 3 $b$.
        
        Sea la gramática $G=\left(V,T,P,S\right)$ dada por:
        \begin{align*}
            V &= \{S, A, B\} \\
            T &= \{a,b\} \\
            S &= S \\
            P &= \left\{
                \begin{array}{rcl}
                    S &\rightarrow & AbAbABA \\
                    A &\rightarrow & aA \mid \veps \\
                    B &\rightarrow & b \mid \veps
                \end{array}
            \right\}
        \end{align*}

        Esta gramática no obstante es de tipo $2$. Busquemos otra que sea de tipo 3.
        Sea la gramática $G'=\left(V',T',P',S'\right)$ dada por:
        \begin{align*}
            V' &= \{S, X,Y,Z,W,V,T\} \\
            T' &= \{a,b\} \\
            S' &= S \\
            P' &= \left\{
                \begin{array}{rcl}
                    S &\rightarrow & aS \mid X \\
                    X &\rightarrow & bY \\
                    Y &\rightarrow & aY \mid Z \\
                    Z &\rightarrow & bW \\
                    W &\rightarrow & aW \mid \veps \mid V \\
                    V &\rightarrow & bT \\
                    T &\rightarrow & aT \mid \veps
                \end{array}
            \right\}
        \end{align*}

        Esta gramática ya es de tipo $3$, pero contiene un número elevado de variables. Veamos si podemos reducirlo:
        Sea la gramática $G''=\left(V'',T'',P'',S''\right)$ dada por:
        \begin{align*}
            V'' &= \{S, X,Y,Z\} \\
            T'' &= \{a,b\} \\
            S'' &= S \\
            P'' &= \left\{
                \begin{array}{rcl}
                    S &\rightarrow & aS \mid bX \\
                    X &\rightarrow & aX \mid bY \\
                    Y &\rightarrow & aY \mid \veps \mid bZ \\
                    Z &\rightarrow & aZ \mid \veps
                \end{array}
            \right\}
        \end{align*}

        Notemos que, en esta gramática de tipo $3$, ya hemos conseguido el menor número de variables posibles, que representan las $4$ etapas. Como la última es opcional, está la regla $Y\rightarrow \veps$, para así no agregar la tercera $b$.

    \end{enumerate}
\end{ejercicio}

\begin{ejercicio}
    Encontrar gramáticas de tipo 2 para los siguientes lenguajes sobre el alfabeto $\{a, b\}$. En cada caso determinar si los lenguajes generados son de tipo 3, estudiando si existe una gramática de tipo 3 que los genera.
    \begin{enumerate}
        \item Palabras que no contienen la subcadena $ab$.
        
        Sea la gramática $G=\left(V,T,P,S\right)$ dada por:
        \begin{align*}
            V &= \{S, A\} \\
            T &= \{a,b\} \\
            S &= S \\
            P &= \left\{
                \begin{array}{rcl}
                    S &\rightarrow & aA \mid bS \mid \veps \\
                    A &\rightarrow & aA \mid \veps \\
                \end{array}
            \right\}
        \end{align*}

        Notemos además que esta gramática es de tipo $3$, y se tiene que:
        \begin{equation*}
            \cc{L}(G) = \{b^i a^j \mid i,j\in \bb{N}\cup \{0\}\}
        \end{equation*}


        \item Palabras que no contienen la subcadena $baa$.
        
        Sea la gramática $G=\left(V,T,P,S\right)$ dada por:
        \begin{align*}
            V &= \{S, B\} \\
            T &= \{a,b\} \\
            S &= S \\
            P &= \left\{
                \begin{array}{rcl}
                    S &\rightarrow & aS \mid bB \mid \veps \\
                    B &\rightarrow & bB \mid abB \mid a \mid \veps
                \end{array}
            \right\}
        \end{align*}
        Notemos además que esta gramática es de tipo $3$.
    \end{enumerate}
\end{ejercicio}

\begin{ejercicio}
    Encontrar una gramática libre de contexto que genere el lenguaje sobre el alfabeto $\{a, b\}$ de las palabras que tienen más $a$ que $b$ (al menos una más).

    Sea la gramática $G=\left(V,T,P,S\right)$ dada por:
    \begin{align*}
        V &= \{S, S'\} \\
        T &= \{a,b\} \\
        S &= S \\
        P &= \left\{
            \begin{array}{rcl}
                S &\rightarrow & S'aS'\\
                S' &\rightarrow & S'aS' \mid aS'bS' \mid bS'aS' \mid \veps
            \end{array}
        \right\}
    \end{align*}
\end{ejercicio}

\begin{ejercicio}
    Encontrar, si es posible, una gramática regular (o, si no es posible, una gramática libre del contexto) que genere el lenguaje $L$ supuesto que $L \subset \{a, b\}^\ast$ y verifica:
    \begin{enumerate}
        \item $u \in L$ si, y solamente si, verifica que $u$ no contiene dos símbolos $b$ consecutivos.
        
        Sea la gramática $G=\left(V,T,P,S\right)$ dada por:
        \begin{align*}
            V &= \{S\} \\
            T &= \{a,b\} \\
            S &= S \\
            P &= \left\{
                \begin{array}{rcl}
                    S &\rightarrow & aS \mid baS \mid b\mid \veps
                \end{array}
            \right\}
        \end{align*}
        \item $u \in L$ si, y solamente si, verifica que $u$ contiene dos símbolos $b$ consecutivos.
        
        Sea la gramática $G=\left(V,T,P,S\right)$ dada por:
        \begin{align*}
            V &= \{S, B, F\} \\
            T &= \{a,b\} \\
            S &= S \\
            P &= \left\{
                \begin{array}{rcl}
                    S &\rightarrow & aS \mid bB \\
                    B &\rightarrow & bF \mid aS \\
                    F &\rightarrow & aF \mid bF \mid \veps
                \end{array}
            \right\}
        \end{align*}
        Notemos que, en este caso, tenemos tres estados:
        \begin{itemize}
            \item $S$: No hemos encontrado dos $b$'s consecutivas.
            \item $B$: Hemos encontrado una $b$, y puede ser que nos encontremos la segunda $b$.
            \item $F$: Hemos encontrado dos $b$'s consecutivas; ya hay libertad.
        \end{itemize}

        Sí es cierto que usamos tres variables. Para usar solo dos variables,
        podemos hacer lo siguiente.
        Sea la gramática $G'=\left(V',T',P',S'\right)$ dada por:
        \begin{align*}
            V' &= \{S, X\} \\
            T' &= \{a,b\} \\
            S' &= S \\
            P' &= \left\{
                \begin{array}{rcl}
                    S &\rightarrow & aS \mid bS \mid bbX \\
                    X &\rightarrow & aX \mid bX \mid \veps
                \end{array}
            \right\}
        \end{align*}
    \end{enumerate}
\end{ejercicio}

\begin{ejercicio}
    Encontrar, si es posible, una gramática regular (o, si no es posible, una gramática libre del contexto) que genere el lenguaje $L$ supuesto que $L \subset \{a, b\}^\ast$ y verifica:
    \begin{enumerate}
        \item $u \in L$ si, y solamente si, verifica que contiene un número impar de símbolos $a$.
        
        Sea la gramática $G=\left(V,T,P,S\right)$ dada por:
        \begin{align*}
            V &= \{S, X\} \\
            T &= \{a,b\} \\
            S &= S \\
            P &= \left\{
                \begin{array}{rcl}
                    S &\rightarrow & aX \mid bS\\
                    X &\rightarrow & aS \mid bX \mid \veps
                \end{array}
            \right\}
        \end{align*}
        \item $u \in L$ si, y solamente si, verifica que no contiene el mismo número de símbolos $a$ que de símbolos $b$.
        
        Sea la gramática $G=\left(V,T,P,S\right)$ dada por:
        \begin{align*}
            V &= \{S, A, B, X\} \\
            T &= \{a,b\} \\
            S &= S \\
            P &= \left\{
                \begin{array}{rcl}
                    S &\rightarrow & AaA \mid BbB \\
                    A &\rightarrow & AaA\mid X \\
                    B &\rightarrow & BbB \mid X \\
                    X &\rightarrow & aXbX \mid bXaX \mid \veps
                \end{array}
            \right\}
        \end{align*}
    \end{enumerate}
\end{ejercicio}


\begin{ejercicio}
    Dado el alfabeto $A = \{a, b\}$ determinar si es posible encontrar una gramática libre de contexto que:
    \begin{enumerate}
        \item Genere las palabras de longitud impar, y mayor o igual que 3, tales que la primera letra coincida con la letra central de la palabra.
        
        Sea la gramática $G=\left(V,T,P,S\right)$ dada por:
        \begin{align*}
            V &= \{S, X, A, B, C, D\} \\
            T &= \{a,b\} \\
            S &= S \\
            P &= \left\{
                \begin{array}{rcl}
                    S &A\mid B \\
                    A &\rightarrow & aCX \\
                    C & \rightarrow & a \mid XCX \\
                    B &\rightarrow & bDX \\
                    D & \rightarrow & b \mid XDX \\
                    X &\rightarrow & a\mid b
                \end{array}
            \right\}
        \end{align*}
        \item Genere las palabras de longitud par, y mayor o igual que 2, tales que las dos letras centrales coincidan.
        
        Sea la gramática $G=\left(V,T,P,S\right)$ dada por:
        \begin{align*}
            V &= \{S, X\} \\
            T &= \{a,b\} \\
            S &= S \\
            P &= \left\{
                \begin{array}{rcl}
                    S &\rightarrow & XSX\mid C \\
                    C &\rightarrow & aa \mid bb \\
                    X &\rightarrow & a \mid b
                \end{array}
            \right\}
        \end{align*}
    \end{enumerate}
\end{ejercicio}

\begin{ejercicio}
    Sea la gramática $G=\left(V,T,P,S\right)$ dada por:
    \begin{align*}
        V &= \{S, X\} \\
        T &= \{a,b\} \\
        S &= S\\
        P &= \left\{
            \begin{array}{rcl}
                S &\rightarrow & SS \\
                S &\rightarrow & XXX \\
                X &\rightarrow & aX \mid Xa \mid b
            \end{array}
        \right\}
    \end{align*}
    Determinar si el lenguaje generado por la gramática es regular. Justificar la respuesta.\\

    Sea la siguiente gramática regular $G'=\left(V',T',P',S'\right)$ dada por:
    \begin{align*}
        V' &= \{S, X\} \\
        T' &= \{a,b\} \\
        S' &= S \\
        P' &= \left\{
            \begin{array}{rcl}
                S &\rightarrow & aS \mid bX \\
                X &\rightarrow & aX \mid bY \\
                Y &\rightarrow & aY \mid bZ \\
                Z &\rightarrow & aZ \mid bW \mid \veps \\
                W &\rightarrow & aW \mid bU \\
                U &\rightarrow & aU \mid bV \\
                V &\rightarrow & aV \mid \veps
            \end{array}
        \right\}
    \end{align*}

    Tenemos que $\cc{L}(G) = \cc{L}(G')$, y como $G'$ es una gramática regular, tenemos que $\cc{L}(G)$ es regular.
    Sí es cierto que en el tema $2$ aprendemos otras maneras de demostrarlo más sencillas, como buscar un autómata finito que lo genere.
\end{ejercicio}

\begin{ejercicio}
    Dado un lenguaje $L$ sobre un alfabeto $A$, ¿es $L^{\ast}$ siempre numerable? ¿nunca lo es? ¿o puede serlo unas veces sí y otras, no? Pon ejemplos en este último caso.\\

    $L^{\ast}$ es siempre numerable, veámos por qué. $L^{\ast}$ es un lenguaje sobre el alfabeto $A$, por lo que $L^{\ast}\subseteq A^{\ast}$ y $A^{\ast}$ es numerable (visto en teoría), luego $L^{\ast}$ también lo es.
\end{ejercicio}

\begin{ejercicio}
    Dado un lenguaje $L$ sobre un alfabeto $A$, caracterizar cuando $L^{\ast} = L$. Esto es, dar un conjunto de propiedades sobre $L$ de manera que $L$ cumpla esas propiedades si y sólo si $L^{\ast} = L$.

    \begin{equation*}
        L = L^{\ast} \Longleftrightarrow \left\{
            \begin{array}{cl}
                \veps \in L \\ \land \\ u,v \in L & \Longrightarrow uv\in L
            \end{array}
        \right.
    \end{equation*}
    Es decir, $L=L^{\ast}$ si y solo si la cadena vacía está en $L$ y además es cerrado para concatenaciones.

    \begin{proof} Demostramos mediante doble implicación.
        \begin{description}
            \item [$\Longleftarrow)$] La inclusión $L\subseteq L^{\ast}$ es obvia, por lo que solo falta demostrar la otra inclusión.\\

                Sea $v\in L^{\ast}$:
                \begin{enumerate}
                    \item Si $v = \veps \Longrightarrow v\in L$ por hipótesis.
                    \item Si $v\neq \veps$, $\exists n\in \mathbb{N}$ tal que 
                        \begin{equation*}
                            v = a_1 a_2 \ldots a_n
                        \end{equation*}
                        con $a_i \in L$ $\forall i \in \{1, \ldots, n\}$, de donde tenemos que $v\in L$, por ser cerrado para concatenaciones. Luego $L^{\ast}\subseteq L$.
                \end{enumerate}
            \item [$\Longrightarrow)$] Hemos de probar dos cosas:
                \begin{enumerate}
                    \item $\veps \in L^{\ast}=L$.
                    \item Sean $u,v\in L=L^{\ast} \Longrightarrow uv\in L^{\ast}=L$.
                \end{enumerate}
        \end{description}
    \end{proof}
\end{ejercicio}

\begin{ejercicio}
    Dados dos homomorfismos $f : A^{\ast} \rightarrow B^{\ast}$, $g : A^{\ast} \rightarrow B^{\ast}$, se dice que son iguales si $f(x) = g(x)$, $\forall x \in A^{\ast}$. ¿Existe un procedimiento algorítmico para comprobar si dos homomorfismos son iguales?\\

    Sí, basta probar que su imagen coincide sobre un conjunto finito de elementos, los de $A$:
    \begin{equation*}
        f(x) = g(x) \quad \forall x\in A^{\ast} \Longleftrightarrow f(a)=g(a) \quad \forall a\in A
    \end{equation*}
    \begin{proof}\ 
        \begin{description}
            \item [$\Longleftarrow)$] Sea $v\in A^{\ast}$, $\exists n\in \mathbb{N}$ tal que $v=a_1a_2\ldots a_n$ con $a_i \in A$ $\forall i \in \{1,\ldots, n\}$
                \begin{equation*}
                    f(v) = f(a_1)f(a_2)\ldots f(a_n) = g(a_1)g(a_2)\ldots g(a_n) = g(v)
                \end{equation*}
            \item [$\Longrightarrow)$] Sea $a\in A \Longrightarrow a\in A^{\ast}\Longrightarrow f(a)=g(a)$.
        \end{description}
    \end{proof}
\end{ejercicio}

\begin{ejercicio}
    Sea $L \subseteq A^{\ast}$ un lenguaje arbitrario. Sea $C_0 = L$ y definamos los lenguajes $S_i$ y $C_i$, para todo $i \geq 1$, por $S_i = C_{i-1}^+$ y $C_i = \ol{S_i}$. 
    \begin{enumerate}
        \item ¿Es $S_1$ siempre, nunca o a veces igual a $C_2$? Justifica la respuesta.
        \item Demostrar que $S_2 = C_3$, cualquiera que sea $L$.
        \begin{observacion}
            Demuestra que $C_2$ es cerrado para la concatenación.
        \end{observacion}
    \end{enumerate}
    % // TODO: Hacer JJ
\end{ejercicio}

\begin{ejercicio}
    Demuestra que, para todo alfabeto $A$, el conjunto de los lenguajes finitos sobre dicho alfabeto es numerable.

    Sea $A=\{a_1, a_2, \ldots, a_n\}$, con $n\in \mathbb{N}$. Definimos el siguiente conjunto:
    \begin{equation*}
        \Gamma = \{L\subseteq A^{\ast} \mid L \text{\ es finito}\}
    \end{equation*}

    Dado un símbolo $z\notin A$, definimos el conjunto $B=\{z\}\cup A$. Sea $B^{\ast}$ numerable, y buscamos una inyección de $\Gamma$ en $B^{\ast}$.
    Dado un lenguaje $L\in \Gamma$, sea $L=\{l_1, l_2, \ldots, l_m\}$, con $m\in \mathbb{N}$ y $l_i\in A^{\ast}$ $\forall i\in \{1, \ldots, m\}$. Definimos la siguiente función:
    \Func{f}{\Gamma}{B^\ast}{L}{zl_1zl_2\ldots zl_mz}

    Veamos que $f$ es inyectiva. Sean $L_1, L_2\in \Gamma$ tales que $f(L_1)=f(L_2)$. Entonces,
    \begin{equation*}
        zl_1zl_2\ldots zl_kz = zl'_1zl'_2\ldots zl'_{k'}z
    \end{equation*}
    Por ser ambas palabras iguales, tenemos que $k=k'$ y $l_i=l'_i$ $\forall i\in \{1, \ldots, k\}$, de donde $L_1=L_2$. Por tanto, $f$ es inyectiva, por lo que $\Gamma$ es inyectivo con un subconjunto de $B^{\ast}$, que es numerable. Por tanto, $\Gamma$ es numerable.
\end{ejercicio}




\subsection{Cálculo de gramáticas}

\begin{ejercicio}[Complejidad: Sencilla]
    Calcula, de forma razonada, gramáticas que generen cada uno de los siguientes lenguajes:
    \begin{enumerate}
        \item $\{ u\in \{0,1\}^\ast \mid |u|\leq 4 \}$
        
        Sea la gramática $G=\left(V,T,P,S\right)$ dada por:
        \begin{align*}
            V &= \{S, X\} \\
            T &= \{0,1\} \\
            S &= S \\
            P &= \left\{
                \begin{array}{rcl}
                    S &\rightarrow & XXXX \\
                    X &\rightarrow & 0 \mid 1 \mid \veps
                \end{array}
            \right\}
        \end{align*}

        No obstante, esta gramática es de tipo $2$. Busquemos una de tipo $3$.
        Sea la gramática $G'=\left(V',T',P',S'\right)$ dada por:
        \begin{align*}
            V' &= \{S, X, Y, Z\} \\
            T' &= \{0,1\} \\
            S' &= S \\
            P' &= \left\{
                \begin{array}{rcl}
                    S &\rightarrow & 0X \mid 1X \mid \veps \\
                    X &\rightarrow & 0Y \mid 1Y \mid \veps \\
                    Y &\rightarrow & 0Z \mid 1Z \mid \veps \\
                    Z &\rightarrow & 0 \mid 1
                \end{array}
            \right\}
        \end{align*}
        Tenemos que $\cc{L}(G) = \cc{L}(G')$, y es igual al lenguaje deseado. Tenemos por tanto que es un lenguaje regular.


        \item Palabras con 0's y 1's que no contengan dos 1's consecutivos y que empiecen por un 1 y que terminen por dos 0's.
        
        Sea la gramática $G=\left(V,T,P,S\right)$ dada por:
        \begin{align*}
            V &= \{S, X, Y\} \\
            T &= \{0,1\} \\
            S &= S \\
            P &= \left\{
                \begin{array}{rcl}
                    S &\rightarrow & 1X00 \\
                    X &\rightarrow & 0Y \mid \veps \\
                    Y &\rightarrow & 0Y \mid 1X \mid \veps \\
                \end{array}
            \right\}
        \end{align*}

        Notemos que esta gramática es de tipo 2 debido a la primera regla de producción. Busquemos una de tipo 3. 
        Sea la gramática $G'=\left(V',T',P',S'\right)$ dada por:
        \begin{align*}
            V' &= \{S, X, Y\} \\
            T' &= \{0,1\} \\
            S' &= S \\
            P' &= \left\{
                \begin{array}{rcl}
                    S &\rightarrow & 1X \\
                    X &\rightarrow & 0Y \mid F \\
                    Y &\rightarrow & 0Y \mid 1X \mid F \\
                    F &\rightarrow & 00
                \end{array}
            \right\}
        \end{align*}

        Tenemos que $\cc{L}(G) = \cc{L}(G')$, y es igual al lenguaje deseado. Tenemos por tanto que es un lenguaje regular. En esta última gramática, tenemos los siguientes estados:
        \begin{itemize}
            \item $S$: Es el estado inicial, empezamos con un $1$.
            \item $X$: Acabamos de escribir un $1$, por lo que ahora tan solo podemos escribir $0$'s.
            \item $Y$: Acabamos de escribir un $0$, por lo que ahora podemos escribir tanto $0$'s como $1$'s.
            \item $F$: Ya hemos terminado, y escribimos los dos $0$'s finales por la restricción impuesta.
        \end{itemize}
        

        \item El conjunto vacío.
        
        Sea la gramática $G=\left(V,T,P,S\right)$ dada por:
        \begin{align*}
            V &= \{S\} \\
            T &= \emptyset \\
            S &= S \\
            P &= \left\{
                \begin{array}{rcl}
                    S &\rightarrow & S
                \end{array}
            \right\}
        \end{align*}

        \item El lenguaje formado por los números naturales.
        
        Sea la gramática $G=\left(V,T,P,S\right)$ dada por:
        \begin{align*}
            V &= \{\langle \text{número no iniciado} \rangle, \langle \text{dígito no cero} \rangle, \langle \text{dígito} \rangle, \langle \text{número iniciado} \rangle\} \\
            T &= \{0,1,2,3,4,5,6,7,8,9\} \\
            S &= \langle \text{número no iniciado} \rangle \\
            P &= \left\{
                \begin{array}{rcl}
                    \langle \text{número no iniciado} \rangle &\rightarrow & \langle \text{dígito no cero} \rangle \mid \langle \text{dígito no cero} \rangle \langle \text{número iniciado} \rangle \\
                    \langle \text{número iniciado} \rangle &\rightarrow & \langle \text{dígito} \rangle \mid \langle \text{dígito} \rangle \langle \text{número iniciado} \rangle \\
                    \langle \text{dígito no cero} \rangle &\rightarrow & 1 \mid 2 \mid 3 \mid 4 \mid 5 \mid 6 \mid 7 \mid 8 \mid 9 \\
                    \langle \text{dígito} \rangle &\rightarrow & 0 \mid \langle \text{dígito no cero} \rangle
                \end{array}
            \right\}
        \end{align*}

        Notemos que esta gramática es similar a la descrita en el Ejercicio~\ref{ej:1.2}.\ref{ej:1.2.b}, pero adaptada para que los números naturales no puedan empezar por $0$.
        No obstante, esta gramática es de tipo $2$. Busquemos una de tipo $3$.
        Sea la gramática $G'=\left(V',T',P',S'\right)$ dada por:
        \begin{align*}
            V' &= \{S, X, Y, Z\} \\
            T' &= \{0,1,2,3,4,5,6,7,8,9\} \\
            S' &= S \\
            P' &= \left\{
                \begin{array}{rcl}
                    S &\rightarrow & 0 \mid 1N \mid 2N \mid 3N \mid 4N \mid 5N \mid 6N \mid 7N \mid 8N \mid 9N\\
                    N &\rightarrow & 0N\mid 1N \mid 2N \mid 3N \mid 4N \mid 5N \mid 6N \mid 7N \mid 8N \mid 9N \mid \veps
                \end{array}
            \right\}
        \end{align*}
        \item $\{ a^n \in \{a,b\}^\ast \mid n\geq 0 \} \cup \{ a^nb^n \in \{a,b\}^\ast \mid n\geq 0 \}$
        
        Sea la gramática $G=\left(V,T,P,S\right)$ dada por:
        \begin{align*}
            V &= \{S, X, Y\} \\
            T &= \{a,b\} \\
            S &= S \\
            P &= \left\{
                \begin{array}{rcl}
                    S &\rightarrow & X \mid Y \mid \veps \\
                    X &\rightarrow & aX \mid \veps \\
                    Y &\rightarrow & aYb \mid \veps
                \end{array}
            \right\}
        \end{align*}
        \item $\{ a^nb^{2n}c^m \in \{a,b,c\}^\ast \mid n,m>0 \}$
        
        Sea la gramática $G=\left(V,T,P,S\right)$ dada por:
        \begin{align*}
            V &= \{S, X, Y, Z\} \\
            T &= \{a,b,c\} \\
            S &= S \\
            P &= \left\{
                \begin{array}{rcl}
                    S &\rightarrow & aXbbcY \\
                    X &\rightarrow & aXbb \mid \veps \\
                    Y &\rightarrow & cY \mid \veps
                \end{array}
            \right\}
        \end{align*}
        \item $\{ a^nb^ma^n \in \{a,b\}^\ast \mid m,n\geq 0 \}$
        
        Sea la gramática $G=\left(V,T,P,S\right)$ dada por:
        \begin{align*}
            V &= \{S, X\} \\
            T &= \{a,b\} \\
            S &= S \\
            P &= \left\{
                \begin{array}{rcl}
                    S &\rightarrow & aSa \mid bX \mid \veps \\
                    X &\rightarrow & bX \mid \veps \\
                \end{array}
            \right\}
        \end{align*}

        \item Palabras con 0's y 1's que contengan la subcadena 00 y 11.
        
        Sea la gramática $G=\left(V,T,P,S\right)$ dada por:
        \begin{align*}
            V &= \{S, X\} \\
            T &= \{0,1\} \\
            S &= S \\
            P &= \left\{
                \begin{array}{rcl}
                    S &\rightarrow & X00X11X \mid X11X00X \\
                    X &\rightarrow & 0X \mid 1X \mid \veps
                \end{array}
            \right\}
        \end{align*}

        Notemos que esta gramática es de tipo $2$. Busquemos una de tipo $3$.
        Sea la gramática $G'=\left(V',T',P',S'\right)$ dada por:
        \begin{align*}
            V' &= \{S, X, A, B, F\} \\
            T' &= \{0,1\} \\
            S' &= S \\
            P' &= \left\{
                \begin{array}{rcl}
                    S &\rightarrow & 0S \mid 1S \mid X\\
                    X &\rightarrow & 00A \mid 11B \\
                    A &\rightarrow & 0A \mid 1A \mid 11F \\
                    B &\rightarrow & 0B \mid 1B \mid 00F \\
                    F &\rightarrow & 0F \mid 1F \mid \veps
                \end{array}
            \right\}
        \end{align*}

        Notemos que:
        \begin{itemize}
            \item $S$: No hemos encontrado ninguna subcadena.
            \item $X$: Hemos encontrado una subcadena, y ahora buscamos la otra.
            \item $A$: Hemos encontrado la subcadena $00$, y ahora buscamos la subcadena $11$.
            \item $B$: Hemos encontrado la subcadena $11$, y ahora buscamos la subcadena $00$.
            \item $F$: Hemos encontrado ambas subcadenas.
        \end{itemize}
        
        \item Palíndromos formados con las letras $a$ y $b$.
        
        Sea la gramática $G=\left(V,T,P,S\right)$ dada por:
        \begin{align*}
            V &= \{S, X, Y\} \\
            T &= \{a,b\} \\
            S &= S \\
            P &= \left\{
                \begin{array}{rcl}
                    S &\rightarrow & aSa \mid bSb \mid \veps \mid a \mid b
                \end{array}
            \right\}
        \end{align*}
        Notemos que las reglas $S\rightarrow a\mid b$ se han añadido para añadir los palíndromos de longitud impar.
    \end{enumerate}
\end{ejercicio}

\begin{ejercicio}[Complejidad: Media]
    Calcula, de forma razonada, gramáticas que generen cada uno de los siguientes lenguajes:
    \begin{enumerate}
        \item $\{uv \in \{0,1\}^\ast \mid u^{-1} \text{ es un prefijo de } v\}$
        
        Sea la gramática $G=\left(V,T,P,S\right)$ dada por:
        \begin{align*}
            V &= \{S, X, Y\} \\
            T &= \{0,1\} \\
            S &= S \\
            P &= \left\{
                \begin{array}{rcl}
                    S &\rightarrow & XY \\
                    X &\rightarrow & 0X0 \mid 1X1 \mid \veps \\
                    Y &\rightarrow & 0Y \mid 1Y \mid \veps
                \end{array}
            \right\}
        \end{align*}
        Notemos que $X$ deriva en el palíndromo, $uu^{-1}$, y $Y$ en el resto de la palabra de $v$.
        \item $\{ucv \in \{a,b,c\}^\ast \mid |u| = |v|\}$
        
        Sea la gramática $G=\left(V,T,P,S\right)$ dada por:
        \begin{align*}
            V &= \{S, X\} \\
            T &= \{a,b,c\} \\
            S &= S \\
            P &= \left\{
                \begin{array}{rcl}
                    S &\rightarrow & XSX \mid c \\
                    X &\rightarrow & a \mid b \mid c
                \end{array}
            \right\}
        \end{align*}

        \item $\{u1^n \in \{0,1\}^\ast \mid |u| = n\}$
        
        Sea la gramática $G=\left(V,T,P,S\right)$ dada por:
        \begin{align*}
            V &= \{S, X\} \\
            T &= \{0,1\} \\
            S &= S \\
            P &= \left\{
                \begin{array}{rcl}
                    S &\rightarrow & XS1 \mid \veps \\
                    X &\rightarrow & 0 \mid 1
                \end{array}
            \right\}
        \end{align*}

        \item $\{a^nb^na^{n+1} \in \{a,b\}^\ast \mid n\geq 0\}$ (observar transparencias de teoría)
        
        Sea la gramática $G=\left(V,T,P,S\right)$ dada por:
        \begin{align*}
            V &= \{S, X, Y\} \\
            T &= \{a,b\} \\
            S &= S \\
            P &= \left\{
                \begin{array}{rcl}
                    S &\rightarrow & a\mid abaa\mid aXbaa\\
                    Xb & \rightarrow & bX\\
                    Xa & \rightarrow & Ybaa\\
                    bY & \rightarrow & Yb\\
                    aY & \rightarrow aa\mid aaX
                \end{array}
            \right\}
        \end{align*}
    \end{enumerate}
\end{ejercicio}


\begin{ejercicio}[Complejidad: Difícil]
    Calcula, de forma razonada, gramáticas que generen cada uno de los siguientes lenguajes:
    \begin{enumerate}
        \item $\{a^nb^mc^k \in \{a,b,c\}^\ast \mid k = m + n\}$
        
        Sea la gramática $G=\left(V,T,P,S\right)$ dada por:
        \begin{align*}
            V &= \{S, X\} \\
            T &= \{a,b,c\} \\
            S &= S \\
            P &= \left\{
                \begin{array}{rcl}
                    S &\rightarrow & aSc \mid X \\
                    X &\rightarrow & bXc \mid \veps
                \end{array}
            \right\}
        \end{align*}
        
        \item Palabras que son múltiplos de 7 en binario.
            \begin{description}
                \item [Opción 1.] 
                    Hacer un autómata que acepte el lenguaje.
                \item [Opción 2.] 
                    Un tanto más complicada, introduce cálculos con números binarios. La idea principal es que si $x$ es un número natural, entonces:
                    \begin{equation*}
                        7x = (8-1)x = 8x - x
                    \end{equation*}
                    \begin{itemize}
                        \item Multiplicar un número en binario por 8 es añadirle tres 0s al final.
                        \item Restar un número binario menos otro es realizarle el complemento a dos al segundo, sumar los números y descartar el primer 1.
                    \end{itemize}
                    Realizar estas dos operaciones es más sencillo que multiplicar un número cualquiera en binario por 7. Procedemos por tanto, a:
                    \begin{enumerate}
                        \item Generar un número cualquiera en binario.
                        \item Multiplicarlo por 8.
                        \item Generar su complemento a 2 en binario.
                        \item Sumar ambos números.
                        \item Descarar el bit de acarreo (el más significativo).
                    \end{enumerate}
                    Para esta opción, construiremos la gramática $G=(V,T,P,S)$ dada por:
                    \begin{align*}
                        V &= \{S, N, \alpha, \beta, \delta,\gamma, Z, Z', A, D, E, E_0, E_1, E_2, E_0', E_1', \overline{E_0}, \overline{E_1}, \overline{E_2}, L_0, L_1, X\} \\
                        T &= \{0,1\} \\
                        S &= S
                    \end{align*}
                    Y $P$ es un conjunto que contiene todas las reglas de producción que se mostrarán a continuación.

                    La idea es:
                    \begin{itemize}
                        \item Generar entre $\alpha$ y $\beta$ cualquier número en binario, mientras generamos entre $\beta$ y $\gamma$ su complemento a 1 en espejo (es decir, el número invertido). Finalmente, multiplicaremos el de la izquierda por 8 y se verá reflejado en la izquierda con 1s.
                        \item Posteriormente, usaremos la variable $Z$ para sumarle 1 al complemento a 1 del número generado. 
                        \item Como no podemos modificar símbolos terminales una vez ya generados, trabajaremos todo el rato hasta el final con variables, de forma que $A$ será un 0 y $B$ un 1.
                        \item Una vez generado el número $8x$ y $x$ en complemento a dos en espejo, pasaremos a sumar ambos números usando para ello las variables $E$ y $L$. Los símbolos del número en complemento a 2 los iremos eliminando y en la izquierda controlaremos los bits del número que ya hemos usado con la variable $\delta$.
                        \item Cuando las variables $\beta$ y $\gamma$ ``se toquen'', habremos terminado de sumar y ya sólo quedará eliminar las variables delimitadoras (las letras griegas) y sustituir $A$ y $B$ por 0 y 1, respectivamente.
                    \end{itemize}
                    Comenzamos ya describiendo las reglas de producción:
                    \begin{itemize}
                        \item En primer lugar, creamos el entorno en el que trabajaremos, aceptando ``0'' como número en binario múltiplo de 7:
                            \begin{equation*}
                                S \rightarrow \alpha BNABBB\gamma\ |\ 0
                            \end{equation*}
                            Iremos usando $N$ para generar nuestro número a su izquierda y el complemento a 1 en espejo a la derecha. 
                            \begin{itemize}
                                \item Las tres $B$s ya introducidas a la derecha son para luego compensar la multiplicación por 8.
                                \item Así mismo, hemos generado ya $B$ a la izquierda y $A$ a la derecha para aceptar sólo números binarios que comiencen por 1.
                            \end{itemize}
                            Usamos ahora la variable $N$ para generar cualquier número en binario a la izquierda, con su complemento a 1 en espejo a la derecha:
                            \begin{equation*}
                                N \rightarrow ANB\ |\ BNA\ |\ AAA\gamma\beta Z
                            \end{equation*}
                            Una vez generado el número, terminaremos añadiendo 3 $A$s a la izquierda (multiplicar por 8), incluyendo los separadores $\gamma$ y $\beta$ y la variable $Z$, que se encargará de sumar 1 al número en complemento a 1 para pasarlo a complemento a 2.
                        \item Usamos ahora $Z$ para pasar el número de la derecha a complemento a 2:
                            \begin{equation*}
                                ZB \rightarrow BZ
                            \end{equation*}
                            Buscamos el primer 0, por lo que saltamos los 1s.
                            \begin{equation*}
                                ZA \rightarrow Z'B
                            \end{equation*}
                            Hemos encontrado el primer 0, lo cambiamos por 1 y volvemos con la variable $Z'$.
                            \begin{equation*}
                                BZ' \rightarrow Z'A
                            \end{equation*}
                            Volvemos a la izquierda, cambiando todos los 1s que saltamos anteriormente por 0s.
                            \begin{equation*}
                                \beta Z'\rightarrow\beta L_0
                            \end{equation*}
                            Una vez llegamos a $\beta$, tenemos el número en complemento a 1 y comenzamos con la aritmética ($L_0$ representa que no hemos cogido ningún número y que no nos llevamos nada de la suma anterior).
                        \item Comenzamos ahora con la aritmética, la parte más complicada de la gramática. Distinguimos dos casos:
                            \begin{enumerate}
                                \item No nos llevamos nada de la operación anterior ($L_0$):
                                    \begin{equation*}
                                        L_0A \rightarrow E_0
                                    \end{equation*}
                                    Cogemos un 0 de la derecha y la variable $E_0$ lo transportará a la izquierda.
                                    \begin{align*}
                                        A E_0 &\rightarrow E_0 A \\
                                        B E_0 &\rightarrow E_0 B \\
                                        \beta E_0 &\rightarrow E_0 \beta
                                    \end{align*}
                                    Nos movemos hacia la izquierda, buscando $\delta$ (que indica por dónde nos quedamos sumando).
                                    \begin{equation*}
                                        \delta E_0 \rightarrow \overline{E_0}
                                    \end{equation*}
                                    Donde la barra indica que hemos ``cogido'' $\delta$, la cual tendremos que soltar en el siguiente dígito.
                                    \begin{align*}
                                        A\overline{E_0} &\rightarrow \delta A E_0' \\
                                        B\overline{E_0} &\rightarrow \delta B E_0' 
                                    \end{align*}
                                    Como estamos sumando 0, dejamos el dígito invariante, sólo movemos $\delta$ hacia la izquierda. Usamos la variable $E_0'$ para volver, que indica que no nos llevamos nada de la suma:
                                    \begin{align*}
                                        E_0' A &\rightarrow AE_0' \\
                                        E_0' B &\rightarrow BE_0' \\
                                        E_0' \beta &\rightarrow \beta L_0
                                    \end{align*}
                                    Nos desplazamos hacia la derecha, hasta encontrar $\beta$, ya que después encontraremos el siguiente dígito con el que operar. Como no nos llevábamos nada, volvemos a $L_0$.

                                    Si ahora no nos llevamos nada y en vez de un 0 (una $A$) hay un 1 (una $B$), repetimos el proceso pero usando para ello $E_1$:
                                    \begin{align*}
                                        L_0 B &\rightarrow E_1 \\
                                        AE_1 &\rightarrow E_1 A \\
                                        BE_1 &\rightarrow E_1 B \\
                                        \beta E_1 &\rightarrow E_1 \beta \\
                                        \delta E_1 &\rightarrow \overline{E_1}
                                    \end{align*}
                                    A continuación, $\overline{E_1}$ se encontrará con el dígito con el que operar:
                                    \begin{align*}
                                        A\overline{E_1} &\rightarrow \delta B E_0' \\
                                        B\overline{E_1} &\rightarrow \delta A E_1'
                                    \end{align*}
                                    \begin{itemize}
                                        \item Si era un 0 (una $A$), lo cambiamos por un 1.
                                        \item Si era un 1 (una $B$), lo cambiamos por un 0 y nos llevamos 1 (que es lo que indica $E_1'$).
                                    \end{itemize}
                                    El comportamiento de $E_1'$ es similar a $E_0'$ pero ahora pasando a $L_1$:
                                    \begin{align*}
                                        E_1' A &\rightarrow A E_1' \\
                                        E_1' B &\rightarrow B E_1' \\
                                        E_1' \beta &\rightarrow \beta L_1 
                                    \end{align*}
                                \item Ahora, estamos en el caso en el que nos llevamos un 1 de la operación anterior ($L_1$), que hemos visto que puede suceder:
                                    \begin{equation*}
                                        L_1A \rightarrow E_1
                                    \end{equation*}
                                    Si nos encontramos un 0 llevando 1, es como si nos hubiéramos encontrado un 1 llevando 0, por lo que no hay nada nuevo que hacer. Sin embargo, si nos encontramos un 1:
                                    \begin{equation*}
                                        L_1B \rightarrow E_2
                                    \end{equation*}
                                    Tenemos que tener en mente que el siguiente dígito con el que realizar la suma lo sumaremos con 2 ($E_2$ es análogo a $E_0$ y $E_1$ pero ahora ``transportando'' un 2):
                                    \begin{align*}
                                        A E_2 &\rightarrow E_2 A \\
                                        B E_2 &\rightarrow E_2 B \\
                                        \beta E_2 &\rightarrow E_2 \beta \\
                                        \delta E_2 &\rightarrow \overline{E_2}
                                    \end{align*}
                                    A continuación, $\overline{E_2}$ se encontrará con el dígito con el que operar:
                                    \begin{align*}
                                        A\overline{E_2} &\rightarrow \delta A E_1' \\
                                        B\overline{E_2} &\rightarrow \delta B E_1'
                                    \end{align*}
                                    Similar al caso de $\overline{E_0}$, dejamos el dígito invariante pero ahora tenemos que llevarnos 1 para la siguiente opereación.
                            \end{enumerate}
                        \item A poco que se piense, como $8x$ y su complemento a 2 tienen la misma cantidad de bits, terminaremos de realizar la operación cuando nos llevemos 1 y no queden bits del número en complemento a 2, dando lugar a:
                            \begin{equation*}
                                \beta L_1 \gamma \rightarrow X
                            \end{equation*}
                            donde $X$ es una variable finalizadora, que usamos para cambiar las $A$s por 0s, las $B$s por 1s y eliminar las variables auxiliares que nos quedan (ya hemos eliminado $\beta$ y $\gamma$ directamente al crear $X$, por lo que nos quedan $\delta$ y $\alpha$):
                            \begin{align*}
                                AX &\rightarrow X0 \\
                                BX &\rightarrow X1 \\
                                \delta X &\rightarrow X \\
                                \alpha X &\rightarrow \veps
                            \end{align*}
                            Cuando lleguemos a $\alpha X$, habremos ``limpiado'' la palaba, por lo que ya podemos quitar todas las variables, generando una palabra de la gramática, que forzosamente tiene que ser un múltiplo de 7 en binario (acabamos de multiplicar cualquier número por 7). Además, como con esta gramática podemos multiplicar cualquier número por 7, esta genera todos los números que son múltiplos de 7 en binario.
                    \end{itemize}
                Mostramos finalmente un ejemplo de producción de una palabra mediante esta gramática. Trataremos de generar ``14'' en binario (la 3ª palabra que usa menos reglas de producción para ser creada, tras 0 y 7):
                \begin{align*}
                    S &\rightarrow \alpha BNABBB \gamma \rightarrow \alpha BANBABBB\gamma \rightarrow \alpha BAAAA \delta \beta Z BABBB\gamma \rightarrow \\
                      &\rightarrow \alpha BAAAA\delta \beta BZABBB\gamma \rightarrow \alpha BAAAA \delta \beta BZ' BBBB \gamma \rightarrow \\
                      &\rightarrow \alpha BAAAA \delta \beta Z' ABBBB \gamma \rightarrow \alpha BAAAA\delta \beta L_0 ABBBB\gamma \rightarrow \\
                      &\rightarrow \alpha BAAAA\delta \beta E_0 BBBB \gamma \rightarrow \alpha BAAAA\delta E_0 \beta BBBB\gamma \rightarrow \\
                      &\rightarrow \alpha BAAAA \overline{E_0}\beta BBBB\gamma \rightarrow \alpha BAAA \delta A E_0' \beta BBBB \gamma \rightarrow \\
                      &\rightarrow \alpha BAAA \delta A\beta L_0 BBBB \gamma \rightarrow \alpha BAAA \delta A\beta E_1 BBB \gamma \rightarrow \\
                      &\rightarrow \alpha BAAA \delta A E_1\beta BBB \gamma \rightarrow \alpha BAAA\delta E_1 A \beta BBB \gamma \rightarrow \\
                      &\rightarrow \alpha B AAA \overline{E_1} A \beta BBB \gamma \rightarrow \alpha BAA \delta B E_0' A \beta BBB \gamma \rightarrow \\ 
                      &\rightarrow \alpha BAA \delta BA E_0' \beta BBB \gamma \rightarrow \alpha BAA \delta BA \beta L_0 BBB \gamma \rightarrow \\
                      &\rightarrow \alpha BAA \delta BA \beta E_1 BB \gamma \rightarrow \alpha BAA \delta BA E_1 \beta BB \gamma \rightarrow \\
                      &\rightarrow \alpha BAA \delta BE_1 A \beta BB \gamma \rightarrow \alpha BAA \delta E_1 BA \beta BB \gamma \rightarrow \\
                      &\rightarrow \alpha BAA \overline{E_1} BA \beta BB \gamma \rightarrow \alpha BA \delta B E_0' BA \beta BB \gamma \rightarrow \\
                      &\rightarrow \alpha BA \delta BBE_0' A \beta BB \gamma \rightarrow \alpha BA \delta BBA E_0' \beta BB \gamma \rightarrow \\
                      &\rightarrow \alpha BA \delta BBA \beta L_0 BB \gamma \rightarrow \alpha BA \delta BBA \beta E_1 B \gamma \rightarrow \alpha BA \delta BBA E_1 \beta B\gamma \rightarrow \\
                      &\rightarrow \alpha BA \delta BBE_1 A \beta B \gamma \rightarrow \alpha BA \delta B E_1 BA \beta B \gamma \rightarrow \alpha BA \delta E_1 BBA \beta B \gamma \rightarrow \\
                      &\rightarrow \alpha BA \overline{E_1} BBA \beta B \gamma \rightarrow \alpha B \delta B E_0' BBA \beta B \gamma \rightarrow \alpha B\delta BB E_0' BA \beta B \gamma \rightarrow \\
                      &\rightarrow \alpha B\delta BBB E_0' A \beta B \gamma \rightarrow \alpha B\delta BBBA E_0' \beta B \gamma \rightarrow \alpha B \delta BBBA \beta L_0 B \gamma \rightarrow \\
                      &\rightarrow \alpha B \delta BBBA \beta E_1 \gamma \rightarrow \alpha B \delta BBBA E_1 \beta \gamma \rightarrow \alpha B \delta BBB E_1 A \beta \gamma \rightarrow \\
                      &\rightarrow \alpha B \delta BB E_1 BA \beta \gamma \rightarrow \alpha B \delta B E_1 BBA \beta \gamma \rightarrow \alpha B \delta E_1 BBBA \delta \gamma \rightarrow \\
                      &\rightarrow \alpha B \overline{E_1} BBBA \beta \gamma \rightarrow \alpha \delta A E_1' BBBA \beta \gamma \rightarrow \alpha \delta AB E_1' BBA \beta \gamma \rightarrow \\
                      &\rightarrow \alpha \delta ABB E_1' BA \beta \gamma \rightarrow \alpha \delta ABBB E_1' A \beta \gamma \rightarrow \alpha \delta ABBBA E_1' \beta \gamma \rightarrow \\
                      &\rightarrow \alpha \delta ABBBA \beta L_1 \gamma \rightarrow \alpha \delta ABBBA X \rightarrow \alpha \delta ABBBX0 \rightarrow \alpha \delta ABBX10 \rightarrow \\
                      &\rightarrow \alpha \delta ABX110 \rightarrow \alpha \delta AX1110 \rightarrow \alpha \delta X01110 \rightarrow \alpha X 01110 \rightarrow 01110
                \end{align*}
            \end{description}
    \end{enumerate}
\end{ejercicio}


\begin{ejercicio}[Complejidad: Extrema (no son libres de contexto)]
    Calcula, de forma razonada, gramáticas que generen cada uno de los siguientes lenguajes:
    \begin{enumerate}
        \item $\{ww \mid w \in \{0,1\}^\ast\}$\\
            Para este lenguaje, hemos construido la gramática $G=(V,T,P,S)$ dada por:
            \begin{align*}
                V &= \{S, \alpha, \beta, \gamma, X, E, E_1, E_0, E', B\} \\
                T &= \{0,1\} \\
                S &= S
            \end{align*}
            $P$ que contiene las reglas de producción que se mostrarán a continuación. 

            La idea principal en la gramática es generar entre las variables $\alpha$ y $\beta$ cualqueir palabra del lenguaje ${\{0,1\}}^{\ast}$. Posteriormente, iremos copiando dicha palabra a la derecha de $\beta$ usando para ello las variables $E$ y $\gamma$, de forma que con $\gamma$ controlaremos la parte de la palabra de la izquierda que ya hayamos copiado a la derecha de $\beta$.

            Finalmente, usaremos $B$ para eliminar cualquier rastro de las variable auxiliares. De esta forma, las reglas de $P$ son:
            \begin{itemize}
                \item Para generar cualquier palabra entre $\alpha$ y $\beta$:
                    \begin{align*}
                        S &\rightarrow \alpha X \beta \\
                        X &\rightarrow 0X\ |\ 1X\ |\ E\gamma
                    \end{align*}
                \item Para coger un 1 y copiarlo a la derecha:

                    Hemos de estar al final de la parte de la palabra no copiada (luego ha de ser $xE\gamma$ siendo $x$ 0 o 1). Posteriormente, avanzamos $\gamma$ a la izquierda para indicar que dicho 1 ya está copiado y cambiamos a la variable que transporta el 1 a la derecha:
                    \begin{equation*}
                        1E\gamma \rightarrow \gamma 1 E_1 
                    \end{equation*}
                    Posteriormente, movemos dicha variable a la derecha:
                    \begin{align*}
                        E_1 1 &\rightarrow 1E_1 \\
                        E_1 0 &\rightarrow 0E_1
                    \end{align*}
                    Cuando lleguemos al final de la palabra de la izquierda, soltamos el 1 al inicio de la palabra de la derecha:
                    \begin{equation*}
                        E_1 \beta \rightarrow E' \beta 1
                    \end{equation*}
                \item Para coger un 0 y copiarlo a la derecha, es una situación análoga pero usamos otra variable:
                    \begin{align*}
                        0E\gamma &\rightarrow \gamma 0 E_0 \\
                        \\
                        E_0 1 &\rightarrow 1E_0 \\
                        E_0 0 &\rightarrow 0E_0 \\
                        \\
                        E_0 \beta &\rightarrow E'\beta 0
                    \end{align*}
                \item Ahora, explicamos $E'$, cuya única funcionalidad es volver al final de la parte no copiada de la palabra de la izquierda:
                    \begin{align*}
                        1E' &\rightarrow E'1 \\
                        0E' &\rightarrow E'0 \\
                        \gamma E' &\rightarrow E\gamma
                    \end{align*}
                \item La copia de la palabra terminará cuando se de $\alpha E\gamma$ (ya que estará toda la palabra copiada a la derecha). En dicho caso, eliminamos todas las variables auxiliares restantes:
                    \begin{align*}
                        \alpha E\gamma &\rightarrow B \\
                        B1 &\rightarrow 1B \\
                        B0 &\rightarrow 0B \\
                        B\beta &\rightarrow \veps 
                    \end{align*}
            \end{itemize}
            Puede demostrarse que el lenguaje generado por esta gramática es el solicitado. Por la complejidad de la gramática, nos limitamos a mostrar un ejemplo para ver de forma intuitiva el buen funcionamiento de la misma.

            Trataremos de generar la cadena: $10111011$ (es decir, ${(1011)}^{2}$):
            \begin{align*}
                S &\rightarrow \alpha X \beta \rightarrow \alpha 1X\beta \rightarrow \alpha 10X\beta \rightarrow \alpha101X \beta \rightarrow \alpha1011X\beta \rightarrow \alpha1011E\gamma\beta \rightarrow \\
                  &\rightarrow \alpha101\gamma1E_1\beta \rightarrow \alpha101\gamma1E'\beta1 \rightarrow \alpha101\gamma E'1\beta1 \rightarrow \alpha101E\gamma1\beta1 \rightarrow \\
                  &\rightarrow \alpha10\gamma1E_11\beta1 \rightarrow \alpha 10\gamma11E_1\beta1 \rightarrow \alpha 10\gamma11E'\beta11 \rightarrow \alpha10\gamma1E'1\beta11 \rightarrow \\
                  &\rightarrow \alpha10\gamma E'11\beta11 \rightarrow \alpha10E\gamma11\beta11 \rightarrow \alpha 1\gamma0E_011\beta11 \rightarrow \alpha1\gamma01E_01\beta11 \rightarrow\\
                  &\rightarrow\alpha1\gamma011E_0\beta11 \rightarrow \alpha1\gamma011E'\beta011 \rightarrow \alpha1\gamma01E'1\beta011 \rightarrow \alpha1\gamma0E'11\beta011 \rightarrow \\
                  &\rightarrow \alpha1\gamma E'011\beta011 \rightarrow \alpha1E\gamma011\beta011 \rightarrow \alpha\gamma1E_1011\beta011 \rightarrow \alpha\gamma10E_111\beta011 \rightarrow \\
                  &\rightarrow \alpha\gamma101E_11\beta011 \rightarrow \alpha\gamma1011E_1\beta011 \rightarrow \alpha\gamma1011E'\beta1011 \rightarrow \alpha\gamma101E'1\beta1011 \rightarrow \\
                  &\rightarrow \alpha\gamma10E'11\beta1011 \rightarrow \alpha\gamma1E'011\beta1011 \rightarrow \alpha\gamma E'1011\beta1011 \rightarrow \alpha E\gamma1011\beta1011 \rightarrow \\
                  &\rightarrow B1011\beta1011 \rightarrow 1B011\beta1011 \rightarrow10B11\beta1011 \rightarrow 101B1\beta1011\rightarrow \\
                  &\rightarrow 1011B\beta1011 \rightarrow10111011
            \end{align*}
        \item $\{a^{n^2} \in \{a\}^{\ast} \mid n\geq 0\}$\\
            La idea que hemos tenido para hacer una gramática que acepte el lenguaje es la siguiente. Si representamos las 5 primeras palabras del lenguaje (ordenándolas por su longitud):
            \begin{gather*}
                \veps \\
                a \\
                aaaa \\
                aaaaaaaaa \\
                aaaaaaaaaaaaaaaa
            \end{gather*}
            Notemos que podemos ordenar las letras de la siguiente forma (olvidándonos de $\veps$, que no será relevante):
            \begin{gather*}
                a \\
                aa\ aa \\
                aaa\ aaa\ aaa\\
                aaaa\ aaaa\ aaaa\ aaaa 
            \end{gather*}
            De forma que tenemos 1 grupo de 1 ``a'', dos grupos de 2 ``a'', 3 grupos de 3 ``a'', \ldots Notemos que dados $n$ grupos de $n$ ``a'', será sencillo construir $n+1$ grupos de $n+1$ ``a'', ya que nos bastará con añadir una ``a'' a cada grupo y con duplicar el último grupo de ``a''.

            Hemos construido una gramática $G = (V, T, S, P)$ que simula este comportamiento inductivo del lenguaje, con lo que el lenguaje generado por la misma es el solicitado. Tenemos:
            \begin{align*}
                V &= \{\alpha, \beta, \delta, \gamma, \sigma, X, A, E, E_{\sigma}, \overline{E}, E', I, R, L, Z \} \\
                T &= \{0,1\} \\
                S &= S
            \end{align*}
            Donde $P$ es el conjunto de reglas de producción que contiene todas las reglas que explicaremos a continuación.

            La idea es que si queremos generar la palabra $a^{n^2}$, que generemos $n-1$ $A$s entre $\alpha$ y $\beta$. Tendremos ya creada una letra $a$ y lo que haremos será que por cada $A$ que hayamos generado, repitamos el proceso inductivo descrito anteriormente. Además, separaremos los ``grupos'' de ``a'' con variables $I$.
            Finalmente, para duplicar un grupo de ``a'', usaremos las variables $\delta$ y $\sigma$.\\

                Empezamos generando nuestro entorno en el que trabajaremos (o la palabra vacía):
                \begin{equation*}
                    S \rightarrow \alpha X \beta Ia\delta \gamma \ |\ \veps
                \end{equation*}
                A continuación, usamos $X$ para generar las $A$s:
                \begin{equation*}
                    X \rightarrow AX\ |\ E
                \end{equation*}
                Una vez terminadas de leer las $A$s, generaremos $E$, que se encargará de ir eliminando una $A$, de realizar el proceso inductivo y de volver al estado inicial, hasta terminar con todas las $A$s generadas.

             Ahora, hacemos que $E$ coja una $A$, con lo que le dejamos salir de la región comprendida por $\alpha$ y $\beta$:
                \begin{equation*}
                    AE\beta \rightarrow \beta E
                \end{equation*}
                Ahora desplazamos la variable $E$ a la derecha, haciendo que cada vez que entre en un grupo de ``a'' (cuando pase una variable $I$) añada una nueva:
                \begin{align*}
                    Ea &\rightarrow aE \\
                    EI &\rightarrow IaE
                \end{align*}
                Cuando la variable $E$ se encuentre con $\delta$, habremos terminado de incrementar las ``a'', con lo que tendremos que duplicar ahora el último grupo de ``a''. Para ello, prepararemos un entorno, de forma que entre $I$ y $\sigma$ vayamos generando el nuevo grupo de ``a'', entre $\sigma$ y $\delta$ se encuentre las ``a'' por copiar; y que entre $\delta$ y $\gamma$ se encuentren las ``a'' que ya hayan sido duplicadas.

                De esta forma, cuando $E$ se encuentre con $\delta$, pasaremos a una variable que busque $I$ para colocar delante suya $\sigma$:
                \begin{align*}
                    E\delta &\rightarrow E_{\sigma} \delta \\
                    aE_{\sigma} &\rightarrow E_{\sigma} a \\
                    IE_{\sigma} &\rightarrow I\sigma R
                \end{align*}
                Ahora, usaremos $R$ para movernos a la derecha tras copiar una letra y $L$ para movernos a la izquierda con el fin de pegar una letra.
                \begin{equation*}
                    Ra \rightarrow aR 
                \end{equation*}
                Nos movemos a la derecha
                \begin{equation*}
                    aR\delta \rightarrow L\delta a
                \end{equation*}
                Cuando lleguemos a $\delta$, guardamos una ``a'' más como copiada.
                \begin{equation*}
                    aL \rightarrow La
                \end{equation*}
                Nos moveremos hacia la izquierda buscando $\sigma$ para crear una $a$:
                \begin{equation*}
                    \sigma L \rightarrow a \sigma R
                \end{equation*}
                Pegaremos una $a$ tras $\sigma$ y repetiremos el proceso.

                El proceso terminará cuando no haya más ``a'' entre $\sigma$ y $\delta$:
                \begin{equation*}
                    \sigma R\delta \rightarrow I\overline{E}
                \end{equation*}
                Cuando hayamos terminado, colocamos una $I$ para hacer efectivo el nuevo grupo. Finalmente, debemos colocar nuevamente $\delta$ a la izquierda de $\gamma$ para la siguiente vez que copiemos. Usamos para ello $\overline{E}$:
                \begin{equation*}
                    \overline{E}a \rightarrow a\overline{E}
                \end{equation*}
                Nos movemos a la izquierda buscando $\gamma$ y cuando la encontremos, colocamos $\delta$:
                \begin{equation*}
                    \overline{E}\gamma \rightarrow E' \delta \gamma
                \end{equation*}
                Usaremos finalmente $E'$ para desplazarnos a la izquierda, tras $\beta$, donde volveremos a la variable $E$, que reiniciará el proceso descrito para realizarlo nuevamente:
                \begin{align*}
                    aE' &\rightarrow E' a \\
                    IE' &\rightarrow E' I \\
                    \beta E' &\rightarrow E\beta
                \end{align*}
                Este proceso terminará cuando no queden $A$s por copiar. En dicho caso, pasaremos a una variable $Z$ que eliminará todas las variables auxiliares:
                \begin{equation*}
                    \alpha E \rightarrow Z
                \end{equation*}
                De esta forma, $Z$ se mueve a la derecha, eliminando todas las variables y pasando a través de las letras:
                \begin{align*}
                    Z\beta &\rightarrow Z \\
                    ZI &\rightarrow Z \\
                    Za &\rightarrow aZ \\
                    Z\delta &\rightarrow Z \\
                    Z\gamma &\rightarrow \veps
                \end{align*}
            Como ejemplo y para comprobar que el lenguaje generado por dicha gramática funciona es el deseado mostramos el siguiente ejemplo, en el que generamos $a^9$:
            \begin{align*}
                S &\rightarrow \alpha X \beta I a \delta\gamma \rightarrow \alpha AX\beta Ia\delta\gamma \rightarrow \alpha AAX\beta Ia\delta\gamma \rightarrow \alpha AAE \beta Ia \delta\gamma \rightarrow \alpha A \beta E Ia\delta\gamma \rightarrow \\
                  &\rightarrow \alpha A\beta IaEa\delta\gamma \rightarrow \alpha A\beta IaaE\delta\gamma \rightarrow \alpha A\beta IaaE_{\sigma}\delta\gamma \rightarrow \alpha A\beta IaE_{\sigma}a\delta\gamma \rightarrow \\ 
                  &\rightarrow\alpha A\beta IE_{\sigma}aa\delta\gamma \rightarrow \alpha A\beta I\sigma Raa\delta\gamma \rightarrow \alpha A\beta I \sigma aRa \delta\gamma \rightarrow \alpha A \beta I \sigma aaR \delta\gamma\rightarrow \\
                  &\rightarrow \alpha A\beta I\sigma aL\delta a\gamma \rightarrow \alpha A \beta I\sigma L a\delta a\gamma \rightarrow \alpha A \beta I a\sigma R a \delta a\gamma  \rightarrow \alpha A \beta Ia\sigma aR\delta a\gamma \rightarrow \\
                  &\rightarrow \alpha A \beta I a\sigma L\delta aa \gamma \rightarrow \alpha A \beta Iaa \sigma R \delta aa \gamma \rightarrow \alpha A\beta IaaI\overline{E}aa\gamma \rightarrow \alpha A \beta IaaIa\overline{E}a\gamma \rightarrow \\
                  &\rightarrow \alpha A\beta IaaIaa\overline{E} \gamma \rightarrow \alpha A \beta IaaIaaE'\delta \gamma \rightarrow \alpha A\beta IaaIaE'a\delta\gamma \rightarrow \alpha A\beta IaaIE'aa\delta\gamma \rightarrow \\
                  &\rightarrow \alpha A\beta IaaE'Iaa\delta\gamma \rightarrow \alpha A \beta IaE'aIaa\delta\gamma \rightarrow \alpha A\beta IE'aaIaa\delta\gamma \rightarrow \\
                  &\rightarrow \alpha A \beta E'IaaIaa\delta\gamma \rightarrow \alpha AE\beta IaaIaa\delta\gamma \rightarrow \ldots
            \end{align*}
            % // TODO: Seguir
        \item $\{a^p \in \{a\}^{\ast} \mid p \text{ es primo}\}$
        \item $\{a^nb^m \in \{a,b\}^{\ast} \mid n\leq m^2\}$
    \end{enumerate}

\end{ejercicio}


    \newpage
\section{Autómatas Finitos}

\begin{comment}
\begin{tikzpicture}
    \node[state, initial] (q1) {$q_1$};
    \node[state, accepting, right of=q1] (q2) {$q_2$};
    \node[state, right of=q2] (q3) {$q_3$};
    \draw (q1) edge[loop above] node{0} (q1)
    (q1) edge[above] node{1} (q2)
    (q2) edge[loop above] node{1} (q2)
    (q2) edge[bend left, above] node{0} (q3)
    (q3) edge[bend left, below] node{0, 1} (q2);
\end{tikzpicture}
\end{comment}

\begin{ejercicio} \label{ej:1.2.1}
    Considera el siguiente Autómata Finito Determinista (AFD) dado por $M = (Q, A, \delta, q_0, F)$, donde:
    \begin{itemize}
        \item $Q = \{q_0, q_1, q_2\}$
        \item $A = \{0, 1\}$
        \item La función de transición viene dada por:
        \begin{align*}
            \delta(q_0, 0) &= q_1, & \delta(q_0, 1) &= q_0 \\
            \delta(q_1, 0) &= q_2, & \delta(q_1, 1) &= q_0 \\
            \delta(q_2, 0) &= q_2, & \delta(q_2, 1) &= q_2
        \end{align*}
        \item $F = \{q_2\}$
    \end{itemize}
    Describe informalmente el lenguaje aceptado.\\

    Su representación gráfica está en la Figura~\ref{fig:ej:1.2.1}.
    \begin{figure}
        \centering
        \begin{tikzpicture}
            \node[state, initial] (q0) {$q_0$};
            \node[state, right of=q0] (q1) {$q_1$};
            \node[state, accepting, right of=q1] (q2) {$q_2$};

            \draw   (q0) edge[loop above] node{1} (q0)
                    (q0) edge[above, bend left] node{0} (q1)
                    (q1) edge[below, bend left] node{1} (q0)
                    (q1) edge[above] node{0} (q2)
                    (q2) edge[loop above] node{0, 1} (q2);
        \end{tikzpicture}
        \caption{Autómata Finito Determinista del Ejercicio~\ref{ej:1.2.1}.}
        \label{fig:ej:1.2.1}
    \end{figure}

    Tenemos que el lenguaje aceptado por el autómata es el conjunto de todas las palabras que contienen
    la cadena $00$ como subcadena. Es decir,
    \begin{align*}
        L = \{u_100u_2 \in \{0, 1\}^* \mid u_1, u_2 \in \{0, 1\}^*\}.
    \end{align*}
\end{ejercicio}

\begin{ejercicio} \label{ej:1.2.2}
    Dado el AFD de la Figura~\ref{fig:ej:1.2.2}, describir el lenguaje aceptado por dicho autómata.
    \begin{figure}
        \centering
        \begin{tikzpicture}
            \node[state, initial] (q0) {$q_0$};
            \node[state, right of=q0] (q1) {$q_1$};
            \node[state, accepting, right of=q1] (q2) {$q_2$};

            \draw   (q0) edge[loop above] node{$b$} (q0)
                    (q0) edge[above] node{$a$} (q1)
                    (q1) edge[loop above] node{$b$} (q1)
                    (q1) edge[bend left, above] node{$a$} (q2)
                    (q2) edge[bend left, below] node{$a$} (q1)
                    (q2) edge[loop above] node{$b$} (q2);
        \end{tikzpicture}
        \caption{Autómata Finito Determinista del Ejercicio~\ref{ej:1.2.2}.}
        \label{fig:ej:1.2.2}
    \end{figure}

    El lenguaje aceptado por el autómata es el conjunto de todas las palabras que contienen
    un número par de $a$'s. Es decir,
    \begin{align*}
        L = \{u \in \{a, b\}^* \mid n_a(u) \text{\ es par},~n_a(u)>0\},
    \end{align*}
\end{ejercicio}

\begin{ejercicio} \label{ej:1.2.3}
    Dibujar AFDs que acepten los siguientes lenguajes con alfabeto $\{0, 1\}$:
    \begin{enumerate}
        \item El lenguaje vacío,
        
        \begin{figure}[H]
            \centering
            \begin{tikzpicture}
                \node[state, initial] (q0) {$q_0$};

                \draw   (q0) edge[loop above] node{0, 1} (q0);
            \end{tikzpicture}
        \end{figure}
        \item El lenguaje formado por la palabra vacía, es decir, $\{\varepsilon\}$,
        
        \begin{figure}[H]
            \centering
            \begin{tikzpicture}
                \node[state, initial, accepting] (q0) {$q_0$};
                \node[state, right of=q0, error] (E) {$E$};

                \draw   (q0) edge[above] node{0, 1} (E);
                \draw   (E) edge[loop above] node{0, 1} (E);
            \end{tikzpicture}
        \end{figure}
        \item El lenguaje formado por la palabra $01$, es decir, $\{01\}$,
        
        \begin{figure}[H]
            \centering
            \begin{tikzpicture}
                \node[state, initial] (q0) {$q_0$};
                \node[state, right of=q0] (q1) {$q_1$};
                \node[state, accepting, right of=q1] (q2) {$q_2$};
                \node[state, below of=q1, error] (E) {$E$};

                \draw   (q0) edge[above] node{0} (q1)
                        (q1) edge[above] node{1} (q2);
                \draw   (q0) edge[below] node{1} (E)
                        (q1) edge[left] node{0} (E)
                        (q2) edge[below right] node{0, 1} (E)
                        (E) edge[loop below] node{0, 1} (E);
            \end{tikzpicture}
        \end{figure}
        \item El lenguaje $\{11, 00\}$,
        
        \begin{figure}[H]
            \centering
            \begin{tikzpicture}
                \node[state, initial] (q0) {$q_0$};
                \node[state, above right of=q0] (q1) {$q_1$};
                \node[state, below right of=q0] (q2) {$q_2$};
                \node[state, accepting, right of=q0] (q3) {$q_3$};
                \node[state, right of=q3, error] (E) {$E$};

                \draw   (q0) edge[above] node{1} (q1)
                        (q0) edge[below] node{0} (q2)
                        (q1) edge[above] node{0} (E)
                        (q1) edge[left] node{1} (q3)
                        (q2) edge[left] node{0} (q3)
                        (q2) edge[below] node{1} (E);
                \draw   (E) edge[loop above] node{0, 1} (E);
                \draw   (q3) edge[above left] node{0, 1} (E);
            \end{tikzpicture}
        \end{figure}
        \item El lenguaje $\{(01)^i \mid i \geq 0\}$,
        
        \begin{figure}[H]
            \centering
            \begin{tikzpicture}
                \node[state, initial, accepting] (q0) {$q_0$};
                \node[state, below right of=q0, error] (E) {$E$};
                \node[state, above right of=E] (q1) {$q_1$};


                \draw   (q0) edge[below] node{0} (q1)
                        (q0) edge[below] node{1} (q2)
                        (q1) edge[left] node{0} (q2)
                        (q1) edge[above, bend right] node{1} (q0)
                        (q2) edge[loop below] node{0, 1} (q2);
            \end{tikzpicture}
        \end{figure}
        \item El lenguaje formado por las cadenas con $0$'s y $1$'s donde el número de unos es divisible por $3$.
        
        \begin{figure}[H]
            \centering
            \begin{tikzpicture}
                \node[state, initial, accepting] (q0) {$q_0$};
                \node[state, right of=q0] (q1) {$q_1$};
                \node[state, right of=q1] (q2) {$q_2$};

                \draw   (q0) edge[loop above] node{0} (q0)
                        (q0) edge[above] node{1} (q1)
                        (q1) edge[loop above] node{0} (q1)
                        (q1) edge[above] node{1} (q2)
                        (q2) edge[loop above] node{0} (q2)
                        (q2) edge[below, bend left] node{1} (q0);
            \end{tikzpicture}
        \end{figure}
    \end{enumerate}
\end{ejercicio}

\begin{ejercicio} \label{ej:1.2.4}
    Obtener a partir de la gramática regular $G = (\{S, B\}, \{1, 0\}, P, S)$, con
    \begin{align*}
        P = \left\{
            \begin{aligned}
                S &\to 110B \\
                B &\to 0B \mid 1B \mid \varepsilon
            \end{aligned}
        \right.
    \end{align*}
    un AFND que reconozca el lenguaje generado por esa gramática.

    El autómata obtenido es el de la Figura~\ref{fig:ej:1.2.4}.
    \begin{figure}[H]
        \centering
        \begin{tikzpicture}
            \node[state, initial] (q0) {$q_0$};
            \node[state, right of=q0] (q1) {$q_1$};
            \node[state, right of=q1] (q2) {$q_2$};
            \node[state, accepting, right of=q2] (q3) {$q_3$};

            \draw   (q0) edge[above] node{1} (q1)
                    (q1) edge[above] node{1} (q2)
                    (q2) edge[above] node{0} (q3)
                    (q3) edge[loop above] node{0, 1} (q3);
        \end{tikzpicture}
        \caption{Autómata Finito No Determinista del Ejercicio~\ref{ej:1.2.4}.}
        \label{fig:ej:1.2.4}
    \end{figure}
\end{ejercicio}

\begin{ejercicio} \label{ej:1.2.5}
    Dada la gramática regular $G = (\{S\}, \{1, 0\}, P, S)$, con
    \begin{align*}
        P = \{S &\to S10, S \to 0\},
    \end{align*}
    obtener un AFD que reconozca el lenguaje generado por esa gramática.\\

    El lenguaje es:
    \begin{equation*}
        L = \{0(10)^n \mid n\in \bb{N}\cup \{0\}\}.
    \end{equation*}

    El autómata obtenido es el de la Figura~\ref{fig:ej:1.2.5}.
    \begin{figure}[H]
        \centering
        \begin{tikzpicture}
            \node[state, initial] (q0) {$q_0$};
            \node[state, below right of=q0, error] (E) {$E$};
            \node[state, accepting, above right of=E] (q1) {$q_1$};

            \draw   (q0) edge[above, bend left] node{0} (q1)
                    (q0) edge[below] node{1} (E)
                    (q1) edge[above, bend left] node{1} (q0)
                    (q1) edge[below] node{0} (E)
                    (E) edge[loop below] node{0, 1} (E);
        \end{tikzpicture}
        \caption{Autómata Finito Determinista del Ejercicio~\ref{ej:1.2.5}.}
        \label{fig:ej:1.2.5}
    \end{figure}
\end{ejercicio}

\begin{ejercicio} \label{ej:1.2.6}
    Construir un AFND o AFD (dependiendo del caso) que acepte las cadenas $u \in \{0, 1\}^*$ que:
    \begin{enumerate}
        \item AFND\@. Contengan la subcadena $010$.~\label{ej:1.2.6.1}
        
        El autómata obtenido es el de la Figura~\ref{fig:ej:1.2.6.1}.
        \begin{figure}[H]
            \centering
            \begin{tikzpicture}
                \node[state, initial] (q0) {$q_0$};
                \node[state, right of=q0] (q1) {$q_1$};
                \node[state, right of=q1] (q2) {$q_2$};
                \node[state, accepting, right of=q2] (q3) {$q_3$};

                \draw   (q0) edge[above] node{0} (q1)
                        (q1) edge[above] node{1} (q2)
                        (q2) edge[above] node{0} (q3)
                        (q3) edge[loop above] node{0, 1} (q3)
                        (q0) edge[loop above] node{0, 1} (q0);
            \end{tikzpicture}
            \caption{Autómata Finito No Determinista del Ejercicio~\ref{ej:1.2.6} apartado~\ref{ej:1.2.6.1}.}
            \label{fig:ej:1.2.6.1}
        \end{figure}
        \item AFND\@. Contengan la subcadena $110$. \label{ej:1.2.6.2}
        
        El autómata obtenido es el de la Figura~\ref{fig:ej:1.2.6.2}.
        \begin{figure}[H]
            \centering
            \begin{tikzpicture}
                \node[state, initial] (q0) {$q_0$};
                \node[state, right of=q0] (q1) {$q_1$};
                \node[state, right of=q1] (q2) {$q_2$};
                \node[state, accepting, right of=q2] (q3) {$q_3$};

                \draw   (q0) edge[above] node{1} (q1)
                        (q1) edge[above] node{1} (q2)
                        (q2) edge[above] node{0} (q3)
                        (q3) edge[loop above] node{0, 1} (q3)
                        (q0) edge[loop above] node{0, 1} (q0);
            \end{tikzpicture}
            \caption{Autómata Finito No Determinista del Ejercicio~\ref{ej:1.2.6} apartado~\ref{ej:1.2.6.2}.}
            \label{fig:ej:1.2.6.2}
        \end{figure}
        \item AFD\@. Contengan simultáneamente las subcadenas $010$ y $110$. \label{ej:1.2.6.3}
        
        El estado $q_0$ representa que no se ha empezado ninguna de las subcadenas, y el estado $q_F$ representa que se han encontrado ambas cadenas.
        Hay dos opciones:
        \begin{description}
            \item[Opción 1] Primero se lee $010$ y luego $110$. Son los siguientes estados:
            \begin{itemize}
                \item $q_0$: Estado inicial, no ha empezado la subcadena $010$.
                \item $q_1$: Se ha leído el $0$ de la subcadena $010$.
                \item $q_2$: Se ha leído la subcadena $01$ de la subcadena $010$.
                \item $q_3$: Se ha leído la subcadena $010$. No ha empezado la subcadena $110$.
                \item $q_4$: Se ha leído el $1$ de la subcadena $110$.
                \item $q_5$: Se ha leído la subcadena $11$ de la subcadena $110$.
                \item $q_F$: Se ha leído la subcadena $110$. Se han leído ambas subcadenas.
            \end{itemize}

            \item[Opción 2] Primero se lee $110$ y luego $010$. Son los siguientes estados:
            \begin{itemize}
                \item $q_0$: Estado inicial, no ha empezado la subcadena $110$.
                \item $q_1'$: Se ha leído el $1$ de la subcadena $110$.
                \item $q_2'$: Se ha leído la subcadena $11$ de la subcadena $110$.
                \item $q_3'$: Se ha leído la subcadena $110$. Se ha leído el $0$ de la subcadena $010$. Notemos que en este caso podemos agruparlo, puesto que el último carácter de la subcadena $110$ es el mismo que el primero de la subcadena $010$.
                \item $q_4'$: Se ha leído la subcadena $01$ de la subcadena $010$.
                \item $q_F$: Se ha leído la subcadena $010$. Se han leído ambas subcadenas.
            \end{itemize}
        \end{description}

        El autómata obtenido es el de la Figura~\ref{fig:ej:1.2.6.3}.
        \begin{figure}
            \centering
            \begin{tikzpicture}[node distance=2.2cm]
                \node[state, initial] (q0) {$q_0$};
                \node[state, above right of=q0] (q1) {$q_1$};
                \node[state, right of=q1] (q2) {$q_2$};
                \node[state, right of=q2] (q3) {$q_3$};
                \node[state, right of=q3] (q4) {$q_4$};
                \node[state, right of=q4] (q5) {$q_5$};
                \node[state, accepting, below right of=q5] (qF) {$q_F$};

                \node[state, below right of=q0] (q1p) {$q_1'$};
                \node[state, right of=q1p] (q2p) {$q_2'$};
                \node[state, right of=q2p] (q3p) {$q_3'$};
                \node[state, right of=q3p, xshift=2.5em] (q4p) {$q_4'$};
                
                % Conexiones directas por arriba
                \draw   (q0) edge[above] node{0} (q1)
                        (q1) edge[above] node{1} (q2)
                        (q2) edge[above] node{0} (q3)
                        (q3) edge[above] node{1} (q4)
                        (q4) edge[above] node{1} (q5)
                        (q5) edge[above] node{0} (qF)
                        (qF) edge[loop above] node{0, 1} (qF);
                
                % Conexiones directas por abajo
                \draw   (q0) edge[below] node{1} (q1p)
                        (q1p) edge[below] node{1} (q2p)
                        (q2p) edge[below] node{0} (q3p)
                        (q3p) edge[below] node{1} (q4p)
                        (q4p) edge[below] node{0} (qF);

                % Completamos los de arriba
                \draw   (q1) edge[loop above] node{0} (q1)
                        (q2) edge[left] node{1} (q2p)
                        (q3) edge[loop above] node{0} (q3)
                        (q4) edge[bend left, below] node{0} (q3)
                        (q5) edge[loop above] node{1} (q5);

                % Completamos los de abajo
                \draw   (q1p) edge[left] node{0} (q1)
                        (q2p) edge[loop below] node{1} (q2p)
                        (q3p) edge[loop below] node{0} (q3p)
                        (q4p) edge[bend right, above] node{1} (q2p);
                
            \end{tikzpicture}
            \caption{Autómata Finito Determinista del Ejercicio~\ref{ej:1.2.6} apartado~\ref{ej:1.2.6.3}.}
            \label{fig:ej:1.2.6.3}
        \end{figure}
    \end{enumerate}
\end{ejercicio}

\begin{ejercicio} \label{ej:1.2.7}
    Construir un AFD que acepte el lenguaje generado por la siguiente gramática:
    \begin{align*}
        S &\to AB, & A &\to aA, & A &\to c, & B &\to bBb, & B &\to b.
    \end{align*}

    El lenguaje generado por la gramática es:
    \begin{align*}
        L = \{a^ncb^{2m+1} \mid n, m \in \bb{N}\cup \{0\}\}.
    \end{align*}

    El autómata obtenido es el de la Figura~\ref{fig:ej:1.2.7}.
    \begin{figure}
        \centering
        \begin{tikzpicture}
            \node[state, initial] (q0) {$q_0$};
            \node[state, right of=q0] (q1) {$q_1$};
            \node[state, right of=q1, accepting] (q2) {$q_2$};
            \node[state, below of=q1, error] (E) {$E$};

            \draw   (q0) edge[loop above] node{$a$} (q0)
                    (q0) edge[above] node{$c$} (q1)
                    (q0) edge[above] node{$b$} (E)
                    (q1) edge[left] node{$a,c$} (E)
                    (q1) edge[above, bend right] node{$b$} (q2)
                    (q2) edge[above, bend right] node{$b$} (q1)
                    (q2) edge[above] node{$a,c$} (E)
                    (E) edge[loop below] node{$a,b,c$} (E);
        \end{tikzpicture}
        \caption{Autómata Finito Determinista del Ejercicio~\ref{ej:1.2.7}.}
        \label{fig:ej:1.2.7}
    \end{figure}

\end{ejercicio}

\begin{ejercicio} \label{ej:1.2.8}
    Construir un AFD que acepte el lenguaje $L \subseteq \{a, b, c\}^*$ de todas las palabras con un número impar de ocurrencias de la subcadena $abc$.

    El autómata tiene los siguientes estados:
    \begin{itemize}
        \item $q_0$: Llevo un número par de ocurrencias de $abc$, y no he empezado la siguiente.
        \item $q_1$: Acabo de empezar una ocurrencia impar de $abc$, llevo solo una $a$.
        \item $q_2$: Estoy en una ocurrencia impar de $abc$, llevo $ab$.
        \item $q_3$: Llevo un número impar de ocurrencias de $abc$, y no he empezado la siguiente.
        \item $q_4$: Acabo de empezar una ocurrencia par de $abc$, llevo solo una $a$.
        \item $q_5$: Estoy en una ocurrencia par de $abc$, llevo $ab$.
    \end{itemize}

    El autómata obtenido es el de la Figura~\ref{fig:ej:1.2.8}.
    \begin{figure}
        \centering
        \begin{tikzpicture}
            \node[state, initial] (q0) {$q_0$};
            \node[state, right of=q0] (q1) {$q_1$};
            \node[state, right of=q1] (q2) {$q_2$};
            \node[state, below of=q2, accepting] (q3) {$q_3$};
            \node[state, left of=q3, accepting] (q4) {$q_4$};
            \node[state, left of=q4, accepting] (q5) {$q_5$};

            % Conexiones básicas
            \draw   (q0) edge[above] node{$a$} (q1)
                    (q1) edge[above] node{$b$} (q2)
                    (q2) edge[right] node{$c$} (q3)
                    (q3) edge[below] node{$a$} (q4)
                    (q4) edge[below] node{$b$} (q5)
                    (q5) edge[left] node{$c$} (q0);

            % Completamos
            \draw   (q0) edge[loop above] node{$b,c$} (q0)
                    (q1) edge[loop above] node{$a$} (q1)
                    (q1) edge[above, bend right] node{$c$} (q0)
                    (q2) edge[above, bend right] node{$a$} (q1)
                    (q2) edge[below, bend left] node{$b$} (q0)
                    (q3) edge[loop below] node{$b,c$} (q3)
                    (q4) edge[bend right, below] node{$c$} (q3)
                    (q4) edge[loop below] node{$a$} (q4)
                    (q5) edge[bend right, below] node{$a$} (q4)
                    (q5) edge[bend left, above] node{$b$} (q3);
        \end{tikzpicture}
        \caption{Autómata Finito Determinista del Ejercicio~\ref{ej:1.2.8}.}
        \label{fig:ej:1.2.8}
    \end{figure}
\end{ejercicio}

\begin{ejercicio} \label{ej:1.2.9}
    Sea $L$ el lenguaje de todas las palabras sobre el alfabeto $\{0, 1\}$ que no contienen dos $1$s que estén separados por un número impar de símbolos. Describir un AFD que acepte este lenguaje.

    Sea $u\in L$. Veamos que, a lo sumo, puede tener dos $1$'s. Supongamos por reducción al absurdo que tiene tres $1$'s. Entonces, entre la primera y la segunda hay un número impar de símbolos, y entre la segunda y la tercera hay un número impar de símbolos. Por lo tanto, entre el primer y el tercer $1$ hay:
    \begin{itemize}
        \item Un número par de símbolos antes del segundo $1$.
        \item El segundo $1$.
        \item Un número par de símbolos entre el segundo y el tercer $1$.
    \end{itemize}
    Por tanto, como el número de símbolos entre el primer y el tercer $1$ es impar, entonces $u\notin L$. Por lo tanto, $u$ tiene a lo sumo dos $1$'s.\\

    Por tanto, los estados son:
    \begin{itemize}
        \item $q_0$: No se ha introducido ningún $1$.
        \item $q_{p}$: Se ha introducido un $1$, después de él y antes del siguiente $1$ hay un número par de símbolos.
        \item $q_{i}$: Se ha introducido un $1$, después de él y antes del siguiente $1$ hay un número impar de símbolos.
        \item $q_{2}$: Se han introducido dos $1$'s, y no se ha introducido ningún otro.
        \item $E$: Estado de error.
    \end{itemize}

    El autómata obtenido es el de la Figura~\ref{fig:ej:1.2.9}.
    \begin{figure}
        \centering
        \begin{tikzpicture}
            \node[state, initial, accepting] (q0) {$q_0$};
            \node[state, right of=q0, accepting] (qp) {$q_p$};
            \node[state, right of=qp, accepting] (qi) {$q_i$};
            \node[state, below of=qp, accepting] (q2) {$q_2$};
            \node[state, below of=qi, error] (E) {$E$};

            \draw   (q0) edge[loop below] node{0} (q0)
                    (q0) edge[above] node{1} (qp)
                    (qp) edge[below, bend right] node{0} (qi)
                    (qp) edge[right] node{1} (q2)
                    (qi) edge[above, bend right] node{0} (qp)
                    (qi) edge[right] node{1} (E)
                    (q2) edge[loop left] node{0} (q2)
                    (q2) edge[above] node{1} (E)
                    (E) edge[loop right] node{0, 1} (E);
        \end{tikzpicture}
        \caption{Autómata Finito Determinista del Ejercicio~\ref{ej:1.2.9}.}
        \label{fig:ej:1.2.9}
    \end{figure}
\end{ejercicio}

\begin{ejercicio} \label{ej:1.2.10}
    Dada la expresión regular $(a + \varepsilon)b^*$, encontrar un AFND asociado y, a partir de este, calcular un AFD que acepte el lenguaje.

    El AFND con transiciones nulas obtenido (siguiendo el algoritmo) es el de la Figura~\ref{fig:ej:1.2.10}.
    \begin{figure}
        \centering
        \begin{tikzpicture}
            \node[state, initial] (q0) {$q_0$};
            \node[state, above right of=q0] (q1) {$q_1$};
            \node[state, below right of=q0] (q2) {$q_2$};
            \node[state, right of=q1] (q3) {$q_3$};
            \node[state, below right of=q3, accepting] (q4) {$q_4$};
            \node[state, right of=q4] (q5) {$q_5$};
            \node[state, right of=q5, accepting] (q6) {$q_6$};

            \draw   (q0) edge[above] node{$\varepsilon$} (q1)
                    (q0) edge[below] node{$\varepsilon$} (q2)
                    (q1) edge[above] node{$a$} (q3)
                    (q2) edge[above, bend right] node{$\veps$} (q4)
                    (q3) edge[above] node{$\varepsilon$} (q4)
                    (q4) edge[above] node{$\varepsilon$} (q5)
                    (q5) edge[above] node{$b$} (q6)
                    (q6) edge[below, bend left] node{$\veps$} (q5);
        \end{tikzpicture}
        \caption{Autómata Finito No Determinista algorítmico del Ejercicio~\ref{ej:1.2.10}.}
        \label{fig:ej:1.2.10}
    \end{figure}

    Podemos simplificar este autómata para que así la transición al AFD sea más sencilla. El autómata simplificado es el de la Figura~\ref{fig:ej:1.2.10.simplified}.
    \begin{figure}
        \centering
        \begin{tikzpicture}
            \node[state, initial] (q0) {$q_0$};
            \node[state, right of=q0, accepting] (q6) {$q_6$};

            \draw   (q0) edge[above] node{$a, \veps$} (q6)
                    (q6) edge[loop above] node{$b$} (q6);
        \end{tikzpicture}
        \caption{Autómata Finito No Determinista simplificado del Ejercicio~\ref{ej:1.2.10}.}
        \label{fig:ej:1.2.10.simplified}
    \end{figure}

    A partir de este autómata simplificado, obtenemos el AFD de la Figura~\ref{fig:ej:1.2.10.afd}.
    \begin{figure}
        \centering
        \begin{tikzpicture}
            \node[state, initial, accepting] (q0q6) {$\{q_0, q_6\}$};
            \node[state, right of=q0q6, accepting] (q6) {$q_6$};
            \node[state, right of=q6, error] (E) {$E$};

            \draw   (q0q6) edge[above] node{$a,b$} (q6)
                    (q6) edge[loop above] node{$b$} (q6)
                    (q6) edge[above] node{$a$} (E)
                    (E) edge[loop above] node{$a,b$} (E);
        \end{tikzpicture}
        \caption{Autómata Finito Determinista del Ejercicio~\ref{ej:1.2.10}.}
        \label{fig:ej:1.2.10.afd}
    \end{figure}
\end{ejercicio}

\begin{ejercicio} \label{ej:1.2.11}
    Obtener una expresión regular para el lenguaje complementario al aceptado por la gramática
    \begin{align*}
        S &\to abA \mid B \mid baB \mid \varepsilon, & A &\to bS \mid b, & B &\to aS.
    \end{align*}
    \begin{observacion}
        Construir un AFD asociado.
    \end{observacion}
    % // TODO: Hecho en clase. Comparar JJ

    Esta gramática es lineal por la derecha. Algorítmicamente, obtenemos el autómata de la Figura~\ref{fig:ej:1.2.11_L} para el lenguaje generado por $S$, $\cc{L}(S)$.
    \begin{figure}
        \centering
        \begin{tikzpicture}
            % Estados: abA, bA, A, B, baB, aB, B, veps, bS, S, b, aS, S
            \node[state, initial] (S) {$[S]$};
            \node[state, right of=S] (abA) {$[abA]$};
            \node[state, right of=abA] (bA) {$[bA]$};
            \node[state, right of=bA] (A) {$[A]$};
            \node[state, below of=abA] (B) {$[B]$};
            \node[state, below of=B] (baB) {$[baB]$};
            \node[state, right of=baB] (aB) {$[aB]$};
            \node[state, above of=S, accepting] (veps) {$[\varepsilon]$};
            \node[state, below of=A] (bS) {$[bS]$};
            \node[state, right of=A] (b) {$[b]$};
            \node[state, right of=B] (aS) {$[aS]$};

            \draw   (S) edge[above] node{$\veps$} (abA)
                    (abA) edge[above] node{$a$} (bA)
                    (bA) edge[above] node{$b$} (A)
                    (S) edge[above] node{$\veps$} (B)
                    (S) edge[below] node{$\veps$} (baB)
                    (baB) edge[above] node{$b$} (aB)
                    (aB) edge[above] node{$a$} (B)
                    (S) edge[left] node{$\veps$} (veps)
                    (A) edge[left] node{$\veps$} (bS)
                    (A) edge[below] node{$\veps$} (b)
                    (b) edge[above] node{$b$} (veps)
                    (bS) edge[above] node{$b$} (S)
                    (B) edge[above] node{$\veps$} (aS)
                    (aS) edge[below] node{$a$} (S);
        \end{tikzpicture}
        \caption{Autómata Finito Determinista del lenguaje $\cc{L}(S)$ del Ejercicio~\ref{ej:1.2.11}.}
        \label{fig:ej:1.2.11_L}
    \end{figure}

    Ahora, tendríamos que eliminar las transiciones nulas para poder así aplicar el algoritmo para hallar la expresión regular. Esto no es sencillo, por lo que vamos a intentar obtener de forma directa el AFD\@. Para ello, la gramática dada genera el mismo lenguaje que las siguientes reglas de producción, donde hemos eliminado la variable $B$:
    \begin{align*}
        S &\to abA \mid aS \mid baaS \mid \varepsilon, & A &\to bS \mid b
    \end{align*}

    Eliminamos ahora la variable $A$:
    \begin{align*}
        S &\to abbS \mid abb \mid aS \mid baaS \mid \varepsilon
    \end{align*}

    Veamos ahora que la regla $S \to abb$ no es relevante, ya que podemos obtenerla a partir de $S \to abbS$ y $S \to \varepsilon$. Por tanto, la gramática dada inicialmente genera el mismo lenguaje que si estas fuesen las reglas de producción:
    \begin{align*}
        S &\to abbS \mid aS \mid baaS \mid \varepsilon
    \end{align*}

    Por tanto, vemos que:
    \begin{equation*}
        \cc{L}(G)=\{abb, a, baa\}^*.
    \end{equation*}

    En consecuencia, la expresión regular asociada a $\cc{L}(G)$ es:
    \begin{equation*}
        (abb + a + baa)^*
    \end{equation*}

    El AFND asociado a esta expresión regular es el de la Figura~\ref{fig:ej:1.2.11_L_AFND}.
    \begin{figure}
        \centering
        \begin{tikzpicture}
            \node[state, initial, accepting] (q0) {$q_0$};
            \node[state, right of=q0] (q20) {$q_{20}$};
            \node[state, above of=q20] (q10) {$q_{10}$};
            \node[state, below of=q20] (q30) {$q_{30}$};
            \node[state, right of=q10] (q11) {$q_{11}$};
            \node[state, right of=q11] (q12) {$q_{12}$};
            \node[state, right of=q12] (q13) {$q_{13}$};
            \node[state, right of=q20] (q21) {$q_{21}$};
            \node[state, right of=q30] (q31) {$q_{31}$};
            \node[state, right of=q31] (q32) {$q_{32}$};
            \node[state, right of=q32] (q33) {$q_{33}$};

            \draw   (q0) edge[above] node{$\veps$} (q10)
                    (q0) edge[above] node[pos=0.8]{$\veps$} (q20)
                    (q0) edge[above] node{$\veps$} (q30)
                    (q10) edge[above] node{$a$} (q11)
                    (q11) edge[above] node{$b$} (q12)
                    (q12) edge[above] node{$b$} (q13)
                    (q13) edge[above] node{$\veps$} (q0)
                    (q20) edge[above] node{$a$} (q21)
                    (q21) edge[above, bend right] node{$\veps$} (q0)
                    (q30) edge[below] node{$b$} (q31)
                    (q31) edge[below] node{$a$} (q32)
                    (q32) edge[below] node{$a$} (q33)
                    (q33) edge[below] node{$\veps$} (q0);
        \end{tikzpicture}
        \caption{AFND asociado a $\cc{L}(G)$ del Ejercicio~\ref{ej:1.2.11}.}
        \label{fig:ej:1.2.11_L_AFND}
    \end{figure}

    Convirtiendo este autómata en un AFD, obtenemos el de la Figura~\ref{fig:ej:1.2.11_L_AFD}.
    \begin{figure}
        \centering
        \begin{tikzpicture}
            \node[state, initial, accepting] (q0_q10_q20_q30) {$q_0q_{10}q_{20}q_{30}$};
            \node[state, right of=q0_q10_q20_q30, accepting, xshift=5em] (q11_q21_q0_q10_q20_q30) {$q_{11}q_{21}q_{0}q_{10}q_{20}q_{30}$};
            \node[state, right of=q11_q21_q0_q10_q20_q30, xshift=5em] (q12_q31) {$q_{12}q_{31}$};
            \node[state, below of=q0_q10_q20_q30] (q31) {$q_{31}$};
            \node[state, below of=q11_q21_q0_q10_q20_q30] (q32) {$q_{32}$};
            \node[state, below of=q12_q31, error] (E) {$E$};

            \draw   (q0_q10_q20_q30) edge[above] node{$a$} (q11_q21_q0_q10_q20_q30)
                    (q0_q10_q20_q30) edge[left] node{$b$} (q31)
                    (q11_q21_q0_q10_q20_q30) edge[above] node{$b$} (q12_q31)
                    (q11_q21_q0_q10_q20_q30) edge[loop below] node{$a$} (q11_q21_q0_q10_q20_q30)
                    (q12_q31) edge[below, bend right] node{$b$} (q0_q10_q20_q30)
                    (q12_q31) edge[above] node{$a$} (q32)
                    (q31) edge[above] node{$a$} (q32)
                    (q31) edge[bend right, above] node{$b$} (E)
                    (q32) edge[above] node{$a$} (q0_q10_q20_q30)
                    (q32) edge[above] node{$b$} (E)
                    (E) edge[loop above] node{$a,b$} (E);
        \end{tikzpicture}
        \caption{AFD asociado a $\cc{L}(G)$ del Ejercicio~\ref{ej:1.2.11}.}
        \label{fig:ej:1.2.11_L_AFD}
    \end{figure}

    Para que buscar la expresión regular del lenguaje complementario de $\cc{L}(G)$ sea más cómoda, cambiaremos la notación de cada estado.
    El AFD del lenguaje complementario de $\cc{L}(G)$ es el de la Figura~\ref{fig:ej:1.2.11_L_comp}.
    \begin{figure}
        \centering
        \begin{tikzpicture}
            \node[state, initial] (q0_q10_q20_q30) {$p_0$};
            \node[state, right of=q0_q10_q20_q30, xshift=5em] (q11_q21_q0_q10_q20_q30) {$p_1$};
            \node[state, right of=q11_q21_q0_q10_q20_q30, xshift=5em, accepting] (q12_q31) {$p_2$};
            \node[state, below of=q0_q10_q20_q30, accepting] (q31) {$p_3$};
            \node[state, below of=q11_q21_q0_q10_q20_q30, accepting] (q32) {$p_4$};
            \node[state, below of=q12_q31, error, accepting] (E) {$E$};

            \draw   (q0_q10_q20_q30) edge[above] node{$a$} (q11_q21_q0_q10_q20_q30)
                    (q0_q10_q20_q30) edge[left] node{$b$} (q31)
                    (q11_q21_q0_q10_q20_q30) edge[above] node{$b$} (q12_q31)
                    (q11_q21_q0_q10_q20_q30) edge[loop below] node{$a$} (q11_q21_q0_q10_q20_q30)
                    (q12_q31) edge[below, bend right] node{$b$} (q0_q10_q20_q30)
                    (q12_q31) edge[above] node{$a$} (q32)
                    (q31) edge[above] node{$a$} (q32)
                    (q31) edge[bend right, above] node{$b$} (E)
                    (q32) edge[above] node{$a$} (q0_q10_q20_q30)
                    (q32) edge[above] node{$b$} (E)
                    (E) edge[loop above] node{$a,b$} (E);
        \end{tikzpicture}
        \caption{Autómata Finito Determinista del lenguaje $\ol{\cc{L}(G)}$ del Ejercicio~\ref{ej:1.2.11}.}
        \label{fig:ej:1.2.11_L_comp}
    \end{figure}

    Buscamos ahora una expresión para $\ol{\cc{L}(G)}$. Resolvemos el siguiente sistema:
    \begin{equation*}
        \begin{cases}
            p_0 &= ap_1 + bp_3\\
            p_1 &= ap_1 + bp_2\\
            p_2 &= ap_4 + bp_0 + \varepsilon\\
            p_3 &= ap_4 + bE + \varepsilon\\
            p_4 &= ap_0 + bE + \varepsilon\\
            E &= aE + bE + \varepsilon
        \end{cases}
    \end{equation*}

    De la última ecuación, obtenemos que $E = (a+b)^*$. El sistema queda:
    \begin{equation*}
        \begin{cases}
            p_0 &= ap_1 + bp_3\\
            p_1 &= ap_1 + bp_2\\
            p_2 &= ap_4 + bp_0 + \varepsilon\\
            p_3 &= ap_4 + b(a+b)^* + \varepsilon\\
            p_4 &= ap_0 + b(a+b)^* + \varepsilon
        \end{cases}
    \end{equation*}

    Sustituyendo $p_4$, obtenemos:
    \begin{equation*}
        \begin{cases}
            p_0 &= ap_1 + bp_3\\
            p_1 &= ap_1 + bp_2\\
            p_2 &= a[ap_0 + b(a+b)^* + \varepsilon] + bp_0 + \varepsilon = (aa+b)p_0 + ab(a+b)^* + a+\veps\\
            p_3 &= a[ap_0 + b(a+b)^* + \varepsilon] + b(a+b)^* + \varepsilon = aap_0 + (a+\veps)(b(a+b)^* + \veps)\\
        \end{cases}
    \end{equation*}

    Tenemos que:
    \begin{equation*}
        p_1 = ap_1 + bp_2 = a^*bp_2
        = a^*b[(aa+b)p_0 + ab(a+b)^* + a+\veps]
    \end{equation*}

    Sustituyendo, tenemos que:
    \begin{align*}
        p_0 &= a[a^*b[(aa+b)p_0 + ab(a+b)^* + a+\veps]] + b[aap_0 + (a+\veps)(b(a+b)^* + \veps)] =\\
        &= a^+b(aa+b)p_0 + a^+bab(a+b)^* + a^+ba + a^+b + baap_0 + b(a+\veps)(b(a+b)^* + \veps) =\\
        &= [a^+b(aa+b)+baa]p_0 + a^+bab(a+b)^* + a^+ba + a^+b + b(a+\veps)(b(a+b)^* + \veps) \AstIg\\
        & \AstIg [a^+b(aa+b)+baa]^*[a^+bab(a+b)^* + a^+ba + a^+b + b(a+\veps)(b(a+b)^* + \veps)]
    \end{align*}
    donde en $(\ast)$ hemos aplicado el Lema de Arden.
    Por tanto, la expresión regular asociada a $\ol{\cc{L}(G)}$ es:
    \begin{equation*}
        [a^+b(aa+b)+baa]^*[a^+bab(a+b)^* + a^+ba + a^+b + b(a+\veps)(b(a+b)^* + \veps)]
    \end{equation*}
\end{ejercicio}

\begin{ejercicio} \label{ej:1.2.12}
    Dar expresiones regulares para los lenguajes sobre el alfabeto $\{a, b\}$ dados por las siguientes condiciones:
    \begin{enumerate}
        \item Palabras que no contienen la subcadena $a$,
        \begin{equation*}
            b^*
        \end{equation*}
        \item Palabras que no contienen la subcadena $ab$.
        \begin{equation*}
            b^*a^*
        \end{equation*}
        \item \label{ej:1.2.12_3} % // TODO: Comprobar JJ
        Palabras que no contienen la subcadena $aba$.
        
        Este lenguaje viene descrito por el autómata de la Figura~\ref{fig:ej:1.2.12_3}.
        \begin{figure}
            \centering
            \begin{tikzpicture}
                \node[state, initial, accepting] (q0) {$q_0$};
                \node[state, right of=q0, accepting] (q1) {$q_1$};
                \node[state, right of=q1, accepting] (q2) {$q_2$};
                \node[state, right of=q2, error] (q3) {$E$};

                \draw   (q0) edge[above] node{$a$} (q1)
                        (q1) edge[above] node{$b$} (q2)
                        (q2) edge[above] node{$a$} (q3)
                        (q0) edge[loop above] node{$b$} (q0)
                        (q1) edge[loop above] node{$a$} (q1)
                        (q2) edge[below, bend left] node{$b$} (q0)
                        (q3) edge[loop above] node{$a,b$} (q3);
            \end{tikzpicture}
            \caption{AFD del lenguaje del Ejercicio~\ref{ej:1.2.12} apartado~\ref{ej:1.2.12_3}.}
            \label{fig:ej:1.2.12_3}
        \end{figure}

        Obtenemos la expresión regular asociada al lenguaje del autómata de la Figura~\ref{fig:ej:1.2.12_3}.
        \begin{equation*}
            \begin{cases}
                q_0 &= bq_0 + aq_1 + \varepsilon\\
                q_1 &= aq_1 + bq_2 + \varepsilon\\
                q_2 &= bq_0 + aE + \varepsilon\\
                E &= aE + bE = (a+b)E + \emptyset
            \end{cases}
        \end{equation*}

        Usando el Lema de Arden, obtenemos que $E=\emptyset(a+b)^*=\emptyset$. Sustituyendo, obtenemos:
        \begin{equation*}
            \begin{cases}
                q_0 &= bq_0 + aq_1 + \varepsilon\\
                q_1 &= aq_1 + bq_2 + \varepsilon\\
                q_2 &= bq_0 + a\emptyset + \varepsilon = bq_0 + \varepsilon
            \end{cases}
        \end{equation*}

        Sustituyendo $q_2$, obtenemos:
        \begin{equation*}
            \begin{cases}
                q_0 &= bq_0 + aq_1 + \varepsilon\\
                q_1 &= aq_1 + b(bq_0 + \varepsilon) + \varepsilon\\
            \end{cases}
        \end{equation*}

        Usando el Lema de Arden, obtenemos que:
        \begin{equation*}
            q_1 = a^*[b(bq_0 + \varepsilon) + \varepsilon]
        \end{equation*}

        Sustituyendo en la primera ecuación, tenemos que:
        \begin{align*}
            q_0 &= bq_0 +aa^*[b(bq_0 + \varepsilon) + \varepsilon] + \veps=\\
            &= bq_0 + a^+[bbq_0 + b + \varepsilon] + \veps =\\
            &= (b+a^+bb)q_0 + a^+[b + \varepsilon] + \veps \AstIg\\
            &\AstIg (b+a^+bb)^*[a^+(b + \varepsilon)+\varepsilon]
        \end{align*}
        donde en $(\ast)$ hemos aplicado el Lema de Arden.
        Por tanto, la expresión regular asociada al lenguaje del autómata de la Figura~\ref{fig:ej:1.2.12_3} es:
        \begin{equation*}
            (b+a^+bb)^*[a^+(b + \varepsilon)+\varepsilon]
        \end{equation*}
    \end{enumerate}
\end{ejercicio}

\begin{ejercicio}
    Determinar si el lenguaje generado por la siguiente gramática es regular:
    \begin{align*}
        S &\to AabB, & A &\to aA \mid bA \mid \varepsilon, & B &\to Bab \mid Bb \mid ab \mid b.
    \end{align*}
    En caso de que lo sea, encontrar una expresión regular asociada.\\

    Es directo ver que el lenguaje generado por la gramática tiene como expresión regular asociada:
    \begin{align*}
        (a+b)^*ab(ab+b)^+.
    \end{align*}

    Por tanto, el lenguaje es regular.
\end{ejercicio}

\begin{ejercicio} \label{ej:1.2.14}
    Sobre el alfabeto $A = \{0, 1\}$ realizar las siguientes tareas:
    \begin{enumerate}
        \item \label{ej:1.2.14.1}
        Describir un autómata finito determinista que acepte todas las palabras que contengan a $011$ o a $010$ (o las dos) como subcadenas.
        
        Tenemos los siguientes estados:
        \begin{itemize}
            \item \ul{$q_0$}: No se ha empezado ninguna subcadena.
            \item \ul{$q_1$}: Se ha empezado una subcadena deseada. Tengo el carácter $0$.
            \item \ul{$q_2$}: Se continúa la subcadena deseada. Tengo los caracteres $01$.
            \item \ul{$q_3$}: Se ha encontrado la subcadena deseada. Tengo los caracteres $011$ o $010$.
        \end{itemize}

        El autómata obtenido es el de la Figura~\ref{fig:ej:1.2.14.1}.
        \begin{figure}
            \centering
            \begin{tikzpicture}
                \node[state, initial] (q0) {$q_0$};
                \node[state, right of=q0] (q1) {$q_1$};
                \node[state, right of=q1] (q2) {$q_2$};
                \node[state, accepting, right of=q2] (q3) {$q_3$};

                \draw   (q0) edge[above] node{0} (q1)
                        (q1) edge[above] node{1} (q2)
                        (q2) edge[above] node{0,1} (q3)
                        (q1) edge[loop above] node{0} (q1)
                        (q0) edge[loop above] node{1} (q0)
                        (q3) edge[loop above] node{0,1} (q3);
            \end{tikzpicture}
            \caption{Autómata Finito Determinista del Ejercicio~\ref{ej:1.2.14} apartado~\ref{ej:1.2.14.1}.}
            \label{fig:ej:1.2.14.1}
        \end{figure}
        
        \item \label{ej:1.2.14.2}
        Describir un autómata finito determinista que acepte todas las palabras que empiecen o terminen (o ambas cosas) por $01$.
        
        Tenemos los siguientes estados:
        \begin{itemize}
            \item \ul{$q_0$}: No hemos leído nada.
            \item \ul{$q_1$}: Hemos empezado con un $0$, por lo que puede comenzar por $01$ (o terminar por $01$).
            \item \ul{$q_2$}: Hemos empezado con $01$, por lo que ya no hay más restricciones.
            \item \ul{$q_3$}: No hemos empezado por $01$, por lo que ha de terminar por $01$.
            \item \ul{$q_4$}: Ha de terminar por $01$, y estamos en $0$, por lo que si introduce un $1$ puede terminar.
            \item \ul{$q_5$}: Ha de terminar por $01$, y acabamos de leer $01$, por lo que podemos terminar.
        \end{itemize}

        El autómata obtenido es el de la Figura~\ref{fig:ej:1.2.14.2}.
        \begin{figure}
            \centering
            \begin{tikzpicture}
                \node[state, initial] (q0) {$q_0$};
                \node[state, above right of=q0] (q1) {$q_1$};
                \node[state, right of=q1, accepting] (q2) {$q_2$};
                \node[state, below right of=q0] (q3) {$q_3$};
                \node[state, right  of=q3] (q4) {$q_4$};
                \node[state, accepting, right of=q4] (q5) {$q_5$};

                % Uniones directas
                \draw   (q0) edge[above] node{0} (q1)
                        (q0) edge[above] node{1} (q3)
                        (q1) edge[above] node{1} (q2)
                        (q2) edge[loop above] node{0,1} (q2)
                        (q3) edge[above] node{0} (q4)
                        (q4) edge[above] node{1} (q5);
                
                % Completamos autómata
                \draw   (q1) edge[left] node{0} (q4)
                        (q3) edge[loop below] node{1} (q3)
                        (q4) edge[loop above] node{0} (q4)
                        (q5) edge[above, bend right] node{0} (q4)
                        (q5) edge[below, bend left] node{1} (q3);
            \end{tikzpicture}
            \caption{Autómata Finito Determinista del Ejercicio~\ref{ej:1.2.14} apartado~\ref{ej:1.2.14.2}.}
            \label{fig:ej:1.2.14.2}
        \end{figure}

        \item Dar una expresión regular para el conjunto de las palabras en las que hay dos ceros separados por un número de símbolos que es múltiplo de $4$ (los símbolos que separan los ceros pueden ser ceros y puede haber otros símbolos delante o detrás de estos dos ceros).
        \begin{equation*}
            (0+1)^*\red{0}\left((0+1)(0+1)(0+1)(0+1)\right)^*\red{0}(0+1)^*
        \end{equation*}
        Notemos que los dos $0$'s en cuestión están marcados en rojo para facilitar la comprensión.
        \item Dar una expresión regular para las palabras en las que el número de ceros es divisible por $4$.
        
        En un primer momento, podríamos pensar en:
        \begin{equation*}
            (1^*01^*01^*01^*01^*)^*
        \end{equation*}
        No obstante, una palabra con $1$'s y sin $0$'s, que es aceptada por el lenguaje, no está contemplada en la expresión regular. La expresión regular correcta es:
        \begin{equation*}
            (1^*01^*01^*01^*0)^*1^*
        \end{equation*}
    \end{enumerate}
\end{ejercicio}

\begin{ejercicio}~\label{ej:1.2.15}
    Construye una gramática regular que genere el siguiente lenguaje:
    \begin{align*}
        L_1 = \{u \in \{0, 1\}^* \mid \text{el número de $1$'s y de $0$'s es impar}\}.
    \end{align*}

    Tenemos los siguientes estados:
    \begin{itemize}
        \item \ul{$E_{01}$}: Tenemos un error en $0$ y $1$, ya que el número de $0$'s y de $1$'s es par.
        \item \ul{$E_0$}: Tenemos un error en $0$, ya que el número de $0$'s es par. El número de $1$'s es impar.
        \item \ul{$E_1$}: Tenemos un error en $1$, ya que el número de $1$'s es par. El número de $0$'s es impar.
        \item \ul{$X$}: No tenemos errores. El número de $0$'s y de $1$'s es impar.
    \end{itemize}

    La gramática obtenida es $G=(\{E_{01},E_0,E_1,X\}, \{0, 1\}, P, E_{01})$, donde $P$ es:
    \begin{align*}
        E_{01} &\to 0E_1 \mid 1E_0, \\
        E_0 &\to 0X \mid 1E_{01}, \\
        E_1 &\to 0E_{01} \mid 1X, \\
        X &\to 0E_0 \mid 1E_1 \mid \veps
    \end{align*}

\end{ejercicio}

\begin{ejercicio}
    Encuentra una expresión regular que represente el siguiente lenguaje:
    \begin{align*}
        L_2 = \{0^n1^m \mid n \geq 1, m \geq 0, n \text{\ múltiplo de\ } 3 \text{\ y\ } m \text{\ es par}\}.
    \end{align*}

    La expresión regular es:
    \begin{align*}
        (000)^+(11)^*
    \end{align*}
\end{ejercicio}

\begin{ejercicio} \label{ej:1.2.17}
    Diseña un autómata finito determinista que reconozca el siguiente lenguaje:
    \begin{align*}
        L_3 = \{u \in \{0, 1\}^* \mid \text{el número de $1$'s no es múltiplo de\ } 3 \text{\ y el número de $0$'s es par}\}.
    \end{align*}

    Sean $n_0$ el número de $0$'s y $n_1$ el número de $1$'s.

    Tenemos la siguiente disposición de estados:
    \begin{itemize}
        \item Los estados de arriba representan $n_0 \text{ mod } 2 = 0$.
        \item Los estados de abajo representan $n_0 \text{ mod } 2 = 1$.
        \item Los estados de la primera columna representan $n_1 \text{ mod } 3 = 0$.
        \item Los estados de la segunda columna representan $n_1 \text{ mod } 3 = 1$.
        \item Los estados de la tercera columna representan $n_1 \text{ mod } 3 = 2$.
    \end{itemize}

    El estado $q_{ij}$ representa $n_0 \text{ mod } 2 = i$ y $n_1 \text{ mod } 3 = j$.
    
    El autómata obtenido es el de la Figura~\ref{fig:ej:1.2.17}.
    \begin{figure}
        \centering
        \begin{tikzpicture}
            \node[state, initial] (q00) {$q_{00}$};
            \node[state, right of=q00, accepting] (q01) {$q_{01}$};
            \node[state, right of=q01, accepting] (q02) {$q_{02}$};
            \node[state, below of=q00] (q10) {$q_{10}$};
            \node[state, right of=q10] (q11) {$q_{11}$};
            \node[state, right of=q11] (q12) {$q_{12}$};

            \draw   (q00) edge[above] node{1} (q01)
                    (q01) edge[above] node{1} (q02)
                    (q02) edge[above, bend right] node{1} (q00)
                    (q10) edge[above] node{1} (q11)
                    (q11) edge[above] node{1} (q12)
                    (q12) edge[below, bend left] node{1} (q10)
                    (q00) edge[left] node{0} (q10)
                    (q10) edge[right, bend right] node{0} (q00)
                    (q01) edge[left] node{0} (q11)
                    (q11) edge[right, bend right] node{0} (q01)
                    (q02) edge[left] node{0} (q12)
                    (q12) edge[right, bend right] node{0} (q02);
        \end{tikzpicture}
        \caption{Autómata Finito Determinista del Ejercicio~\ref{ej:1.2.17}.}
        \label{fig:ej:1.2.17}
    \end{figure}
\end{ejercicio}

\begin{ejercicio} \label{ej:1.2.18}
    Dar una expresión regular para el lenguaje aceptado por el autómata de la Figura~\ref{fig:ej:1.2.18}.
    \begin{figure}
        \centering
        \begin{tikzpicture}
            \node[state, initial, accepting] (q0) {$q_0$};
            \node[state, accepting, below right of=q0] (q2) {$q_2$};
            \node[state, above right of=q2] (q1) {$q_1$};

            \draw   (q1) edge[loop above] node{$a$} (q1)
                    (q0) edge[above] node{$a,b$} (q1)
                    (q1) edge[below, bend left] node{$b$} (q2)
                    (q2) edge[below, bend left] node{$b$} (q1)
                    (q2) edge[above] node{$a$} (q0);
        \end{tikzpicture}
        \caption{Autómata Finito Determinista del Ejercicio~\ref{ej:1.2.18}.}
        \label{fig:ej:1.2.18}
    \end{figure}

    Establecemos una ecuación por cada uno de los estados. El sistema inicial es:
    \begin{equation*}
        \begin{cases}
            q_0 = \veps + aq_1 + bq_1, \\
            q_1 = aq_1 + bq_2, \\
            q_2 = \veps + aq_0 + bq_1.
        \end{cases}
    \end{equation*}

    Buscamos obtener la expresión regular asociada a $q_1$:
    \begin{align*}
        q_1 &= aq_1 + b + baq_0 +bbq_1=\\
        &= baq_0 + b + (a+bb)q_1\AstIg\\
        &\AstIg (a+bb)^*(baq_0 + b)
    \end{align*}
    donde en $(\ast)$ hemos aplicado el Lema de Arden. Sustituyendo en la ecuación de $q_0$ obtenemos:
    \begin{align*}
        q_0 &= \veps + (a+b)q_1 =\\
        &= \veps + (a+b)(a+bb)^*(baq_0 + b) =\\
        &= \veps + (a+b)(a+bb)^*b + (a+b)(a+bb)^*baq_0 \AstIg\\
        &\AstIg ((a+b)(a+bb)^*ba)^*(\veps + (a+b)(a+bb)^*b)
    \end{align*}
    donde, de nuevo, en $(\ast)$ hemos aplicado el Lema de Arden. Por tanto, la expresión regular asociada al autómata es:
    \begin{align*}
        ((a+b)(a+bb)^*ba)^*(\veps + (a+b)(a+bb)^*b).
    \end{align*}
\end{ejercicio}


\begin{ejercicio} \label{ej:1.2.19}
    Dado el lenguaje
    \begin{align*}
        L = \{u110 \mid u \in \{1, 0\}^*\},
    \end{align*}
    encontrar la expresión regular, la gramática lineal por la derecha, la gramática lineal por la izquierda y el AFD asociado.

    La expresión regular es:
    \begin{align*}
        (0+1)^*110.
    \end{align*}

    La gramática lineal por la derecha es $G=(\{S,A\}, \{0, 1\}, P, S)$, donde $P$ es:
    \begin{align*}
        S &\to 0S \mid 1S \mid A\\
        A &\to 110.
    \end{align*}

    La gramática lineal por la izquierda es $G=(\{S,A\}, \{0, 1\}, P', S)$, donde $P'$ es:
    \begin{align*}
        S &\to X110\\
        X &\to X0 \mid X1 \mid \varepsilon.
    \end{align*}

    El AFD asociado es el de la Figura~\ref{fig:ej:1.2.19}. Sus estados son:
    \begin{itemize}
        \item \ul{$q_0$}: No estoy en la cadena $110$ final.
        \item \ul{$q_1$}: He leído un $1$ de la cadena final.
        \item \ul{$q_2$}: He leído un $11$ de la cadena final.
        \item \ul{$q_3$}: He leído un $110$ de la cadena final.
    \end{itemize}
    \begin{figure}
        \centering
        \begin{tikzpicture}
            \node[state, initial] (q0) {$q_0$};
            \node[state, right of=q0] (q1) {$q_1$};
            \node[state, right of=q1] (q2) {$q_2$};
            \node[state, accepting, right of=q2] (q3) {$q_3$};

            \draw   (q0) edge[above] node{1} (q1)
                    (q1) edge[above] node{1} (q2)
                    (q2) edge[above] node{0} (q3)
                    (q0) edge[loop above] node{0} (q0)
                    (q1) edge[above, bend right] node{0} (q0)
                    (q2) edge[loop above] node{1} (q2)
                    (q3) edge[below, bend left] node{0} (q0)
                    (q3) edge[below, bend left] node[pos=0.9]{1} (q1);
        \end{tikzpicture}
        \caption{Autómata Finito Determinista del Ejercicio~\ref{ej:1.2.19}.}
        \label{fig:ej:1.2.19}
    \end{figure}
\end{ejercicio}

\begin{ejercicio}
    Dado un AFD, determinar el proceso que habría que seguir para construir una gramática lineal por la izquierda capaz de generar el Lenguaje aceptado por dicho autómata.\\

    Sea $M=(Q, A, \delta, q_0, F)$ un AFD\@. Como $Q$ es finito, podemos enumerar los estados como $Q=\{q_1, q_2, \ldots, q_n\}$.    
    La Gramática Lineal por la Izquierda asociada es $G=(Q\cup \{S\}, A, P, S)$, donde hemos supuesto $S\notin Q$ debido a nuestra enumeración de los estados. Las reglas de producción son:
    \begin{equation*}
        P=\left\{
            \begin{array}{rcll}
                S &\to& q_i &\qquad \forall q_i\in F\\
                q_i &\to& q_j a &\qquad \forall q_i,q_j\in Q, a\in A \mid \delta(q_j, a)=q_i \\
                q_0 &\to& \veps.
            \end{array}
        \right.
    \end{equation*}
    Notemos que lo que hacemos es invertir el autómata, obtener la gramática lineal por la derecha, y después invertir las reglas de producción de esta última.
\end{ejercicio}

\begin{ejercicio} \label{ej:1.2.21}
    Construir un autómata finito determinista que acepte el lenguaje de todas las palabras sobre el alfabeto $\{0, 1\}$ que no contengan la subcadena $001$.
    Construir una gramática regular por la izquierda a partir de dicho autómata.\\

    Los estados son los siguientes:
    \begin{itemize}
        \item \ul{$q_0$}: No se ha empezado la subcadena $001$
        \item \ul{$q_1$}: Se ha leído un $0$ de la subcadena $001$.
        \item \ul{$q_2$}: Se ha leído un $00$ de la subcadena $001$.
        \item \ul{$E$}: Se ha leído la subcadena $001$, por lo que es el estado de error.
    \end{itemize}

    El autómata obtenido es el de la Figura~\ref{fig:ej:1.2.21}.
    \begin{figure}
        \centering
        \begin{tikzpicture}
            \node[state, initial, accepting] (q0) {$q_0$};
            \node[state, right of=q0, accepting] (q1) {$q_1$};
            \node[state, right of=q1, accepting] (q2) {$q_2$};
            \node[state, error, right of=q2] (E) {$E$};

            \draw   (q0) edge[below] node{0} (q1)
                    (q0) edge[loop above] node{1} (q0)
                    (q1) edge[below] node{0} (q2)
                    (q1) edge[above, bend right] node{1} (q0)
                    (q2) edge[loop above] node{0} (q0)
                    (q2) edge[below] node{1} (E)
                    (E) edge[loop above] node{0,1} (E);
        \end{tikzpicture}
        \caption{Autómata Finito Determinista del Ejercicio~\ref{ej:1.2.21}.}
        \label{fig:ej:1.2.21}
    \end{figure}
    
    Respecto a la gramática regular por la izquierda, usando el algoritmo descrito en el apartado anterior, tenemos que la gramática es $G=(Q\cup \{S\}, \{0, 1\}, P, S)$, donde $P$ es:
    \begin{align*}
        P=\left\{
            \begin{array}{rcll}
                S &\to& q_0 \mid q_1 \mid q_2,\\
                q_0 &\to& q_0 1 \mid q_1 1 \red{\mid \veps}\\
                q_1 &\to& q_0 0,\\
                q_2 &\to& q_1 0 \mid q_2 0,\\
                E &\to& E0 \mid E1 \mid q_2 1,
            \end{array}
        \right.
    \end{align*}

    % // TODO: Revisar JJ
\end{ejercicio}

\begin{ejercicio}
    Sea $B_n = \{a^k \mid k \text{\ es múltiplo de\ } n\}$. Demostrar que $B_n$ es regular para todo $n$.\\

    Fijado $n\in \bb{N}$, la expresión regular correspondiente es:
    \begin{align*}
        (a\overbrace{\cdots }^{n \text{\ veces}}a)^* = (a^n)^*
    \end{align*}

    Equivalentemente, usando la notación de las expresiones regulares de UNIX, la expresión regular sería:
    \begin{equation*}
        (a\{n\})^*
    \end{equation*}
\end{ejercicio}

\begin{ejercicio}
    Sea $A$ un alfabeto. Decimos que $u\in A^*$ es un prefijo de $v\in A^*$ si existe $w\in A^*$ tal que $uw = v$. Decimos que $u$ es un prefijo propio de $v$ si además $u \neq v$ y $u \neq \varepsilon$. Demostrar que si $L$ es regular, también lo son los lenguajes siguientes:
    \begin{enumerate}
        \item $\text{NOPREFIJO}(L) = \{u \in L \mid \text{ningún prefijo propio de } u \text{ pertenece a } L\}$,
        
        Como $L$ es regular, existe un AFD $M=(Q, A, \delta, q_0, F)$ tal que $L=\cc{L}(M)$. Construimos un AFD $M'=(Q\cup \{E\}, A, \delta', q_0, F)$,
        donde $E$ es un estado de error ($E\notin Q$) y $\delta'$ es:
        \begin{equation*}
            \begin{cases}
                \delta'(q, a) = \delta(q, a)&\qquad \forall q\in Q\setminus F, a\in A\\
                \delta'(q, a) = E&\qquad \forall q\in F, a\in A\\
                \delta'(E, a) = E&\qquad \forall a\in A
            \end{cases}
        \end{equation*}

        Demostramos mediante doble inclusión que $\text{NOPREFIJO}(L)=\cc{L}(M')$.
        \begin{description}
            \item[$\subseteq)$] Sea $u\in \text{NOPREFIJO}(L)$. Entonces, por definición de $\text{NOPREFIJO}(L)$, $u\in L$. Por tanto, $\exists q\in F$ tal que $\delta^{\ast}(q_0, u)=q$. Para ver que $u\in \cc{L}(M')$, basta ver que $(\delta')^{\ast}(q_0, u)\in F$.

            Como $u$ no tiene prefijos propios en $L$, entonces $\delta^{\ast}(q_0, u')\notin F$ para todo prefijo propio $u'$ de $u$; es decir, en los pasos de cálculo desde $q_0$ hasta $\delta^{\ast}(q_0, u)$ no se pasa por ningún estado final. Por tanto, como en esos casos $\delta'=\delta$, entonces $\delta^{\ast}(q_0, u)=q\in F$, por lo que $u\in \cc{L}(M')$.

            \item[$\supseteq)$] Sea $u\in \cc{L}(M')$. En primer lugar, tenemos que $(\delta')^{\ast}(q_0, u)\in F$. Veamos ahora que los pasos de cálculo desde $q_0$ hasta $(\delta')^{\ast}(q_0, u)$ leyendo $u$ no son ninguno finales.
            
            Si alguno de ellos fuese final (si $u$ tuviese algún prefijo propio $v\in L$), entonces desde él pasaríamos a $E$, y de este estado no final no saldríamos, llegando a contradicción. Por tanto, $u$ no tiene prefijos propios pertenecientes a $L$. Además, como en estos casos $\delta'=\delta$, tenemos que $\delta^{\ast}(q_0, u)\in F$, luego $u\in L$. De esta forma, $u\in \text{NOPREFIJO}(L)$.
        \end{description}
        
        \item $\text{NOEXTENSION}(L) = \{u \in L \mid u \text{ no es un prefijo propio de ninguna palabra de } L\}$.
        
        Como $L$ es regular, existe un AFD $M=(Q, A, \delta, q_0, F)$ tal que $L=\cc{L}(M)$. Construimos un AFD $M'=(Q, A, \delta, q_0, F')$, donde:
        \begin{equation*}
            F'=\{q\in F\mid \delta^{\ast}(q,u)\notin F,\forall u\in A^*\}
        \end{equation*}

        Demostramos mediante doble inclusión que $\text{NOEXTENSION}(L)=\cc{L}(M')$.
        \begin{description}
            \item[$\subseteq)$] Sea $u\in \text{NOEXTENSION}(L)$. Entonces, por definición de $\text{NOEXTENSION}(L)$, $u\in L$. Por tanto, $\exists q\in F$ tal que $\delta^{\ast}(q_0, u)=q$. Para ver que $u\in \cc{L}(M')$, basta ver que $q\in F'$. 
            
            Supongamos por reducción al absurdo $q\notin F'$. Entonces, $\exists v\in A^*$ tal que $\delta^{\ast}(q, v)\in F$. Pero entonces, $\delta^{\ast}(q_0, uv)=\delta^{\ast}(\delta^{\ast}(q_0, u), v)=\delta^{\ast}(q, v)\in F$, por lo que $uv\in L$ y, por tanto, $u$ es prefijo propio de $uv$, lo cual es una contradicción. Por tanto, $q\in F'$ y, por tanto, $u\in \cc{L}(M')$.

            \item[$\supseteq)$] Sea $u\in \cc{L}(M')$. En primer lugar, tenemos que $\delta^{\ast}(q_0, u)=q\in F'\subset F$, luego $u\in L$. Veamos ahora que $u$ no es prefijo propio de ninguna palabra de $L$.
            
            Supongamos por reducción al absurdo que $u$ es prefijo propio de alguna palabra de $L$. Entonces, $\exists v\in A^*\setminus \{\veps\}$ tal que $uv\in L$. Por tanto, $\delta^{\ast}(q_0, uv)\in F$. Pero entonces, $\delta^{\ast}(q_0, uv)=\delta^{\ast}(\delta^{\ast}(q_0, u), v)=\delta^{\ast}(q, v)\in F$.
            No obstante, hemos demostrado entonces que $q\notin F'$, lo cual es una contradicción. Por tanto, $u$ no es prefijo propio de ninguna palabra de $L$ y, por tanto, $u\in \text{NOEXTENSION}(L)$.
        \end{description}
    \end{enumerate}
\end{ejercicio}

\begin{ejercicio}
    Si $L \subseteq A^*$, define la relación $\equiv$ en $A^*$ como sigue: si $u, v \in A^*$, entonces $u \equiv v$ si y solo si para toda $z \in A^*$, tenemos que $(uz \in L \Leftrightarrow vz \in L)$.
    \begin{enumerate}
        \item Demostrar que $\equiv$ es una relación de equivalencia.
        
        Veamos las tres propiedades de las relaciones de equivalencia:
        \begin{itemize}
            \item \ul{Reflexiva}: Sea $u\in A^*$. Entonces, para todo $z\in A^*$, tenemos trivialmente que $(uz\in L\Leftrightarrow uz\in L)$. Por tanto, $u\equiv u$.
            \item \ul{Simétrica}: Sean $u,v\in A^*$ tales que $u\equiv v$. Entonces, para todo $z\in A^*$, tenemos que $(uz\in L\Leftrightarrow vz\in L)$. Por tanto, para todo $z\in A^*$, tenemos que $(vz\in L\Leftrightarrow uz\in L)$, lo cual implica que $v\equiv u$.
            \item \ul{Transitiva}: Sean $u,v,w\in A^*$ tales que $u\equiv v$ y $v\equiv w$. Entonces, para todo $z\in A^*$, tenemos que $(uz\in L\Leftrightarrow vz\in L)$ y $(vz\in L\Leftrightarrow wz\in L)$. Por tanto, para todo $z\in A^*$, tenemos que $(uz\in L\Leftrightarrow wz\in L)$, lo cual implica que $u\equiv w$.
        \end{itemize}
        Tenemos por tanto que $\equiv$ es una relación de equivalencia.

        \item Calcular las clases de equivalencia de $L = \{a^ib^i \mid i \geq 0\}$.
        
        En este caso, $A=\{a,b\}$.
        La primera clase de equivalencia que encontramos es las palabras que, le añadamos al final lo que le añadamos, no pertenecen al lenguaje. Es decir:
        \begin{equation*}
            [u\in A^* \mid \text{en $u$ hay una $a$ después de una $b$} \qquad \lor\qquad u=a^ib^j,~j>i]
        \end{equation*}
        Por comodidad, ya que $b\notin L$, considerando este representante de clase de equivalencia, notaremos a esta clase de equivalencia con $[b]$.
        Además, para cada $k\in \bb{N}\cup \{0\}$, tenemos:
        \begin{equation*}
            [a_k] =: \{a^{k+j}b^j\mid j\in \bb{N}\cup \{0\}\}
        \end{equation*}

        Veamos en primer lugar que no hay más clases de equivalencia. Dado $u\in A^*$, si $u$ tiene una $a$ después de una $b$, entonces $u\in [b]$. En caso contrario, tenemos $u=a^ib^j$ con $i,j\in \bb{N}\cup \{0\}$.
        \begin{itemize}
            \item Si $j>i$, entonces $u\in [b]$.
            \item Si $j\leq i$, entonces $u\in [a_{i-j}]$.
        \end{itemize}
        Por tanto, tenemos que no hay más clases de equivalencia. Veamos ahora que estas clases de equivalencia son disjuntas.
        \begin{itemize}
            \item Sea $u\in [b]$. Si $u$ tiene una $a$ después de una $b$, entonces de forma directa $\nexists k\in \bb{N}\cup \{0\}$ tal que $u\in [a_k]$, ya que las palabras de estas clase de equivalencia son casos particulares de $a^*b^*$. Por otro lado, si $u=a^ib^j$ con $j>i$, entonces para que sea de la forma $a^{k+j}b^j$ es necesario que $k+j=i$, por lo que $k=i-j<0$, luego $\nexists k\in \bb{N}\cup \{0\}$ tal que $u\in [a_k]$.
            
            En conclusión, vemos que $[b]\cap [a_k]=\emptyset$ para todo $k\in \bb{N}\cup \{0\}$.

            \item Sean $k_1,k_2\in \bb{N}\cup \{0\}$ tales que $k_1\neq k_2$. Supongamos que $[a_{k_1}]\cap [a_{k_2}]\neq \emptyset$. Entonces, $\exists u\in A^*$ tal que $u\in [a_{k_1}]\cap [a_{k_2}]$. Por tanto, tenemos que $u=a^{k_1+j_1}b^{j_1}=a^{k_2+j_2}b^{j_2}$ con $j_1,j_2\in \bb{N}\cup \{0\}$. Igualando las potencias de $a$ y $b$, tenemos que $k_1+j_1=k_2+j_2$ y $j_1=j_2$. Por tanto, $k_1=k_2$, lo cual es una contradicción. Por tanto, $[a_{k_1}]\cap [a_{k_2}]=\emptyset$ para todo $k_1,k_2\in \bb{N}\cup \{0\}$ tales que $k_1\neq k_2$.
        \end{itemize}

        Por tanto, vemos que las clases de equivalencia de $L$ son $[b]$ y $[a_k]$ para todo $k\in \bb{N}\cup \{0\}$. Notemos que:
        \begin{itemize}
            \item $[b]$ es la clase de equivalencia de las palabras que, le añadamos al final lo que le añadamos, no pertenecen al lenguaje. Para cualquier par $u,v\in [b]$, tenemos que, para todo $z\in A^*$, $uz\in L\Leftrightarrow vz\in L$ es cierto por vacuidad, puesto que en ambos casos $uz,vz\notin L$.
            \item $[a_0]$ es la clase de equivalencia de las palabras que pertenecen al lenguaje. Para cualquier par $u,v\in [a_0]$, tenemos que, para todo $z\in A^*$, $uz\in L\Leftrightarrow vz\in L$ es cierto. Tomando $z=\veps$, tenemos que $u=uz,v=vz\in L$; mientras que si $z\neq \veps$, entonces $uz,vz\notin L$.
            \item $[a_k]$ para $k\neq 0$ es la clase de equivalencia de las palabras que tan solo pertenecen al lenguaje al concatenarles $b^k$. Para cualquier par $u,v\in [a_k]$, tenemos que, para todo $z\in A^*$, $uz\in L\Leftrightarrow vz\in L$ es cierto. Tomando $z=b^k$, tenemos que $ub^k,vb^k\in L$; mientras que si $z\neq b^k$, entonces $uz,vz\notin L$.
        \end{itemize}

        \item Calcular las clases de equivalencia de $L = \{a^ib^j \mid i, j \geq 0\}$.
        
        La primera clase de equivalencia que encontramos es las palabras que, le añadamos al final lo que le añadamos, no pertenecen al lenguaje. Es decir:
        \begin{equation*}
            [u\in A^* \mid \text{en $u$ hay una $a$ después de una $b$}]
        \end{equation*}
        Por comodidad, ya que $ba\notin L$, considerando este representante de clase de equivalencia, notaremos a esta clase de equivalencia con $[ba]$. Además, tenemos dos clases de equivalencia más:
        \begin{align*}
            [a] &= \{a^i\mid i\in \bb{N}\cup \{0\}\}\\
            [ab] &= \{a^ib^j\mid i\in \bb{N}\cup \{0\},j\in \bb{N}\}
        \end{align*}
        
        Veamos en primer lugar que no hay más clases de equivalencia. Dado $u\in A^*$, si $u$ tiene una $a$ después de una $b$, entonces $u\in [ba]$. En caso contrario, tenemos $u=a^ib^j$ con $i,j\in \bb{N}\cup \{0\}$.
        \begin{itemize}
            \item Si $j=0$, entonces $u\in [a]$.
            \item Si $j>0$, entonces $u\in [ab]$.
        \end{itemize}
        Por tanto, tenemos que no hay más clases de equivalencia. Veamos ahora que estas clases de equivalencia son disjuntas.
        \begin{itemize}
            \item Sea $u\in [ba]$. Como $u$ tiene una $a$ después de una $b$, entonces de forma directa tenemos que $u\notin [a], [ab]$.
            En conclusión, vemos que $[ba]\cap [a]=[ba]\cap [ab]=\emptyset$.

            \item Sea $u\in [ab]$. Como $u=a^ib^j$ con $i,j\in \bb{N}\cup \{0\}$ con $j\neq 0$, entonces $u\notin [a]$. Por tanto, vemos que $[ab]\cap [a]=\emptyset$.
        \end{itemize}

        Por tanto, vemos que las clases de equivalencia de $L$ son $[ba]$, $[a]$ y $[ab]$; y estas son disjuntas.

        \item Demostrar que $L$ es aceptado por un autómata finito determinista si y solo si el número de clases de equivalencia es finito.
        
        Demostramos mediante doble inclusión.
        \begin{description}
            \item[$\Longrightarrow)$] Supongamos que $L$ es aceptado por un autómata finito determinista. Sea $M=(Q, A, \delta, q_0, F)$ su AFD minimal que acepta $L$.
            Supongamos $u,v\in A^*$ tales que $u\equiv v$. Sean:
            \begin{align*}
                q_u &:= \delta^{\ast}(q_0, u)\in F,\\
                q_v &:= \delta^{\ast}(q_0, v)\in F.
            \end{align*}

            Veamos ahora que $q_u,q_v$ son indistingibles, es decir, $q_u=q_v$.
            \begin{itemize}
                \item Para todo $z\in A^*$, como $u\equiv v$, se tiene que:
                \begin{equation*}
                    uz\in L\Leftrightarrow vz\in L.
                \end{equation*}
    
                Equivalentemente, tenemos que:
                \begin{equation*}
                    \delta^{\ast}(\delta^{\ast}(q_0, u), z)\in F\Longleftrightarrow \delta^{\ast}(\delta^{\ast}(q_0, v), z)\in F
                \end{equation*}

                Es decir:
                \begin{equation*}
                    \delta^{\ast}(q_u, z)\in F\Longleftrightarrow \delta^{\ast}(q_v, z)\in F
                \end{equation*}

                Por tanto, $q_u$ y $q_v$ son indistinguibles; y como el autómata es minimal, $q_u=q_v$.
            \end{itemize}
            
            Por tanto, hemos demostrado que si $u\equiv v$, entonces $\delta^{\ast}(q_0, u)=\delta^{\ast}(q_0, v)$, por lo que el número de clases de equivalencia es menor o igual que el número de estados de $Q$. Como $Q$ es finito, el número de clases de equivalencia también lo es finito.

            \item[$\Longleftarrow)$] Supongamos que el número de clases de equivalencia es finito. Sea el autómata $M=(Q, A, \delta, q_0, F)$, donde:
            \begin{itemize}
                \item $Q$ es el conjunto de clases de equivalencia de $L$,
                \item $q_0=[\varepsilon]$,
                \item $\delta([u], a)=[ua]$. Veamos que está bien definida.
                
                Sea $u,v\in A^*$ tales que $u\equiv v$, y veamos que, para todo $a\in A$, $ua\equiv va$.
                Como $u\equiv v$, para todo $z\in A^*$, tenemos que:
                \begin{equation*}
                    uz\in L\Leftrightarrow vz\in L.
                \end{equation*}

                Por tanto, tomando $z=az'$, con $z'\in L$, tenemos que:
                \begin{equation*}
                    uaz'\in L\Leftrightarrow vaz'\in L.
                \end{equation*}

                Es decir, $ua\equiv va$. Por tanto, $\delta$ está bien definida.
                \item $F=\{[u]\in Q\mid u\in L\}$.
                
                Para ver que $F$ está bien definida, veamos que, si $u\equiv v$, entonces $u\in L\Longleftrightarrow v\in L$.
                Esto es directo tomando $z=\veps$.
            \end{itemize}

            Veamos ahora que $L=\cc{L}(M)$.
            \begin{align*}
                u\in \cc{L}(M)&\Longleftrightarrow \delta^{\ast}(q_0, u)\in F\Longleftrightarrow \delta^{\ast}([\varepsilon], u)\in F\Longleftrightarrow [\veps u]\in F\Longleftrightarrow [u]\in F \Longleftrightarrow \\ &\Longleftrightarrow u\in L
            \end{align*}

        \end{description}
        \item ¿Qué relación existe entre el número de clases de equivalencia y el autómata finito minimal que acepta $L$?
        
        Veamos que el autómata descrito en el apartado anterior es minimal. En primer lugar, hemos de demostrar que no tiene estados inaccesibles.
        \begin{itemize}
            \item Sea $q\in Q$ una clase de equivalencia de $L$. Tomando un representante $u\in q$, tenemos que $\delta^{\ast}(q_0, u)=q$. Por tanto, $q$ es accesible.
        \end{itemize}

        Veamos ahora que no tiene estados indistingibles.
        \begin{itemize}
            \item Sean $q_1,q_2\in Q$ clases de equivalencia distintas de $L$. Tomando representantes $u_1\in q_1,u_2\in q_2$, tenemos que:
            \begin{align*}
                \delta^{\ast}(q_0, u_1)=[u_1]=q_1,\\
                \delta^{\ast}(q_0, u_2)=[u_2]=q_2.
            \end{align*}
            Entonces, como son clases de equivalencia distintas, $u_1\not\equiv u_2$, lo cual implica que $\exists z\in A^*$ tal que $u_1z\in L$ y $u_2z\notin L$ o viceversa. Supongamos sin pérdida de generalidad el primer caso, luego:
            \begin{equation*}
                \exists z\in A^*\mid u_1z\in L\quad \land \quad u_2z\notin L
                \Longleftrightarrow
                \delta^{\ast}(q_0, u_1z)\in F\quad \land \quad \delta^{\ast}(q_0, u_2z)\notin F
            \end{equation*}
            Por tanto, $q_1$ y $q_2$ no son indistinguibles.
        \end{itemize}

        Por tanto, hemos visto que todos los estados de $Q$ son accesibles y no son indistinguibles, por lo que el autómata $M$ del ejercicio anterior es minimal. Por tanto, la relación es que el número de clases de equivalencia es igual al número de estados del autómata finito minimal que acepta $L$.
    \end{enumerate}
\end{ejercicio}

\begin{ejercicio}
    Dada una palabra $u = a_1 \cdots a_n \in A^*$, se llama $\operatorname{Per}(u)$ al conjunto
    \begin{align*}
        \{a_{\sigma(1)}, \ldots, a_{\sigma(n)} \mid \sigma \text{\ es una permutación de\ } \{1, \ldots, n\}\}.
    \end{align*}
    Dado un lenguaje $L$, se llama $\operatorname{Per}(L) = \bigcup\limits_{u \in L} \operatorname{Per}(u)$.
    Dar expresiones regulares y autómatas minimales para $\operatorname{Per}(L)$ en los siguientes casos:
    \begin{enumerate}% // TODO: Comprobar entero
        \item $L = (00 + 1)^*$,
        
        Tenemos que:
        \begin{align*}
            \operatorname{Per}(L) &= \{u\in A^*\mid \text{\ $n_0(u)$ es par}\}
        \end{align*}

        Su autómata finito minimal es:
        \begin{figure}[H]
            \centering
            \begin{tikzpicture}
                \node[state, initial, accepting] (q0) {$q_0$};
                \node[state, right of=q0] (q1) {$q_1$};

                \draw   (q0) edge[loop above] node{1} (q0)
                        (q0) edge[above, bend right] node{0} (q1)
                        (q1) edge[loop above] node{1} (q1)
                        (q1) edge[above, bend right] node{0} (q0);
            \end{tikzpicture}
        \end{figure}
        
        Este es de forma directa minimal, puesto que sus dos estados son distinguibles al ser uno final y el otro no. Para obtener la expresión regular, resolvemos el sistema de ecuaciones:
        \begin{align*}
            q_0 &= 1q_0 + 0q_1 + \veps\\
            q_1 &= 1q_1 + 0q_0
        \end{align*}
        De la segunda ecuación, obtenemos $q_1=1^*0q_0$. Sustituyendo en la primera, obtenemos:
        \begin{align*}
            q_0 &= 1q_0 + 0(1^*0q_0) + \veps\\
            q_0 &= 1q_0 + 01^*0q_0 + \veps\\
            q_0 &= (1+01^*0)q_0 + \veps\\
            q_0 &= (1+01^*0)^*
        \end{align*}

        De forma directa, podríamos haber obtenido la siguiente expresión regular:
        \begin{align*}
            1^*(01^*01^*)^*
        \end{align*}
        \item $L = (0 + 1)^*0$,
        
        Tenemos que $\operatorname{Per}(L) = \{u\in A^*\mid n_0(u)>0\}$.
        Su autómata finito minimal es:
        \begin{figure}[H]
            \centering
            \begin{tikzpicture}
                \node[state, initial] (q0) {$q_0$};
                \node[state, right of=q0, accepting] (q1) {$q_1$};

                \draw   (q0) edge[loop above] node{1} (q0)
                        (q0) edge[above] node{0} (q1)
                        (q1) edge[loop above] node{0,1} (q1);
            \end{tikzpicture}
        \end{figure}
        Este es de forma directa minimal, puesto que sus dos estados son distinguibles al ser uno final y el otro no. Para obtener la expresión regular, resolvemos el sistema de ecuaciones:
        \begin{align*}
            q_0 &= 1q_0+0q_1\\
            q_1 &= (0+1)q_1 + \veps
        \end{align*}
        Tenemos por tanto que $q_1=(0+1)^*$, y sustituyendo en la primera ecuación, obtenemos:
        \begin{equation*}
            q_0 = 1q_0 + 0(0+1)^* = 1^*\red{0}(0+1)^*
        \end{equation*}

        \item $L = (01)^*$.
        
        Tenemos que $\operatorname{Per}(L) = \{u\in A^*\mid n_0(u)=n_1(u)\}$.
        En este caso, veamos que el lenguaje no es regular usando el Lema de Bombeo. Para todo $n\in \bb{N}$, consideramos la palabra $z=0^n1^n$, que cumple $|z|\geq n$. Si consideramos una descomposición $z=uvw$ con $|uv|\leq n$ y $|v|\geq 1$, tenemos que:
        \begin{equation*}
            u=0^k,\quad v=0^l,\quad w=0^{n-k-l}1^n, \qquad \text{\ con\ } k+l\leq n, l\geq 1
        \end{equation*}
        Entonces, tomando $i=2$, tenemos que $uv^2w=0^{n+l}1^n\notin L$. Por tanto, $L$ no es regular.
    \end{enumerate}
    ¿Es posible que, siendo $L$ regular, $\operatorname{Per}(L)$ no lo sea?\\

    Como hemos visto en el apartado anterior, el lenguaje $L=(01)^*$ es regular, pero su permutación no lo es. Por tanto, es posible que, siendo $L$ regular, $\operatorname{Per}(L)$ no lo sea.
\end{ejercicio}

\subsection{Preguntas Tipo Test}
Se pide discutir la veracidad o falsedad de las siguientes cuestiones.
\begin{enumerate}
    \item Si $r$ y $s$ son expresiones regulares, siempre se verifica que ${(rs)}^{\ast}=r^\ast s^\ast$.\\

        Falso, si $r = 0$ y $s = 1$, tenemos que:
        \begin{itemize}
            \item La expresión regular ${(01)}^{\ast}$ está asociada al lenguaje $L_1 = \{{(01)}^{i}\mid i \in \mathbb{N}\}$.
            \item La expresión regular $0^\ast 1^\ast$ está asociada al lenguaje $L_2 = \{0^i 1^j \mid i,j\in \mathbb{N}\}$.
        \end{itemize}
        Y los lenguajes no son iguales, ya que $011\in L_2 \setminus L_1$.
    \item Si $r$ y $s$ son expresiones regulares, siempre se verifica que ${(r+s)}^{\ast} = r^\ast + s^\ast$.\\

        Falso, si $r = 0$ y $s= 1$, tenemos que:
        \begin{itemize}
            \item La expresión regular ${(0+1)}^{\ast}$ está asociada al lenguaje $L_1={\{0,1\}}^{\ast}$.
            \item La expresión regular $0^\ast + 1^\ast$ está asociada al lenguaje $L_2 = {\{0\}}^{\ast}\cup {\{1\}}^{\ast}$.
        \end{itemize}
        Y como $010\in L_1\setminus L_2$, no son iguales.
    \item Si $r_1$ y $r_2$ son expresiones regulares tales que su lenguaje asociado contiene la palabra vacía, entonces ${(r_1r_2)}^{\ast}={(r_2r_1)}^{\ast}$\\

        Verdadero, sean $R_1$ y $R_2$ los lenguajes asociados a $r_1$ y a $r_2$ respectivamente, tratamos de comprobar si
        \begin{equation*}
            {(R_1R_2)}^{\ast} = {(R_2R_1)}^{\ast}
        \end{equation*}
        con la condición de que $\veps \in R_1 \cap R_2$.
        \begin{itemize}
            \item Sea $u\in {(R_1R_2)}^{\ast}$, entonces $u$ será de la forma
                \begin{equation*}
                    u = v_1w_1v_2w_2\ldots v_nw_n \qquad v_i \in R_1, \quad w_i \in R_2 \qquad \forall i \in \{1,\ldots,n\}
                \end{equation*}
                y podemos reescribir $u$ como:
                \begin{equation*}
                    u = \veps v_1w_1v_2w_2 \ldots v_nw_n\veps \in {(R_2R_1)}^{\ast}
                \end{equation*}
            \item De forma análoga, si $u\in {(R_2R_1)}^{\ast}$, podemos llegar a que $u\in {(R_1R_2)}^{\ast}$.
        \end{itemize}
    \item Si $r$ y $s$ son expresiones regulares, siempre se verifica que ${(r+\veps )}^{\ast}=r^\ast$.\\

        Verdadero, sea $R$ el lenguaje asociado a $r$, tenemos que
        \begin{equation*}
            {(R\cup \{\veps\})}^{\ast} = R^\ast
        \end{equation*}
        con lo que ${(r+\veps)}^{\ast}=r^\ast$.
    \item Si $r$ y $s$ son expresiones regulares, siempre se verifica que $r{(r+s)}^{\ast}={(r+s)}^{\ast}r$.\\

        Falso, sean $r = 0$ y $s = 1$, tenemos que:
        \begin{itemize}
            \item $0{(0+1)}^{\ast}$ es la expresión regular asociada al lenguaje \newline$L_1=\{0u \mid u\in {\{0,1\}}^{\ast}\}$.
            \item ${(0+1)}^{\ast}0$ es la expresión regular asociada al lenguaje \newline$L_2 = \{u0 \mid u\in {\{0,1\}}^{\ast}\}$.
        \end{itemize}
        Como $0111 \in L_1\setminus L_2$, los lenguajes no son iguales.
    \item Si $r_1$ y $r_2$ son expresiones regulares, entonces $r_1^\ast r_2^\ast \subseteq {(r_1r_2)}^{\ast}$, en el sentido de que los lenguajes asociados están incluidos.\\

        Falso, si $r_1 = 1$ y $r_2 = 0$, estaríamos afirmando que $1^\ast 0^\ast \subseteq {(10)}^{\ast}$, lo cual es falso, ya que $110$ es una palabra del lenguaje asociado a la primera expresión regular pero no al lenguaje asociado a la segunda.
    \item \label{item:preg2_7}
    Si $r_1$, $r_2$ y $r_3$ son expresiones regulares, entonces ${(r_1+r_2)}^{\ast}r_3=r_1^\ast r_3 + r_2^\ast r_3$.\\

        Falso, si $r_1 = a$, $r_2 = b$ y $r_3 = c$, estaríamos afirmando que 
        \begin{equation*}
            {(a+b)}^{\ast}c = a^\ast c + b^\ast c
        \end{equation*}
        Pero $abc$ es una palabra del lenguaje asociado a la primera expresión regular que no es del lenguaje asociado a la segunda expresión.
    \item \label{item:preg2_8}
     Si $r_1$ y $r_2$ son expresiones regulares entonces: ${(r_1^\ast r_2^\ast)}^{\ast}={(r_1+r_2)}^{\ast}$.\\

        Verdadero, sean $R_1$ y $R_2$ los lenguajes asociados a $r_1$ y a $r_2$ respectivamente, tratamos de ver que
        \begin{equation*}
            {(R_1^\ast R_2^\ast)}^{\ast} = {(R_1\cup R_2)}^{\ast}
        \end{equation*}
        Lo cual puede demostrarse que es cierto.
    \item Si $r_1$ y $r_2$ son expresiones regulares, entonces ${(r_1r_2)}^{\ast}={(r_1+r_2)}^{\ast}$.\\

        Falso, ya que si $r_1 = 0$ y $r_2 = 1$, estaríamos afirmando que 
        \begin{equation*}
            {(01)}^{\ast} = {(0+1)}^{\ast}
        \end{equation*}
        Pero $0$ es una palabra del lenguaje asociado a la segunda expresión regular y no del lenguaje asociado a la primera.
    \item Si $r$ es una expresión regular, entonces $r^\ast r^\ast = r^\ast$.\\

        Verdadero, si $R$ es el lenguaje asociado, tenemos que ver que $R^\ast R^\ast = R^\ast$:
        \begin{itemize}
            \item Sea $u\in R^\ast R^\ast$, entonces $u$ es de la forma:
                \begin{equation*}
                    u = v_1v_2\ldots v_nv_{n+1}v_{n+2}\ldots v_m \qquad v_1v_2\ldots v_n, v_{n+1}v_{n+2},\ldots v_m \in R^\ast
                \end{equation*}
                por lo que $v_i \in R$ $\forall i \in \{1,\ldots, m\}$ con lo que $u\in R^\ast$.

                Si por contrario $u = \veps$, entonces $u\in R^\ast$.
            \item Sea $u\in R^\ast$, entonces $u$ será de la forma:
                \begin{equation*}
                    u = v_1v_2\ldots v_n \qquad v_i \in R \quad \forall i \in \{1,\ldots, n\} \qquad n > 1
                \end{equation*}
                por lo que podemos tomar $w = v_2\ldots v_n \in R^\ast$ y $v_1\in R \subseteq R^\ast$, llegando a que
                \begin{equation*}
                    u = v_1w \qquad v_1,w \in R^\ast \Longrightarrow u\in R^\ast R^\ast
                \end{equation*}
                \begin{itemize}
                    \item En el caso en el que $u\in R\subseteq R^\ast$, entonces $u = \veps u$ con $\veps \in R^\ast$.
                    \item Además, si $u = \veps$, entonces $u = \veps \veps$ con $\veps \in R^\ast$.
                \end{itemize}
        \end{itemize}
    \item Si $r$ es una expresión regular, entonces $r\emptyset =r+\emptyset $.\\

        Falso, sea $R$ el lenguaje asociado a $r$, estaríamos afirmando que
        \begin{equation*}
            \emptyset  = R\emptyset  = R \cup \emptyset  = R
        \end{equation*}
        lo cual no es cierto, a no ser que $r = \emptyset $.
    \item Si $r$ es una expresión regular, entonces se verifica que $r^\ast\veps =r^+\veps$.\\

        Verdadero, sea $R$ el lenguaje asociado a $r$, entonces:
        \begin{equation*}
            R^\ast \cup \{\veps\} = R^\ast = R^+ \cup \{\veps\}
        \end{equation*}
    \item \label{item:preg13}
    Si $r_1$ y $r_2$ son expresiones regulares, entonces siempre $r_1{(r_2r_1)}^{\ast}={(r_1r_2)}^{\ast}r_1$.\\

        Verdadero, sean $R_1$ y $R_2$ los lenguajes asociados a $r_1$ y a $r_2$ respectivamente, entonces, tratamos de probar que
        \begin{equation*}
            R_1{(R_2R_1)}^{\ast} = {(R_1R_2)}^{\ast}R_1
        \end{equation*}
        \begin{itemize}
            \item Sea $u\in R_1{(R_2R_1)}^{\ast}$, entonces $u$ es de la forma
                \begin{equation*}
                    u = v_0w_1v_1w_2v_2 \ldots w_n v_n \qquad v_i \in R_1 \quad w_i \in R_2
                \end{equation*}
                con lo que podemos tomar:
                \begin{equation*}
                    v = v_0w_1v_1w_2v_2 \ldots w_n \in {(R_1R_2)}^{\ast}
                \end{equation*}
                y tenemos que $u = vv_n \in {(R_1R_2)}^{\ast}R_1$.
            \item Puede probarse de forma análoga que si $u\in {(R_1R_2)}^{\ast}R_1$, entonces $u\in R_1{(R_2R_1)}^{\ast}$.
        \end{itemize}
    \item Si $r_1$ y $r_2$ son expresiones regulares, siempre se verifica que $r_1{(r_2r_1)}^{\ast}={(r_1r_2)}^{\ast}r_1$.\\

        Verdadero, la pregunta es idéntica a la Pregunta~\ref{item:preg13}.
    \item Si $r$ y $s$ son expresiones regulares, entonces ${(r^\ast s^\ast)}^{\ast}={(r+s)}^{\ast}$.\\

        Verdadero, la pregunta es idéntica a la Pregunta~\ref{item:preg2_8}.
    \item Si $r$ es una expresión regular, entonces ${(rr)}^{\ast}\subseteq r^\ast$.\\

        Verdadero, sea $R$ el lenguaje asociado a $r$, tratamos de ver que 
        \begin{equation*}
            {(RR)}^{\ast} \subseteq  R^\ast
        \end{equation*}
        Sea $u\in {(RR)}^{\ast}$, entonces $u$ es de la forma
        \begin{equation*}
            u = v_1v_1'v_2v_2'\ldots v_nv_n' \qquad v_i, v_i' \in R \quad \forall i \in \{1,\ldots, n\}
        \end{equation*}
        por lo que $u\in R^\ast$.
    \item Si $r_1$ y $r_2$ son expresiones regulares, tales que su lenguaje asociado contiene la palabra vacía, entonces ${(r_1r_2)}^{\ast}={(r_1+r_2)}^{\ast}$.\\

        Verdadero, puede probarse de forma similar a como lo hicimos en la pregunta~\ref{item:preg8}.
    \item Si $r_1$, $r_2$, $r_3$ son expresiones regulares, entonces $r_1(r_2^\ast + r_3^\ast) = r_1r_2^\ast + r_1r_3^\ast$.\\

        Falso, podemos dar un contraejemplo de forma similar a como lo hicimos en la pregunta~\ref{item:preg2_7}.
    \item La demostración de que la clase de lenguajes aceptados por los autómatas no deterministas es la misma que la aceptada por los autómatas determistas, se basa en dado un autómata no determinista construir uno determinista que, ante una palabra de entrada, explore todas las posibles opciones que puede seguir el no determinista.\\

        Verdadero, así se comprueba que todo lenguaje aceptado por un autómata no determinista puede ser aceptado por uno determinista. No hace falta demostrarlo en la otra dirección, ya que un autómata determinista es a su vez no determinista.
    \item Un autómata finito puede ser determinista y no-determinista a la vez.\\

        Verdadero, ya que todos los autómatas deterministas son a su vez no deterministas.
    \item Para transformar un autómata que acepta el lenguaje $L$ en uno que acepte $L^\ast$, basta unir los estados finales con el inicial mediante transiciones nulas.\\

        Falso, tenemos que hacer que el estado inicial sean también final, para que la palabra $\veps$ sea aceptada por el autómata, en caso de que $\veps \notin L$.
    \item Para pasar de un autómata que acepte el lenguaje asociado a $r$ a uno que acepte $r^\ast$ basta con unir con transiciones nulas sus estados finales con el estado inicial.\\

        Falso, es la misma pregunta que la 21.
    \item Existe un lenguaje reconocido por un AFD y no generado por una gramática independiente del contexto.\\

        Falso, si un lenguaje es reconocido por un AFD, hemos visto en teoría que hay una gramática de tipo 3 que genera dicho lenguaje y como las gramáticas tipo 3 son a su vez independientes del contexto, el enunciado es falso.
    \item Existen lenguajes aceptados por AFD que no pueden ser aceptados por AF no determinísticos.\\

        Falso, el conjunto de lenguajes aceptados por un AFD coincide con el conjunto de lenguajes aceptados por AF no determinísticos, tal y como se ha visto en teoría.
    \item La clausura de un lenguaje aceptado por un AFD puede ser representado con una expresión regular.\\

        Verdadero, sea $L$ un lenguaje aceptado por un AFD, conocemos por teoría que podemos encontrar una expresión regular $r$ asociada a dicho lenguaje. Si ahora consideramos $r^\ast$, esta expresión regular estará asociada al lenguaje $L^\ast$.
    \item Un lenguaje representado por una expresión regular siempre puede ser reconocido por un AF no determinista.\\

        Verdadero, hemos visto en teoría que el conjunto de lenguajes representados por expresiones regulares coincide con el conjunto de lenguajes reconocidos por cualquier tipo de AF\@.
    \item Todo lenguaje regular puede ser generado por una gramática libre de contexto.\\

        Verdadero, ya que todo lenguaje regular puede ser generado por una gramática lineal por la derecha, que a su vez es independiente del contexto.
    \item Un lenguaje con un número finito de palabras siempre puede ser reconocido por un AF no determinista.\\

        Verdadero, sea $L = \{u_1, u_2, \ldots, u_n\}$ un lenguaje finito de $n$ palabras, entonces podemos considerar el autómata formado por un estado inicial $q_0$ y para cada palabra $u_i$ con $i \in \{1,\ldots,n\}$ de longitud $k=|u_i|$, añadimos $k$ nuevos estados a los que llamaremos por ejemplo $q_{ij}$ con $j\in \{1,\ldots,k\}$. De esta forma, si 
        \begin{equation*}
            u_i = a_{i1}a_{i2}\ldots a_{ik} \qquad a_{ij} \in A \quad \forall j \in \{1,\ldots,k\}
        \end{equation*}
        Entonces, añadimos las transiciones $\delta(q_{ij-1}, a_j) = q_{ij}$, entendiendo que $q_{i0} = q_0$ para cualquier $i \in \{1,\ldots,n\}$ y consideraremos que $q_{ik}$ es un estado final.

        De manera más formal, dado un lenguaje finito $L=\{u_1,u_2, \ldots, u_n\}$ sobre un alfabeto $A$, construimos el autómata $M=(Q,A,\delta, q_0, F)$ formado por los conjuntos:
        \begin{align*}
            Q &= \{q_0\} \cup \{q_{ij_i} \mid i \in \{1,\ldots,n\}, j_i \in \{1,\ldots, |u_i|\}\} \\
            F &= \{q_{ij_i} \mid j_i = |u_i|\}
        \end{align*}
        donde la función de transición $\delta$ viene dada por:
        \begin{align*}
            \delta(q_0, a) &= \left\{\begin{array}{ccl}
                    \{q_{i1}\} & \text{\ si\ } & a = a_{i1} \\
                    \emptyset & \text{\ si\ } & a\neq a_{i1} \quad \forall i \in \{1,\ldots,n\}
            \end{array}\right. \\
            \delta(q_{ij_i}, a) &= \left\{\begin{array}{ccl}
                    \{q_{ij_i+1}\} & \text{\ si\ } & a = a_{ij_i} \\
                    \emptyset & \text{\ si\ } & a \neq a_{ij_i} \quad \forall i \in \{1,\ldots,n\}, j_i\in \{1,\ldots,|u_i|\}
            \end{array}\right.
        \end{align*}
        Es decir, para cada palabra creamos un camino desde el estado inicial a un estado final que se recorra leyendo dicha palabra.

        Alternativamente, podríamos haber dicho que si $L=\{u_1,u_2,\ldots, u_n\}$ es un lenguaje finito, entonces podemos considerar la gramática $G=(\{S\}, A, P, S)$ donde $P$ es el conjunto:
        \begin{equation*}
            P = \{S\rightarrow u_1,  S\rightarrow u_2,  \ldots,  S\rightarrow u_n\}
        \end{equation*}
        Y argumentar que dicha gramática puede pasarse a un AF no determinista por la teoría vista en el Tema 2.
    \item Todo autómata finito determinista de $n$ estados, cuyo alfabeto $A$ contiene $m$ símbolos debe tener $m\cdot n$ transiciones.\\

        Verdadero, todo estado de un autómata finito determinista debe tener tantas transiciones como símbolos tenga su alfabeto $A$ de entrada, en este caso, $m$. Como en este caso el autómata finito determinista tiene $n$ estados, entonces el número de transiciones del autómata es:
        \begin{equation*}
            \sum_{i=1}^{n} m = n\cdot m
        \end{equation*}
    \item Para que un autómata con pila sea determinista es necesario que no tenga transiciones nulas.\\

        No hemos visto autómatas con pila todavía, esta pregunta no es parte del Tema 2.
    \item Si $r_1$ y $r_2$ son expresiones regulares, entonces siempre se tiene que ${(r_1+r_2)}^{\ast}={(r_1^\ast r_2)}^{\ast}r_1^\ast$.\\

        Verdadero, puede razonarse buscando la igualdad entre los lenguajes que representan ambas expresiones.
    \item Si un lenguaje es infinito no se puede encontrar una expresión regular que lo represente.\\

        Falso, el lenguaje $L=\{1^n \mid n \in \mathbb{N}\}$ es infinito por ser $\mathbb{N}$ infinito y puede representarse por la expresión regular $1^\ast$.
    \item Si $r_1$ y $r_2$ son expresiones regulares, entonces se verifica que ${(r_1+\veps)}^{+}r_2^+ = r_1^+ {(r_2+\veps)}^{+}$.\\

        Falso, si fuera cierto, entonces:
        \begin{equation*}
            r_1^\ast r_2^+ = {(r_1+\veps)}^{+}r_2^+ = r_1^+ {(r_2+\veps)}^{+} = r_1^+ r_2^\ast
        \end{equation*}
        para cualesquiera $r_1$ y $r_2$ expresiones regulares.

        Sin embargo, si tomamos $r_1 = 0$ y $r_2 = 1$, $1$ es una palabra del lenguaje asociado a la expresión regular $0^\ast 1^+$ pero no lo es del lenguaje asociado a la expresión regular $0^+ 1^\ast$, al no contener ningún ``0''.
    \item El conjunto de palabras sobre el alfabeto $\{0,1\}$ tales que eliminando los tres últimos símbolos, en la palabra resultante no aparece el patrón $0011$ es un lenguaje regular.\\

        Verdadero, ya que podemos dar un AFD que reconozca dicho lenguaje:

    \begin{figure}[H]
        \centering
    \begin{tikzpicture}
        \node[state, accepting, initial] (q0) {$q_0$};
        \node[state, accepting, right of=q0] (q1) {$q_1$};
        \node[state, accepting, right of=q1] (q2) {$q_2$};
        \node[state, accepting, right of=q2] (q3) {$q_3$};
        \node[state, accepting, right of=q3] (q4) {$q_4$};
        \node[state, accepting, below of=q4] (q5) {$q_5$};
        \node[state, accepting, left of=q5] (q6) {$q_6$};
        \node[state, error, left of=q6] (q7) {$E$};

        \draw (q0) edge[loop above] node{1} (q0)
              (q0) edge[bend left, above] node{0} (q1)
              (q1) edge[bend left, below] node{1} (q0)
              (q1) edge[left, above] node{0} (q2)
              (q2) edge[loop below] node{0} (q2)
              (q2) edge[right, above] node{1} (q3)
              (q3) edge[right, above] node{1} (q4)
              (q3) edge[bend right, above] node{0} (q1)
              (q4) edge[below, right] node{0, 1} (q5)
              (q5) edge[left, above] node{0, 1} (q6)
              (q6) edge[left, above] node{0, 1} (q7)
              (q7) edge[loop left] node{0, 1} (q7);
    \end{tikzpicture}
    \caption{Autómata Finito Determinista para la pregunta 34.}
    \end{figure}

    \item El lenguaje formado por las cadenas sobre $\{0,1\}$ que tienen un número impar de $0$ y un número par de $1$ no es regular.\\

        Falso, ya que podemos dar una gramática $G=(\{A,B,C,D\}, \{0,1\}, P, A)$ lineal por la derecha que genere dicho lenguaje. Si entendemos los estados $\{A,B,C,D\}$ como:
        \begin{itemize}
            \item $A$ quiere decir que la palabra generada hasta el momento contiene un número par de 0s y de 1s.
            \item $B$ quiere decir que la palabra generada hasta el momento contiene un número impar de 0s y par de 1s.
            \item $C$ quiere decir que la palabra generada hasta el momento contiene un número par de 0s e impar de 1s.
            \item $D$ quiere decir que la palabra generada hasta el momento contiene un número impar de 0s y de 1s.
        \end{itemize}
        Entonces, podemos dar las siguientes reglas de producción, que son las únicas reglas de producción que contiene $P$:
        \begin{align*}
            A &\rightarrow 0B\ |\ 1C \\
            B &\rightarrow \veps\ |\ 0A\ |\ 1D \\
            C &\rightarrow 1A\ |\ 0D \\
            D &\rightarrow 1B\ |\ 0C 
        \end{align*}
        Notemos que la variable inicial es $A$ ya que $\veps$ contiene un número par de 0s y de 1s.
\end{enumerate}

\end{document}
